% ============================================================================
% QuantileFlow — Propuesta técnica (versión de investigación)
% Compilación: latexmk -pdf main.tex   (o `make pdf` desde la raíz)
% ============================================================================
\documentclass[11pt,letterpaper,oneside]{article}
% ============================================================================
% Preámbulo: tipografía, color, cajas y estilos de diagramas
% ============================================================================
\usepackage[T1]{fontenc}
\usepackage[utf8]{inputenc}
\usepackage[spanish,es-noshorthands,es-tabla,es-nolists,es-nodecimaldot]{babel}
\usepackage[sc]{mathpazo}
\linespread{1.06}
\usepackage[scaled=0.90]{helvet}
\renewcommand\ttdefault{lmtt}
\usepackage{microtype}
\usepackage[letterpaper,textwidth=6.5in,textheight=8.9in,top=1.0in,headsep=0.28in,footskip=0.42in]{geometry}

\usepackage{amsmath,amssymb,mathtools,bm}
\usepackage{graphicx}
\usepackage[table]{xcolor}
\usepackage{tikz}
\usetikzlibrary{babel,positioning,arrows.meta,calc,fit,backgrounds,shapes.geometric,
  shapes.misc,decorations.pathreplacing,patterns,matrix,shadows}
\usepackage[most]{tcolorbox}
\usepackage{booktabs,tabularx,array,multirow,makecell}
\usepackage{enumitem}
\usepackage{caption}
\usepackage{subcaption}
\usepackage{float}
\usepackage{fancyhdr}
\usepackage{titlesec}
\usepackage{titletoc}
\usepackage{pifont}
\usepackage{listings}
\usepackage[round,sort]{natbib}
\usepackage{xurl}
\usepackage[hidelinks]{hyperref}
\usepackage{ragged2e}

% ---------------------------------------------------------------------------
% Paleta (misma que las figuras; orden categórico validado para daltonismo)
% ---------------------------------------------------------------------------
\definecolor{qazul}{HTML}{2A78D6}
\definecolor{qazuloscuro}{HTML}{104281}
\definecolor{qnaranja}{HTML}{EB6834}
\definecolor{qaqua}{HTML}{1BAF7A}
\definecolor{qamarillo}{HTML}{EDA100}
\definecolor{qvioleta}{HTML}{4A3AA7}
\definecolor{tinta}{HTML}{0B0B0B}
\definecolor{tinta2}{HTML}{52514E}
\definecolor{tenue}{HTML}{898781}
\definecolor{rejilla}{HTML}{E1E0D9}
\definecolor{eje}{HTML}{C3C2B7}
\definecolor{grisclaro}{HTML}{F0EFEC}
\definecolor{grisfondo}{HTML}{F7F7F5}
\definecolor{bueno}{HTML}{0CA30C}
\definecolor{alerta}{HTML}{FAB219}
\definecolor{serio}{HTML}{EC835A}
\definecolor{critico}{HTML}{D03B3B}
\definecolor{azulfondo}{HTML}{EEF4FC}
\definecolor{naranjafondo}{HTML}{FDF1EB}
\definecolor{aquafondo}{HTML}{E8F7F1}
\definecolor{alertafondo}{HTML}{FFF7E6}
\definecolor{buenofondo}{HTML}{EAF6EA}

\hypersetup{colorlinks=true,linkcolor=qazuloscuro,citecolor=qazuloscuro,urlcolor=qazuloscuro,
  pdftitle={QuantileFlow: seguimiento de distribuciones implícitas y gestión diaria de posiciones},
  pdfsubject={Propuesta técnica (versión de investigación)},pdfauthor={QuantileFlow},
  pdfkeywords={SSVI, Breeden-Litzenberger, Wasserstein, Kalman, BOCPD, CVaR, conformal}}

% ---------------------------------------------------------------------------
% Títulos, encabezados e índice
% ---------------------------------------------------------------------------
\setcounter{secnumdepth}{3}
\setcounter{tocdepth}{2}
\titleformat{\section}{\sffamily\LARGE\bfseries\color{tinta}}
  {\color{qazul}\thesection}{0.55em}{}[\vspace{3pt}{\color{eje}\titlerule[0.6pt]}]
\titlespacing*{\section}{0pt}{0pt}{12pt}
\titleformat{\subsection}{\sffamily\large\bfseries\color{tinta}}{\color{qazul}\thesubsection}{0.6em}{}
\titlespacing*{\subsection}{0pt}{16pt}{6pt}
\titleformat{\subsubsection}{\sffamily\normalsize\bfseries\color{tinta}}{\color{qazul}\thesubsubsection}{0.6em}{}
\titlespacing*{\subsubsection}{0pt}{11pt}{4pt}
\titleformat{\paragraph}[runin]{\sffamily\normalsize\bfseries\color{tinta}}{}{0pt}{}[.]

\titlecontents{section}[1.8em]{\addvspace{2.5pt}\sffamily\bfseries}
  {\contentslabel[\color{qazul}\thecontentslabel]{1.8em}}{}{\hfill\contentspage}
\titlecontents{subsection}[4.2em]{\small}
  {\contentslabel[\color{tinta2}\thecontentslabel]{2.4em}}{}{\titlerule*[0.6pc]{.}\contentspage}

\pagestyle{fancy}
\fancyhf{}
\renewcommand{\headrulewidth}{0.4pt}
\renewcommand{\headrule}{{\color{eje}\hrule width\headwidth height\headrulewidth}}
\fancyhead[L]{\sffamily\scriptsize\color{tinta2}QuantileFlow \textperiodcentered{} Propuesta técnica \textperiodcentered{} Versión de investigación}
\fancyhead[R]{\sffamily\scriptsize\color{tinta2}\nouppercase{\leftmark}}
\fancyfoot[C]{\sffamily\footnotesize\color{tinta2}\thepage}
\renewcommand{\sectionmark}[1]{\markboth{\thesection\quad #1}{}}
\fancypagestyle{plain}{\fancyhf{}\renewcommand{\headrulewidth}{0pt}\fancyfoot[C]{\sffamily\footnotesize\color{tinta2}\thepage}}

% ---------------------------------------------------------------------------
% Pies de figura, listas y tablas
% ---------------------------------------------------------------------------
\captionsetup{font={small,sf},labelfont={bf,color=qazuloscuro},labelsep=period,
  justification=justified,singlelinecheck=false,skip=6pt}
\captionsetup[table]{position=top}
\setlist{itemsep=2pt,topsep=4pt,parsep=0pt,leftmargin=1.4em}
\setlist[itemize,1]{label={\color{qazul}\small\textbullet}}
\setlist[itemize,2]{label={\color{tenue}\textendash}}
\setlist[enumerate,1]{label={\sffamily\bfseries\color{qazul}\arabic*.}}
\renewcommand{\arraystretch}{1.2}
\setlength{\aboverulesep}{0.3ex}
\setlength{\belowrulesep}{0.5ex}
\newcolumntype{Y}{>{\RaggedRight\arraybackslash}X}
\newcommand{\encab}[1]{\textsf{\bfseries\footnotesize #1}}

\setlength{\parindent}{0pt}
\setlength{\parskip}{5pt plus 1pt}
\setlength{\emergencystretch}{2em}
\raggedbottom
\renewcommand{\floatpagefraction}{0.75}
\renewcommand{\topfraction}{0.9}
\renewcommand{\textfraction}{0.08}

% ---------------------------------------------------------------------------
% Iconos
% ---------------------------------------------------------------------------
\newcommand{\iconoidea}{\tikz[baseline=-0.7ex]\node[circle,fill=qazul,text=white,
  font=\sffamily\bfseries\fontsize{6}{6}\selectfont,inner sep=0pt,minimum size=3mm]{i};}
\newcommand{\iconoalerta}{\tikz[baseline=-0.35ex]{\fill[alerta,rounded corners=0.3pt]
  (0,0)--(0.34,0)--(0.17,0.30)--cycle;\node[font=\sffamily\bfseries\fontsize{5.5}{5.5}\selectfont,
  text=tinta] at (0.17,0.105){!};}}
\newcommand{\iconocheck}{\tikz[baseline=-0.7ex]\node[circle,fill=bueno,text=white,
  font=\fontsize{6}{6}\selectfont,inner sep=0pt,minimum size=3mm]{\ding{51}};}
\newcommand{\iconopendiente}{\tikz[baseline=-0.7ex]\node[circle,draw=tenue,line width=0.6pt,
  text=tinta2,font=\sffamily\bfseries\fontsize{6}{6}\selectfont,inner sep=0pt,minimum size=3mm]{?};}
\newcommand{\icononota}{\tikz[baseline=-0.7ex]\node[rectangle,rounded corners=0.6pt,fill=qvioleta,
  text=white,font=\sffamily\bfseries\fontsize{6}{6}\selectfont,inner sep=0pt,minimum size=3mm]{N};}
\newcommand{\si}{\textcolor{bueno}{\ding{51}}}
\newcommand{\no}{\textcolor{critico}{\ding{55}}}

% ---------------------------------------------------------------------------
% Cajas temáticas
% ---------------------------------------------------------------------------
\tcbset{
  qfbase/.style={enhanced,breakable,boxrule=0pt,arc=0pt,outer arc=0pt,left=9pt,right=9pt,
    top=6pt,bottom=6pt,fontupper=\small,before skip=9pt,after skip=9pt,
    fonttitle=\sffamily\bfseries\small,coltitle=tinta,attach title to upper={\par\smallskip}},
}
\newtcolorbox{ideaclave}[1][Idea clave]{qfbase,colback=azulfondo,
  borderline west={2.2pt}{0pt}{qazul},title={\iconoidea\hspace{0.45em}#1}}
\newtcolorbox{ideaclavefinal}[1][Idea clave]{qfbase,unbreakable,colback=azulfondo,
  borderline west={2.2pt}{0pt}{qazul},title={\iconoidea\hspace{0.45em}#1}}
\newtcolorbox{noconcluir}[1][Qué no concluir]{qfbase,colback=alertafondo,
  borderline west={2.2pt}{0pt}{alerta},title={\iconoalerta\hspace{0.45em}#1}}
\newtcolorbox{criterio}[1][Criterio para avanzar]{qfbase,colback=buenofondo,
  borderline west={2.2pt}{0pt}{bueno},title={\iconocheck\hspace{0.45em}#1}}
\newtcolorbox{pendiente}[1][Pendiente de definir]{qfbase,colback=grisfondo,
  borderline west={2.2pt}{0pt}{tenue},title={\iconopendiente\hspace{0.45em}#1}}
\newtcolorbox{notatecnica}[1][Nota técnica]{qfbase,colback=grisfondo,
  borderline west={2.2pt}{0pt}{qvioleta},title={\icononota\hspace{0.45em}#1}}
\newtcolorbox{algoritmo}[1]{qfbase,colback=white,colframe=eje,boxrule=0.5pt,
  borderline west={0pt}{0pt}{white},title={#1},fontupper=\small}

% Ficha de etapa: entradas, proceso, salidas y controles
\newcommand{\ficha}[4]{%
  \begin{tcolorbox}[enhanced,boxrule=0pt,arc=1.5pt,colback=grisfondo,left=7pt,right=7pt,top=6pt,
    bottom=5pt,before skip=4pt,after skip=12pt]
  \small
  \begin{tabularx}{\linewidth}{@{}>{\RaggedRight}X@{\hspace{10pt}}>{\RaggedRight}X@{\hspace{10pt}}>{\RaggedRight}X@{\hspace{10pt}}>{\RaggedRight\arraybackslash}X@{}}
  {\sffamily\bfseries\footnotesize\color{qazul}ENTRADAS} &
  {\sffamily\bfseries\footnotesize\color{qazul}PROCESO} &
  {\sffamily\bfseries\footnotesize\color{qazul}SALIDAS} &
  {\sffamily\bfseries\footnotesize\color{qazul}CONTROLES}\\[2pt]
  \footnotesize #1 & \footnotesize #2 & \footnotesize #3 & \footnotesize #4
  \end{tabularx}
  \end{tcolorbox}}

% ---------------------------------------------------------------------------
% Estilos TikZ
% ---------------------------------------------------------------------------
\tikzset{
  >={Stealth[length=2.1mm,width=1.5mm]},
  base/.style={font=\sffamily\scriptsize,text=tinta,align=center},
  caja/.style={base,draw=#1,fill=#1!7,rounded corners=2pt,line width=0.6pt,inner sep=4pt,
    minimum height=7mm},
  caja/.default=tenue,
  cajallena/.style={base,draw=none,fill=#1,text=white,rounded corners=2pt,inner sep=4pt,
    minimum height=7mm,font=\sffamily\scriptsize\bfseries},
  flecha/.style={->,line width=0.6pt,draw=tinta2},
  flechagruesa/.style={->,line width=1.1pt,draw=tinta2},
  guia/.style={line width=0.5pt,draw=eje},
  nota/.style={font=\sffamily\fontsize{6.4}{7.6}\selectfont,text=tinta2,align=center},
  titulo/.style={font=\sffamily\scriptsize\bfseries,text=tinta},
  capa/.style={rounded corners=3pt,fill=#1!5,draw=#1!45,line width=0.5pt},
}

% Listados (registro de corridas)
\lstdefinestyle{registro}{basicstyle=\ttfamily\scriptsize,columns=fullflexible,keepspaces=true,
  frame=none,backgroundcolor=\color{grisfondo},xleftmargin=6pt,xrightmargin=6pt,aboveskip=6pt,
  belowskip=6pt,showstringspaces=false,stringstyle=\color{qazuloscuro},
  literate={á}{{\'a}}1 {é}{{\'e}}1 {í}{{\'i}}1 {ó}{{\'o}}1 {ú}{{\'u}}1 {ñ}{{\~n}}1}

% ---------------------------------------------------------------------------
% Notación
% ---------------------------------------------------------------------------
\newcommand{\Q}{\mathbb{Q}}
\renewcommand{\P}{\mathbb{P}}
\newcommand{\E}{\mathbb{E}}
\newcommand{\dd}{\,\mathrm{d}}
\newcommand{\CVaR}{\operatorname{CVaR}}
\newcommand{\VaR}{\operatorname{VaR}}
\newcommand{\sintetico}{\textsf{\scriptsize\color{tenue}[datos sintéticos]}}
\newcommand{\grafica}[4][0.66]{%
  \begin{figure}[H]\centering
  \includegraphics[width=#1\linewidth]{figuras/#2}
  \caption{#3}\label{#4}
  \end{figure}}
\newcommand{\figref}[1]{figura~\ref{#1}}
\newcommand{\Figref}[1]{Figura~\ref{#1}}
\newcommand{\secref}[1]{sección~\ref{#1}}
\newcommand{\tabref}[1]{tabla~\ref{#1}}


\begin{document}

% ---------------------------------------------------------------------------
% Portada
% ---------------------------------------------------------------------------
\begin{titlepage}
\thispagestyle{empty}
\begin{tikzpicture}[remember picture,overlay]
  \fill[qazul] (current page.north west) rectangle ([yshift=-0.32in]current page.north east);
\end{tikzpicture}
\vspace*{0.35in}

{\sffamily\small\bfseries\color{qazul}QUANTILEFLOW\hfill
\color{tinta2}\mdseries Versión de investigación \textperiodcentered{} 28 de septiembre de 2026}

\vspace{0.9in}
{\sffamily\fontsize{27}{32}\selectfont\bfseries\color{tinta}
Seguimiento de distribuciones\\[2pt] implícitas y gestión diaria\\[2pt] de posiciones\par}

\vspace{0.22in}
{\sffamily\large\color{tinta2}Guía técnica del proceso: de las cotizaciones de opciones\\
a la decisión sobre la posición\par}

\vspace{0.55in}
\begin{center}
\begin{tikzpicture}[
  capa/.style={rounded corners=4pt,minimum width=4.35cm,minimum height=2.2cm,align=center,
    font=\sffamily,text=white},
  sub/.style={font=\sffamily\footnotesize,text=white,align=center}]
\node[capa,fill=qazul] (q) {\Large\bfseries Medición $\mathbb{Q}$\\[3pt]
  \footnotesize lo que descuentan\\ \footnotesize las opciones};
\node[capa,fill=qnaranja,right=0.55cm of q] (p) {\Large\bfseries Predicción $\mathbb{P}$\\[3pt]
  \footnotesize lo que se espera\\ \footnotesize que ocurra};
\node[capa,fill=qaqua,right=0.55cm of p] (a) {\Large\bfseries Acción\\[3pt]
  \footnotesize mantener o\\ \footnotesize modificar la posición};
\draw[-{Stealth[length=3mm,width=2.4mm]},line width=1.2pt,tinta2] (q) -- (p);
\draw[-{Stealth[length=3mm,width=2.4mm]},line width=1.2pt,tinta2] (p) -- (a);
\node[font=\sffamily\small,text=tinta2,below=4mm of p] {tres capas separadas, cada una con su propia evidencia};
\end{tikzpicture}
\end{center}

\vfill
\begin{tcolorbox}[enhanced,boxrule=0pt,arc=0pt,colback=grisfondo,borderline west={2.2pt}{0pt}{qazul},
  left=10pt,right=10pt,top=8pt,bottom=8pt]
\small\sffamily
\textbf{Alcance.} Este documento explica, etapa por etapa, la arquitectura y el programa de
experimentos de un sistema independiente que mide lo que descuentan las opciones, evalúa si esa
información mejora los pronósticos de rendimiento y riesgo, y compara mantener o modificar una posición.
\textbf{No presenta resultados predictivos ni modelos entrenados.} Todas las gráficas numéricas se
construyen con datos \emph{sintéticos} y semillas fijas (apéndice~\ref{ap:reproducibilidad}); sirven para
explicar el método, no para sugerir su desempeño. La sección~\ref{sec:estado} resume el estado del proyecto
y las primeras cifras reales de calidad de datos y de operación, que no dicen nada sobre la señal.
\end{tcolorbox}
\end{titlepage}

% ---------------------------------------------------------------------------
% Índice
% ---------------------------------------------------------------------------
\pagenumbering{roman}
\setcounter{page}{2}
{\hypersetup{linkcolor=tinta}
\renewcommand{\contentsname}{Contenido}
\linespread{1.0}\selectfont
\tableofcontents}
\clearpage
\pagenumbering{arabic}

% ============================================================================
\section*{Resumen ejecutivo}
\addcontentsline{toc}{section}{Resumen ejecutivo}
\markboth{Resumen ejecutivo}{}
% ============================================================================

Cada sesión, el sistema hace cinco cosas en un orden fijo. Primero \emph{mide} lo que descuentan
las opciones: reconstruye, a partir de cotizaciones de compra y venta, la distribución implícita
$\Q$ del rendimiento a un plazo constante, con bandas que reflejan los spreads. Después
\emph{describe el cambio} respecto de la sesión anterior y separa la parte que ya explicaban el
precio, la volatilidad y el mercado. Luego \emph{pronostica} la distribución física $\P$ de los
rendimientos posteriores al instante de decisión y comprueba si la información de opciones la
mejora. Con ello \emph{decide}, comparando mantener la posición actual con modificarla bajo
incertidumbre, costos y riesgo conjunto de la cartera. Por último \emph{registra} todo lo necesario
para reproducir la corrida.

\begin{ideaclave}[La decisión arquitectónica central]
Mantener separadas tres capas: la \textbf{medición $\Q$}, la \textbf{predicción $\P$} y la
\textbf{acción sobre la posición}. $\Q$ entra a los modelos como variable, nunca como probabilidad
física por decreto; $\P$ se estima con rendimientos realizados y se calibra fuera de muestra; la
acción pasa siempre por costos, riesgo y una regla explícita de abstención.
\end{ideaclave}

\paragraph{Hipótesis que se quiere contrastar} Las \emph{innovaciones} en la forma de las
distribuciones implícitas, una vez descontada la información ya observable en precio, volatilidad y
mercado, pueden mejorar los pronósticos de riesgo o de rendimiento. Es una hipótesis que hay que
contrastar, no una ventaja demostrada: el documento describe cómo medirla y qué resultados la
refutarían. La primera pregunta que conviene responder es si la innovación de las colas, medida con
incertidumbre y relativa al mercado, mejora la gestión de una posición existente.

\paragraph{Qué contiene y qué no} Una arquitectura, un conjunto de métodos y un protocolo de
evidencia. No contiene resultados predictivos ni modelos entrenados. Las gráficas se construyen con un
mundo sintético de parámetros conocidos (apéndice~\ref{ap:reproducibilidad}); su función es mostrar
qué calcula cada etapa y qué puede salir mal, no anticipar el desempeño.

\paragraph{Estado al 28 de septiembre de 2026} Desde que se escribió esta propuesta se construyó la captura
diaria con datos reales de Alpaca, y quedó instalada y verificada en un repositorio privado. La
\secref{sec:estado} resume qué está hecho, qué cambió respecto de la propuesta y qué falta. Todavía no hay
resultados sobre la capacidad predictiva.

\paragraph{Cómo leer este documento} Cada sección corresponde a una etapa del proceso y empieza con
una ficha de \emph{entradas, proceso, salidas y controles}. Cinco tipos de recuadro se repiten:

\begin{center}\small
\begin{tabularx}{\linewidth}{@{}l>{\RaggedRight\arraybackslash}X@{}}
\iconoidea\ \textsf{\textbf{Idea clave}} & el concepto que conviene retener de la etapa.\\
\iconoalerta\ \textsf{\textbf{Qué no concluir}} & interpretaciones tentadoras que el método no permite.\\
\iconocheck\ \textsf{\textbf{Criterio para avanzar}} & condición verificable antes de pasar a la siguiente etapa.\\
\iconopendiente\ \textsf{\textbf{Pendiente de definir}} & decisiones del usuario que parametrizan el sistema.\\
\icononota\ \textsf{\textbf{Nota técnica}} & detalle de implementación o advertencia matemática.\\
\end{tabularx}
\end{center}

La \figref{fig:mapa} resume el recorrido completo. Los colores se mantienen en todos los diagramas:
azul para la medición $\Q$, naranja para la predicción $\P$ y verde para la acción.

\begin{figure}[H]
\centering
\resizebox{\linewidth}{!}{% Mapa del proceso: medición Q, predicción P y acción, con capas transversales.
\begin{tikzpicture}[
  n/.style={caja=#1,text width=2.28cm,minimum height=11.5mm},
  etq/.style={cajallena=#1,text width=1.62cm,minimum height=11.5mm,font=\sffamily\tiny\bfseries},
  node distance=4.2mm and 4.2mm]
% --- Medición Q -------------------------------------------------------------
\node[etq=qazul] (lq) {MEDICIÓN\\ Q};
\node[n=qazul,right=of lq] (q1) {Snapshots bid/ask\\ sincronizados};
\node[n=qazul,right=of q1] (q2) {Calidad, forward\\ y desamericanización};
\node[n=qazul,right=of q2] (q3) {Superficie sin\\ arbitraje (SSVI)};
\node[n=qazul,right=of q3] (q4) {Densidad, cuantiles\\ y bandas};
\node[n=qazul,right=of q4] (q5) {Transporte, factores\\ e innovación};
\foreach \a/\b in {q1/q2,q2/q3,q3/q4,q4/q5} \draw[flecha] (\a) -- (\b);
% --- Predicción P -----------------------------------------------------------
\node[etq=qnaranja,below=10mm of lq] (lp) {PREDICCIÓN\\ P};
\node[n=qnaranja,right=of lp] (p1) {Información base:\\ precio, vol., mercado};
\node[n=qnaranja,right=of p1] (p2) {Modelos P por\\ objetivo y horizonte};
\node[n=qnaranja,right=of p2] (p3) {Combinación y\\ calibración temporal};
\node[n=qnaranja,right=of p3] (p4) {Pronósticos P\\ calibrados};
\node[n=qnaranja,right=of p4] (p5) {Regímenes y\\ confianza medible};
\foreach \a/\b in {p1/p2,p2/p3,p3/p4,p4/p5} \draw[flecha] (\a) -- (\b);
% --- Acción -----------------------------------------------------------------
\node[etq=qaqua,below=10mm of lp] (la) {ACCIÓN};
\node[n=qaqua,right=of la] (a1) {Escenarios\\ conjuntos P};
\node[n=qaqua,right=of a1] (a2) {Costo, riesgo y\\ rendimiento (CVaR)};
\node[n=qaqua,right=of a2] (a3) {Zona de no\\ operación};
\node[n=qaqua,right=of a3] (a4) {Recomendación\\ o abstención};
\node[n=qaqua,right=of a4] (a5) {Paper, registro y\\ reconciliación};
\foreach \a/\b in {a1/a2,a2/a3,a3/a4,a4/a5} \draw[flecha] (\a) -- (\b);
% --- Conexiones entre capas -------------------------------------------------
\draw[flecha] (q5.south) -- ++(0,-0.32) -| (p2.north);
\draw[flecha] (p4.south) -- ++(0,-0.32) -| (a1.north);
\draw[flecha] (p5.south) -- ++(0,-0.6) -| (a4.north);
% --- Capas transversales ----------------------------------------------------
\node[caja=tenue,below=7mm of a1.south west,anchor=north west,minimum width=14.85cm,text width=14.5cm,
  minimum height=7mm] (reg) {\textbf{Registro reproducible}: hashes de entrada, versión de modelo, parámetros,
  hora de corte y motivo de abstención en cada corrida};
\node[caja=tenue,below=2mm of reg,minimum width=14.85cm,text width=14.5cm,minimum height=7mm] (conf)
  {\textbf{Confianza}: calidad de datos $\cdot$ sensibilidad del ajuste $\cdot$ incertidumbre predictiva
  $\cdot$ estabilidad de la decisión};
\end{tikzpicture}
}
\caption{Mapa del proceso diario. Las flechas entre capas son deliberadamente estrechas: de $\Q$ a
$\P$ solo pasan variables candidatas, y de $\P$ a la acción solo pasan distribuciones físicas
calibradas. El registro y la evaluación de confianza atraviesan todas las etapas.}
\label{fig:mapa}
\end{figure}

% ============================================================================
\clearpage
\section{Producto y pregunta de investigación}\label{sec:producto}
% ============================================================================

\ficha{Cadenas de opciones con bid/ask, precio del subyacente, posiciones actuales, calendario de
resultados y dividendos.}
{Medir $\Q$, describir su cambio, pronosticar $\P$ y comparar mantener con modificar.}
{Recomendación con bandas de incertidumbre, o abstención motivada; registro reproducible.}
{Separación estricta entre $\Q$, $\P$ y acción; ninguna decisión sin calibración ni riesgo conjunto.}

\subsection{Qué hace el sistema en cada sesión}

El producto responde, al inicio de cada sesión, una pregunta operativa concreta: \emph{dado lo que
descuentan hoy las opciones y lo que ya se sabe por el precio, ¿conviene mantener la posición o
modificarla?} Para llegar a esa respuesta encadena cinco pasos, cada uno con su propia salida
verificable:

\begin{enumerate}
\item \textbf{Reconstruir lo que descuentan las opciones.} A partir de un \emph{snapshot} de
cotizaciones se obtiene una superficie de volatilidad sin arbitraje estático y, de ella, la
distribución implícita $\Q_t$ del rendimiento a plazo constante, con bandas que reflejan el ancho de
los spreads.
\item \textbf{Cuantificar cómo cambió.} Se compara $\Q_t$ con $\Q_{t-1}$ mediante distancias entre
funciones cuantil, cambios firmados por percentil y factores funcionales filtrados.
\item \textbf{Estimar si esa información mejora el pronóstico.} Los cambios --- en especial su parte
no explicada por la información base --- entran como variables candidatas a modelos que pronostican la
distribución física $\P$ de rendimientos futuros, evaluados fuera de muestra.
\item \textbf{Comparar mantener con modificar.} Con escenarios conjuntos de la cartera se evalúa el
rendimiento esperado, el CVaR y los costos de cada alternativa; los cambios pequeños que no superan
costos e incertidumbre se descartan.
\item \textbf{Registrar.} Cada resultado conserva las entradas, versiones, parámetros, hora de corte y
motivo de abstención, de modo que la corrida pueda repetirse exactamente.
\end{enumerate}

\subsection{Hipótesis principal y cómo se contrasta}\label{sec:hipotesis}

Sea $Y_{t,h}$ una variable objetivo medida entre el instante de decisión $t$ y $t+h$ (rendimiento,
pérdida o volatilidad realizada), $\mathcal{I}^{base}_t$ la información ya observable en precio,
volatilidad y mercado, y $\epsilon_t$ la innovación de la distribución implícita definida en la
\secref{sec:innovacion}. Con una regla de puntuación propia $S$ (menor es mejor) y dos pronósticos de
la \emph{misma familia} de modelos, uno con $\mathcal{I}^{base}_t$ y otro con
$\mathcal{I}^{base}_t\cup\{\epsilon_t\}$, se define el diferencial de pérdida
\begin{equation}
d_{t,h}=S\bigl(\widehat F^{\,base}_{t,h},\,Y_{t,h}\bigr)-S\bigl(\widehat F^{\,base+\epsilon}_{t,h},\,Y_{t,h}\bigr).
\label{eq:diferencial}
\end{equation}
La hipótesis de trabajo es $H_1:\ \E[d_{t,h}]>0$ fuera de muestra, frente a $H_0:\ \E[d_{t,h}]\le 0$.
Se evalúa por separado para cada objetivo y horizonte, con intervalos que respetan la dependencia
temporal (\secref{sec:evidencia}).

\begin{noconcluir}
\begin{itemize}
\item Una mejora en el pronóstico del \emph{riesgo} no implica una mejora en el \emph{signo} del
rendimiento, ni al revés; ambas conclusiones son útiles y se reportan por separado.
\item Una mejora en el pronóstico no implica una mejor política de posición: puede no pagar sus
costos. Son dos pruebas distintas (\secref{sec:dos_pruebas}).
\item La hipótesis no presupone una dirección rentable; esa dirección, si existe, se mide.
\end{itemize}
\end{noconcluir}

\subsection{Alcance de la primera versión}

La primera versión se limita a SPY y entre tres y cinco acciones líquidas, con posiciones en
acciones o ETF y sin apalancamiento como supuesto inicial de investigación. Gestionar opciones como
posición exigiría una fase específica de valoración y escenarios conjuntos (\secref{sec:opciones}).

\begin{table}[H]
\centering\small
\caption{Alcance de la primera versión.}
\label{tab:alcance}
\begin{tabularx}{\linewidth}{@{}l l Y@{}}
\toprule
\encab{Elemento} & \encab{Estado} & \encab{Comentario}\\
\midrule
Universo & Propuesto & SPY y 3--5 acciones líquidas; lista definitiva pendiente.\\
Instrumentos gestionados & Supuesto inicial & Acciones y ETF; opciones como fase aparte.\\
Apalancamiento & Supuesto de investigación & Sin apalancamiento en la primera versión.\\
Capital de referencia & Pendiente & 50,000; falta definir la moneda.\\
Posiciones actuales & Pendiente & Necesarias para evaluar ``mantener frente a modificar''.\\
Horizonte de decisión & Pendiente & Candidatos: 1, 5 y 20 sesiones.\\
Presupuesto de datos & Pendiente & Condiciona la fuente y la historia disponible.\\
\bottomrule
\end{tabularx}
\end{table}

\begin{pendiente}
Universo definitivo, moneda de los 50,000, posiciones, horizonte y presupuesto de datos. No son
necesarios para \emph{diseñar} el motor, pero sí para \emph{parametrizar} su uso: sin ellos no se
pueden fijar límites de concentración, costos ni el benchmark comparable.
\end{pendiente}

\subsection{Tres horizontes que no deben mezclarse}\label{sec:horizontes}

El sistema maneja tres escalas de tiempo distintas (\figref{fig:horizontes}). La \textbf{frecuencia
de observación} es cada sesión, incluidas las ventanas de apertura. El \textbf{horizonte de la
distribución implícita} es el plazo de las opciones que se interpolan; inicialmente 30 días
calendario y constante, de modo que en $t+1$ se sigue midiendo a 30 días y no a 29. El
\textbf{horizonte de decisión y pronóstico} es aquel sobre el que se evalúa el rendimiento de una
decisión, por ejemplo 1, 5 y 20 sesiones.

\begin{figure}[H]
\centering
\resizebox{\linewidth}{!}{% Tres horizontes: observación, distribución implícita y decisión/pronóstico.
\begin{tikzpicture}[x=0.36cm,y=1cm,nota/.style={font=\sffamily\scriptsize,text=tinta2,align=left}]
% Etiquetas de fila
\node[titulo,anchor=east,align=right] at (-0.8,0) {Frecuencia de\\ observación};
\node[titulo,anchor=east,align=right] at (-0.8,-1.35) {Horizonte de la\\ distribución $\Q$};
\node[titulo,anchor=east,align=right] at (-0.8,-2.95) {Horizonte de decisión\\ y pronóstico $\P$};
% Fines de semana sombreados
\foreach \s in {5,12,19,26} \fill[grisclaro] (\s-0.5,0.45) rectangle (\s+1.5,-3.75);
% Instante de decisión
\draw[tinta,line width=0.6pt] (0,0.55) -- (0,-3.75);
\node[nota,anchor=south] at (0,0.55) {instante de decisión $t$};
% Observación: cada sesión hábil
\foreach \d in {0,...,30}{%
  \pgfmathtruncatemacro{\w}{mod(\d,7)}%
  \ifnum\w<5 \fill[tenue] (\d,0) circle (2.1pt);\fi}
\node[nota,anchor=west] at (31.2,0) {cada sesión: apertura\\ y cierre, sin operar\\ necesariamente};
% Distribución implícita a plazo constante
\draw[line width=6pt,qazul!55] (0,-1.15) -- (30,-1.15);
\draw[line width=6pt,qazul!22] (1,-1.55) -- (31,-1.55);
\node[nota,anchor=west] at (31.2,-1.15) {$\tau=30$ días calendario};
\node[nota,anchor=west] at (31.2,-1.62) {en $t+1$ el plazo sigue\\ siendo 30 días (interpolado)};
% Horizontes de decisión
\draw[flecha,draw=qnaranja,line width=0.9pt] (0,-2.45) -- (1,-2.45);
\node[nota,anchor=west] at (1.3,-2.45) {1 sesión};
\draw[flecha,draw=qnaranja,line width=0.9pt] (0,-2.95) -- (7,-2.95);
\node[nota,anchor=west] at (7.3,-2.95) {5 sesiones};
\draw[flecha,draw=qnaranja,line width=0.9pt] (0,-3.45) -- (28,-3.45);
\node[nota,anchor=west] at (28.3,-3.45) {20 sesiones};
% Eje de tiempo
\draw[guia] (-0.5,-3.9) -- (31,-3.9);
\foreach \d in {0,7,14,21,28}{\draw[guia] (\d,-3.9) -- (\d,-4.0);
  \node[nota,anchor=north] at (\d,-4.02) {\d};}
\node[nota,anchor=north] at (15,-4.35) {días calendario desde $t$ (sombreado: fines de semana)};
\end{tikzpicture}
}
\caption{Los tres horizontes. Observar cada sesión no obliga a operar cada sesión; la distribución
implícita se mide siempre a 30 días calendario; los pronósticos se evalúan a 1, 5 y 20 sesiones.}
\label{fig:horizontes}
\end{figure}

La distinción importa porque una distribución a 30 días \emph{no se convierte por simple escalado} en
una distribución a una sesión. La varianza no se acumula de manera uniforme: los fines de semana
aportan poco, un día de resultados trimestrales aporta mucho, y la estructura temporal de
volatilidades rara vez es plana. Las figuras~\ref{fig:acumulacion} y~\ref{fig:riesgo_sesion}
lo ilustran con un calendario sintético de 30 días que incluye un anuncio de resultados.

\grafica{s1_acumulacion_varianza}{Acumulación de varianza en 30 días calendario (datos sintéticos).
La escalera muestra una acumulación plausible: casi nada en fines de semana y un salto el día de
resultados. La recta es lo que supone un escalado lineal en días calendario.}{fig:acumulacion}

\grafica{s1_riesgo_sesion}{Desviación estándar de una sesión (datos sintéticos). Las dos primeras
barras son los valores del mundo sintético para una sesión normal y para la sesión de resultados; las
dos últimas, lo que se obtiene escalando la varianza a 30 días por días calendario o por sesiones.
Ningún escalado recupera el riesgo de la sesión de resultados.}{fig:riesgo_sesion}

\begin{ideaclave}
El plazo de $\Q$ describe qué se mide; el horizonte de decisión describe qué se evalúa. Pasar de uno a
otro exige un modelo explícito (calendario de eventos, estructura temporal, supuestos sobre fines de
semana), nunca una regla de raíz del tiempo aplicada por defecto.
\end{ideaclave}

% ============================================================================
\clearpage
\section{Datos y reconstrucción: la primera capa de robustez}\label{sec:datos}
% ============================================================================

\ficha{Cotizaciones bid/ask de cadenas completas, subyacente sincronizado, tasas, dividendos y eventos
corporativos, con sellos de tiempo de mercado y de recepción.}
{Filtrar, estimar forward y descuento, convertir opciones americanas, ajustar una superficie sin
arbitraje y extraer la densidad.}
{Distribución $\Q_t$ a plazo constante, cuantiles, bandas de sensibilidad y soporte identificado.}
{Nada se sobrescribe; densidad no negativa, masa próxima a uno y media igual al forward; colas sin
soporte marcadas.}

Todo lo que viene después depende de esta etapa. Un error aquí no produce una señal débil: produce una
señal falsa con apariencia de precisión. Por eso la reconstrucción se diseña como una cadena de pasos
auditables (\figref{fig:flujo_q}), cada uno con un control explícito.

\begin{figure}[H]
\centering
\resizebox{\linewidth}{!}{% Reconstrucción de Q: de cotizaciones a cuantiles con bandas.
\begin{tikzpicture}[
  n/.style={caja=qazul,text width=3.15cm,minimum height=15mm},
  chk/.style={nota,text width=3.15cm,align=center,text=tinta2},
  node distance=6mm and 5.5mm]
\node[n] (s1) {\textbf{1. Snapshot inmutable}\\ bid, ask, tamaños, sellos de tiempo, subyacente};
\node[n,right=of s1] (s2) {\textbf{2. Filtros de calidad}\\ cruces, bid nulo, spreads, datos rancios};
\node[n,right=of s2] (s3) {\textbf{3. Forward y descuento}\\ paridad, tasas y dividendos};
\node[n,right=of s3] (s4) {\textbf{4. Desamericanización}\\ equivalente europeo o exclusión};
\node[n,below=12mm of s4] (s5) {\textbf{5. Ajuste en precios}\\ SSVI restringido, pérdida por banda};
\node[n,left=of s5] (s6) {\textbf{6. Contraste}\\ ajuste monótono y convexo};
\node[n,left=of s6] (s7) {\textbf{7. Densidad}\\ Breeden--Litzenberger y controles};
\node[n,left=of s7] (s8) {\textbf{8. Cuantiles y bandas}\\ plazo constante, colas marcadas};
\foreach \a/\b in {s1/s2,s2/s3,s3/s4,s5/s6,s6/s7,s7/s8} \draw[flecha] (\a) -- (\b);
\draw[flecha] (s4) -- (s5);
\node[chk,below=0.5mm of s1] {sin sobrescribir nada};
\node[chk,below=0.5mm of s2] {motivo de cada exclusión};
\node[chk,below=0.5mm of s3] {forward coherente con paridad};
\node[chk,above=0.5mm of s4.north] {error de conversión acotado};
\node[chk,below=0.5mm of s5] {restricciones por construcción};
\node[chk,below=0.5mm of s6] {discrepancia = riesgo de modelo};
\node[chk,below=0.5mm of s7] {$q\ge 0$, masa $\approx 1$, media $=F$};
\node[chk,below=0.5mm of s8] {sensibilidad, no confianza};
\end{tikzpicture}
}
\caption{Cadena de reconstrucción de $\Q$. Debajo de cada paso, el control que debe superar antes de
alimentar al siguiente.}
\label{fig:flujo_q}
\end{figure}

\subsection{Snapshots inmutables y datos puntuales en el tiempo}\label{sec:snapshots}

Cada captura se guarda como un \emph{snapshot inmutable}: nunca se corrige en sitio; una corrección es
un registro nuevo con su propia hora de publicación. Esto permite responder, para cualquier instante
pasado, la única pregunta válida en un backtest: \emph{¿qué se sabía en ese momento?}
La \tabref{tab:campos} resume los campos mínimos y por qué importan.

\begin{table}[H]
\centering\small
\caption{Campos mínimos de cada snapshot.}
\label{tab:campos}
\begin{tabularx}{\linewidth}{@{}l Y@{}}
\toprule
\encab{Campo} & \encab{Por qué importa}\\
\midrule
Bid, ask y tamaños & La banda $[b,a]$ es la observación; el mid es solo un punto interior.\\
Hora de mercado y de recepción & Separan cuándo ocurre un dato de cuándo se conoce (\figref{fig:puntual}).\\
Subyacente sincronizado & Un desfase de segundos contamina el forward y la volatilidad implícita.\\
Strike, vencimiento, estilo, multiplicador & Definen el contrato; el estilo americano exige conversión.\\
Dividendos, tasas y eventos corporativos & Determinan forward, descuento y ajustes de contrato.\\
Revisiones y open interest & Se guardan con su hora efectiva de disponibilidad, no con la del dato.\\
\bottomrule
\end{tabularx}
\end{table}

\begin{figure}[H]
\centering
\resizebox{0.96\linewidth}{!}{% Datos puntuales en el tiempo: cuándo ocurre un dato frente a cuándo se conoce.
\begin{tikzpicture}[x=2.05cm,y=1cm]
% Región posterior al corte
\fill[qazul!6] (3.75,2.75) rectangle (6.3,-0.25);
% Pistas: hora de mercado (arriba) y hora de recepción (abajo); 06:00 -> x = 0
\draw[guia,line width=0.7pt] (0,2.2) -- (6.3,2.2);
\draw[guia,line width=0.7pt] (0,0) -- (6.3,0);
\node[titulo,anchor=east,align=right] at (-0.08,2.2) {ocurre\\[-1pt] {\mdseries\tiny(hora de mercado)}};
\node[titulo,anchor=east,align=right] at (-0.08,0) {se conoce\\[-1pt] {\mdseries\tiny(hora de recepción)}};
\foreach \h/\e in {0/06:00,1/07:00,2/08:00,3/09:00,4/10:00,5/11:00,6/12:00}{
  \draw[guia] (\h,2.2) -- (\h,2.28); \draw[guia] (\h,0) -- (\h,-0.08);
  \node[nota,anchor=north] at (\h,-0.1) {\e};}
% Corte de decisión
\draw[qazul,line width=0.9pt] (3.75,2.95) -- (3.75,-0.45);
\node[nota,anchor=south,text=qazuloscuro] at (3.75,2.95) {\textbf{corte de decisión: 09:45 ET}};
\node[nota,anchor=north west,text=qazuloscuro] at (3.8,-0.45) {recibido después del corte: no existe para esa decisión};
% Open interest del cierre anterior, publicado antes de la apertura
\draw[flecha,draw=tenue] (0.08,2.2) -- (1.0,0);
\node[nota,anchor=south west] at (0.0,2.3) {open interest del cierre de $t-1$};
% Fuente con 15 minutos de retraso: dato de 09:30 recibido a las 09:45
\draw[flecha,draw=qnaranja,line width=0.8pt] (3.5,2.2) -- (3.74,0);
\node[nota,anchor=east,align=right] at (3.52,1.05) {fuente con 15 min de retraso:\\ dato de 09:30 recibido 09:45};
% Cotización en tiempo real
\draw[flecha,draw=qazul,line width=0.8pt] (3.71,2.2) -- (3.715,0);
\node[nota,anchor=south east] at (3.7,2.3) {cotización 09:44:58};
% Revisión posterior de un dato anterior
\draw[flecha,draw=critico,line width=0.8pt] (3.6,2.2) -- (5.08,0);
\node[nota,anchor=west,align=left] at (4.55,1.35) {revisión publicada a las 11:05\\ de un dato de las 09:36};
\end{tikzpicture}
}
\caption{Hora en que ocurre un dato frente a hora en que se conoce. En el corte de las 09:45 solo son
utilizables las flechas que llegan a la línea inferior antes del corte: la cotización en tiempo real,
el open interest del cierre anterior y, con una fuente retrasada, lo ocurrido a las 09:30. La revisión
publicada a las 11:05 no existe para esa decisión aunque se refiera a un dato anterior.}
\label{fig:puntual}
\end{figure}

\subsection{Contrato de datos y fuentes candidatas}\label{sec:contrato}

Antes de adquirir historia extensa conviene validar una muestra contra un contrato explícito. El
primer contrato debe demostrar, como mínimo:

\begin{itemize}
\item cadenas \emph{completas}, incluidos contratos ya vencidos, y paginación sin huecos;
\item cobertura histórica \emph{a la hora de decisión}, no solo al cierre;
\item acceso a cotizaciones de compra y venta. Las barras OHLC de operaciones no sustituyen una cadena
contemporánea de cotizaciones: una operación puede ser antigua, aislada o fuera del spread;
\item sellos de tiempo suficientes para reconstruir la sincronía con el subyacente.
\end{itemize}

\begin{table}[H]
\centering\small
\caption{Fuentes candidatas según la documentación citada en la propuesta. Cobertura, licencia y precio
deben cotizarse; no se asumen.}
\label{tab:fuentes}
\begin{tabularx}{\linewidth}{@{}l Y Y@{}}
\toprule
\encab{Fuente} & \encab{Lo que documenta} & \encab{Uso previsto}\\
\midrule
Alpaca & Históricos de opciones desde febrero de 2024; el feed indicativo modifica cotizaciones
\citep{alpaca}. & Explorar la integración. No se asume suficiente para estudiar cambios diarios pequeños
de distribución.\\
Cboe DataShop & Snapshots NBBO por intervalos con historia desde enero de 2012; entrega intradía estándar
con 15 minutos de retraso \citep{cboe}. & Investigación histórica. Una decisión en apertura exige una
fuente con latencia adecuada o retrasar explícitamente la decisión.\\
\bottomrule
\end{tabularx}
\end{table}

\begin{criterio}
Cadenas reconstruibles en el instante de decisión, con cobertura y costo conocidos, validadas sobre una
muestra antes de comprar historia extensa.
\end{criterio}

\subsection{Cadena cruda: controles, paridad y asimetría observada}\label{sec:cadena_cruda}

Las primeras preguntas del plan de trabajo (\texttt{docs/plan\_de\_trabajo.md}) se responden con la cadena
\emph{cruda}: si el encarecimiento relativo de calls frente a puts, o el residuo de paridad del mismo
strike, anticipa algo. Estas medidas se calculan antes de cualquier ajuste y se guardan como una serie
propia: una superficie sin arbitraje impone la paridad y borraría justamente el residuo que se quiere
estudiar. El módulo \texttt{quantileflow/cadenas.py} implementa esta sección. Su clase \texttt{Captura} es
la versión mínima del contrato de datos: una fila por actualización de cada contrato, con bid, ask,
tamaños y sello de tiempo, más el subyacente sincronizado, la tasa, los dividendos, el estilo de
ejercicio, la hora de corte y la fecha. Sus arreglos son de solo lectura.

\paragraph{Controles por fila} De cada contrato cuenta solo la última actualización anterior o igual al
corte; si esa actualización falla un control, el contrato queda fuera y no se recurre a una cotización más
vieja. Cada fila excluida conserva \emph{todos} sus motivos (\tabref{tab:motivos}), de modo que el informe
diario dice qué se excluyó y por qué. Los umbrales son provisionales hasta auditar una muestra real en la
fase~0. La captura recibe además alertas si el sello del subyacente se aleja del corte o cae en la
ventana de apertura, cuando el valor de un índice puede incluir componentes que aún no han abierto.

\begin{table}[H]
\centering\small
\caption{Motivos de exclusión de una fila de la cadena cruda.}
\label{tab:motivos}
\begin{tabularx}{\linewidth}{@{}l Y@{}}
\toprule
\encab{Motivo} & \encab{Condición y razón}\\
\midrule
\texttt{posterior\_al\_corte} & Sello posterior al corte: no se conocía al decidir.\\
\texttt{reemplazada} & Hay una actualización posterior del mismo contrato antes del corte.\\
\texttt{antes\_de\_apertura} & Cotización de preapertura o de la sesión anterior.\\
\texttt{ventana\_de\_apertura} & Primeros minutos de la sesión (cinco por omisión): cotizaciones inestables.\\
\texttt{desfasada} & Más antigua que la edad máxima al corte (60~s): no refleja el subyacente actual.\\
\texttt{sin\_bid} & Bid nulo: el precio solo queda acotado por arriba.\\
\texttt{cruzada}, \texttt{bloqueada} & Bid mayor que ask, o igual: dato inconsistente o de fuentes
desincronizadas.\\
\texttt{sin\_tamano} & Tamaño menor que el mínimo en bid o en ask.\\
\texttt{spread\_ancho} & Spread mayor que el 50\% del mid.\\
\texttt{fuera\_de\_cotas} & Viola cotas sin modelo: call por encima del subyacente, put por encima del
strike o, si es americana, ask menor que el ejercicio inmediato.\\
\bottomrule
\end{tabularx}
\end{table}

\paragraph{Residuo de paridad fuera de muestra} Para cada par $j$ de call y put válidos del mismo strike,
\begin{equation}
r_j=(C_j-P_j)-D_{-j}\,(F_{-j}-K_j),\qquad
\bigl[\underline r_j,\ \overline r_j\bigr]=\bigl[C^{b}_j-P^{a}_j,\ C^{a}_j-P^{b}_j\bigr]-D_{-j}\,(F_{-j}-K_j),
\label{eq:residuo_paridad}
\end{equation}
donde $C_j$ y $P_j$ son mids, los superíndices $b$ y $a$ indican bid y ask, y $(F_{-j},D_{-j})$ se estiman
por mínimos cuadrados ponderados de $C-P$ sobre $K$ con los \emph{demás} pares, con pesos $1/\ell^2$ y
$\ell_j=(C^a_j-C^b_j)+(P^a_j-P^b_j)$ el ancho de la banda. Ninguna cotización evalúa su propia paridad, de
modo que la medida no es circular. Un par sale de la estimación del forward de los demás si su residuo
supera tres medios anchos y cinco desviaciones robustas del resto, pero conserva su residuo. El residuo se
reporta en unidades de precio y en múltiplos del ancho: si $|r_j|\le\ell_j/2$, la banda contiene el cero y
la diferencia no supera el costo de negociar, así que no se etiqueta como señal. La
\figref{fig:paridad_desfasadas} muestra por qué el control de edad es imprescindible: tres puts de hace
cinco minutos, con el subyacente un 1\% más bajo, dan residuos de $-0.6$, $-1.9$ y $-3.2$ anchos que,
sin ese control, pasarían por señales; los demás pares no se contaminan.


\paragraph{Dividendos y ejercicio anticipado} El forward implícito incorpora los dividendos sin necesidad
de declararlos. Un forward contractual que omite un dividendo en efectivo fabrica, en todos los strikes, un
residuo igual a su valor presente ($-1.2$ en las pruebas, con un dividendo de 1.2). Con opciones americanas
la paridad deja de ser una igualdad: cada cotización se desamericaniza con el árbol binomial y el modelo de
dividendo \emph{escrowed} (\secref{sec:americanas}), valorando la americana y la europea en el mismo
árbol, y el residuo solo se interpreta si supera la banda \emph{más} las primas de ejercicio estimadas y
el subyacente está sincronizado. Tratar como europeas unas americanas con dividendo sesga el forward
implícito (0.2 sobre 98.8 en las pruebas) aunque ningún residuo salga de su banda.

\grafica[0.62]{s2_paridad_desfasadas}{Residuo de paridad del mismo strike, en múltiplos del ancho bid--ask,
con el forward estimado sin el propio par. Los círculos son los pares que superan los controles; los
triángulos, tres puts desfasados que solo aparecerían si no se comprobara la edad de la cotización. Entre
las líneas discontinuas, la diferencia cabe en el spread.}{fig:paridad_desfasadas}

\paragraph{Asimetría call/put} Un call a $+3$\% y un put a $-3$\% no forman un par de paridad: su
comparación mide asimetría de precios de cola. Se usan strikes simétricos en logaritmo, $Fe^{\pm d}$ con
$d=\ln 1.03$, y se compara primero en volatilidad implícita, $\sigma_P(-d)-\sigma_C(d)$, con una banda que
sale de las volatilidades de bid y ask. Las primas brutas no deben compararse con la unidad: con una
sonrisa simétrica en $k$,
\begin{equation}
P\bigl(Fe^{-d}\bigr)=e^{-d}\,C\bigl(Fe^{d}\bigr),\qquad
\varrho_d=\frac{P(Fe^{-d})}{e^{-d}\,C(Fe^{d})}-1,
\label{eq:simetria_pc}
\end{equation}
de modo que $\varrho_d=0$ indica ausencia de asimetría. Como una distancia fija cambia de significado con
el nivel de volatilidad y el plazo, la comparación se repite con deltas comparables,
$\sigma_P(\Delta=-0.25)-\sigma_C(\Delta=0.25)$, usando la delta forward de cada opción. La pendiente local
$\partial\sigma/\partial k$ en $k=0$ se estima por mínimos cuadrados ponderados con $|k|\le5$\%
(\figref{fig:asimetria_sonrisa}). Solo se interpola entre strikes válidos separados a lo sumo 2\% en
$k$; si el tramo necesario tiene un hueco mayor, la medida queda «no identificada».

\grafica[0.62]{s2_asimetria_sonrisa}{Volatilidades implícitas de bid y ask de las cotizaciones OTM válidas y
las tres medidas de asimetría: strikes log-simétricos en $\pm\ln 1.03$, deltas de 25 y pendiente local. En
este mundo sintético, la diferencia en $\pm3$\% es de 5.8 puntos, frente a 5.76 en la sonrisa
verdadera.}{fig:asimetria_sonrisa}

\paragraph{Series separadas y cambios diarios} La fila diaria del informe lleva las medidas observadas
(\texttt{obs\_*}) y, si se ajusta una superficie, las que se leen de ella (\texttt{aj\_*}), nunca
mezcladas: en el ajuste, el residuo de paridad es cero por construcción. Los cambios respecto de la
captura anterior registran días naturales, sesiones transcurridas y tipo de comparación
(apertura--apertura o cierre--apertura); un fin de semana cuenta como tres días y una sesión, no como una
sesión ordinaria.

\begin{noconcluir}
Las pruebas de esta sección usan cadenas sintéticas con cotizaciones desfasadas, dividendos, strikes
ausentes y ventanas de apertura. Verifican que el procesamiento separa efectos de los datos de una posible
señal; no demuestran que el residuo o la asimetría anticipen rendimientos, ni la calidad de datos reales.
\end{noconcluir}

\subsection{De cotizaciones a precios ajustados}\label{sec:ajuste_precios}

\paragraph{Forward y descuento} Para opciones europeas, la paridad put--call
$C(K)-P(K)=D\,(F-K)$ permite estimar $D$ y $F$ por regresión ponderada de $C-P$ sobre $K$
(\secref{sec:cadena_cruda}). Con opciones americanas con dividendos la paridad es solo una desigualdad: se
aplica a los precios desamericanizados, con la prima de ejercicio estimada como parte de la
incertidumbre, y se contrasta con el forward que dan la curva de tasas y los dividendos esperados.

\paragraph{Selección y filtros} Se trabaja con opciones fuera del dinero (puts bajo el forward, calls por
encima), más líquidas y con menor valor intrínseco. Se excluyen, registrando el motivo, las cotizaciones
cruzadas, rancias o con spread anómalo (\tabref{tab:motivos}). Una cotización con bid nulo no se descarta sin más: su ask sigue
siendo una cota superior útil, aunque no identifica la cola (\figref{fig:cotizaciones}).

\grafica{s2_cotizaciones_precio}{Cotizaciones OTM a 30 días en el mundo sintético. Cada segmento vertical
es un intervalo bid--ask; los triángulos son cotizaciones con bid nulo, que solo acotan el precio por
arriba. La curva es el ajuste SSVI en espacio de precios.}{fig:cotizaciones}

\paragraph{Ajuste en espacio de precios} El midpoint no es una verdad exacta: cualquier precio dentro de
$[b_i,a_i]$ es compatible con la cotización. Por eso el residuo de cada contrato es su \emph{distancia a la
banda}, medida en medios spreads $s_i=\max\{(a_i-b_i)/2,\ \text{tick}\}$:
\begin{equation}
\min_{\vartheta\in\Theta}\ \sum_i \omega_i\left[\frac{(b_i-C_i(\vartheta))^+ + (C_i(\vartheta)-a_i)^+}{s_i}\right]^2
+\eta\sum_i \omega_i\left[\frac{C_i(\vartheta)-m_i}{s_i}\right]^2+\mathcal{R}(\vartheta),
\label{eq:perdida_banda}
\end{equation}
donde $C_i(\vartheta)$ es el precio del modelo, $m_i$ el mid, $\eta$ un peso pequeño que solo desempata
entre precios igualmente compatibles y $\mathcal R$ una regularización, por ejemplo hacia el ajuste de la
sesión anterior. Ajustar en precios y no en volatilidades evita dar peso excesivo a las alas, donde un
spread de un tick equivale a varios puntos de volatilidad (\figref{fig:cotizaciones_vol}).

\grafica{s2_cotizaciones_vol}{La misma cadena en volatilidad implícita. Los intervalos se ensanchan en las
alas y las cotizaciones con bid nulo solo dan cotas superiores cada vez más altas: en volatilidad, esas
cotas parecen ``información'' sobre la cola derecha que en realidad no existe.}{fig:cotizaciones_vol}

\grafica[0.6]{s2_perdida_banda}{Pérdida de una cotización con bid 1.00 y ask 1.20. La pérdida cuadrática
al mid penaliza precios perfectamente compatibles; la distancia a la banda es casi nula dentro del
intervalo y crece fuera de él.}{fig:perdida_banda}

\subsection{SSVI y ausencia de arbitraje estático}\label{sec:ssvi}

El modelo base es SSVI \citep{gatheral2014}, que describe la varianza implícita total $w=\sigma^2T$ como
función del log-moneyness $k=\ln(K/F)$ y de la varianza total en el dinero $\theta_T$ de cada vencimiento:
\begin{equation}
w(k,\theta_T)=\frac{\theta_T}{2}\Bigl\{1+\rho\,\varphi(\theta_T)\,k+\sqrt{\bigl(\varphi(\theta_T)k+\rho\bigr)^2+1-\rho^2}\Bigr\},
\qquad \varphi(\theta)=\frac{\eta}{\theta^{\gamma}(1+\theta)^{1-\gamma}}.
\label{eq:ssvi}
\end{equation}
El parámetro $\rho$ controla la asimetría, $\varphi$ la curvatura y el peso de las alas, y $\theta_T$ el
nivel por vencimiento. Una SVI sin restricciones no garantiza ausencia de arbitraje; SSVI admite
condiciones suficientes sencillas:

\begin{itemize}
\item \textbf{Calendario:} $\theta_T$ no decreciente en $T$ y
$0\le\partial_\theta\bigl(\theta\varphi(\theta)\bigr)\le\rho^{-2}\bigl(1+\sqrt{1-\rho^2}\bigr)\varphi(\theta)$.
\item \textbf{Mariposa:} $\theta\varphi(\theta)(1+|\rho|)<4$ y $\theta\varphi(\theta)^2(1+|\rho|)\le 4$.
\item \textbf{Función de potencia:} con $0<\gamma\le 1/2$ y $\eta(1+|\rho|)\le 2$ se cumplen ambas familias
de condiciones, siempre que $\theta_T$ sea no decreciente.
\end{itemize}

La primera condición de mariposa es coherente con la fórmula de momentos de \citet{lee2004}: la pendiente
asintótica de la varianza total no puede superar 2.

\begin{notatecnica}[Restricciones por construcción]
En la implementación de referencia las restricciones no se verifican a posteriori: se imponen con una
reparametrización sin restricciones, $\rho=0.999\tanh(a)$, $\gamma=\tfrac12\,\sigma(b)$,
$\eta=2\sigma(c)/(1+|\rho|)$ y $\theta_i=\sum_{j\le i}e^{d_j}$, con $\sigma$ la función logística. Cuando
el ajuste queda en la frontera de la región suficiente (como en la \figref{fig:ssvi_region}), conviene
registrarlo: la región es suficiente, no necesaria, y una restricción activa puede estar sesgando la
forma.
\end{notatecnica}

\grafica[0.6]{s2_ssvi_rebanadas}{Rebanadas de varianza total ajustadas para 7, 30, 60 y 91 días en el mundo
sintético. No se cruzan: esa es la forma gráfica de la ausencia de arbitraje de calendario.}{fig:ssvi_rebanadas}

\grafica[0.6]{s2_ssvi_region}{Región suficiente de la función de potencia en el plano $(|\rho|,\eta)$ con
$\gamma\le 1/2$. El cuadrado es el valor verdadero del mundo sintético; el círculo, el ajuste, que queda
en la frontera.}{fig:ssvi_region}

La prueba directa de mariposa usa la función $g$ de Durrleman, en la forma de \citet{gatheral2014}:
con $w'$ y $w''$ derivadas respecto de $k$,
\begin{equation}
g(k)=\Bigl(1-\frac{k\,w'}{2w}\Bigr)^2-\frac{(w')^2}{4}\Bigl(\frac1w+\frac14\Bigr)+\frac{w''}{2},
\label{eq:g}
\end{equation}
y la densidad de $x=\ln(S_T/F)$ es $p(k)=g(k)\,(2\pi w)^{-1/2}\exp\{-d_-(k)^2/2\}$ con
$d_-(k)=-k/\sqrt{w}-\sqrt{w}/2$. Una rebanada tiene densidad no negativa si y solo si $g\ge 0$
(\figref{fig:funcion_g}).

\grafica[0.6]{s2_funcion_g}{Función $g(k)$ para una rebanada SSVI restringida y para una SVI sin
restricciones con parámetros publicados como contraejemplo en \citet{gatheral2014}. Donde $g<0$ la
densidad implícita sería negativa: un arbitraje de mariposa.}{fig:funcion_g}

\paragraph{Comparación independiente} Como contraste se ajusta, sin forma paramétrica de volatilidad, una
densidad discreta $q_j\ge 0$ sobre una malla $x_j$:
\begin{equation}
C(K)=D\sum_j q_j\,(x_j-K)^+,\qquad \sum_j q_j=1,\qquad \sum_j q_j x_j=F.
\label{eq:convexo}
\end{equation}
Cualquier $q\ge0$ produce precios decrecientes y convexos en $K$, y las dos igualdades fijan la masa y la
media del forward. Con un precio auxiliar $p_i\in[b_i,a_i]$ por cotización, la pérdida por banda se vuelve
un problema de mínimos cuadrados lineales con cotas, convexo y rápido de resolver, al que se añade una
penalización de suavidad sobre $q$. Donde SSVI y este ajuste discrepan hay \emph{riesgo de modelo}, no un
dato \citep{aitsahalia1998}.

\subsection{Opciones americanas y equivalentes europeos}\label{sec:americanas}

SPY y las acciones individuales tienen opciones americanas; la identidad de extracción de la
\secref{sec:bl} es europea. Hay dos caminos: convertir cada cotización en un equivalente europeo con un
modelo que incluya ejercicio anticipado y dividendos, o excluir las observaciones donde la conversión sea
demasiado incierta. La conversión estándar (\emph{desamericanización}) tiene dos pasos: invertir el precio
americano con un árbol binomial \citep{cox1979} o diferencias finitas con dividendos discretos para obtener
una volatilidad, y valorar con ella la opción europea equivalente \citep{burkovska2018}.

Seleccionar opciones fuera del dinero reduce el problema pero no lo elimina
(figuras~\ref{fig:americanas_puts} y~\ref{fig:americanas_calls}). El error de conversión depende del
modelo y de los dividendos supuestos; por eso se contrasta dentro de los escenarios de incertidumbre
(\secref{sec:bandas}) en lugar de tratarlo como exacto.

\grafica[0.6]{s2_americanas_puts}{Sesgo, en puntos de volatilidad, de tratar puts americanas como europeas
(líneas continuas), comparado con medio spread expresado en volatilidad (discontinuas). A 91 días el sesgo
de puts cerca del dinero ya es comparable con medio spread, aunque estén fuera del dinero.}{fig:americanas_puts}

\grafica[0.6]{s2_americanas_calls}{Lo mismo para calls con un dividendo en efectivo dentro del plazo. Las
calls dentro del dinero antes de la fecha ex-dividendo tienen una prima de ejercicio anticipado grande; las
calls fuera del dinero apenas se ven afectadas.}{fig:americanas_calls}

\subsection{Extracción de la densidad}\label{sec:bl}

Para calls europeas y descuento determinista $D(t,T)$ se cumple \citep{breeden1978}
\begin{equation}
q^{\Q}_t(K,T)=\frac{1}{D(t,T)}\,\frac{\partial^2 C_t(K,T)}{\partial K^2},\qquad
\Q_t(S_T\le K)=1+\frac{1}{D(t,T)}\,\frac{\partial C_t(K,T)}{\partial K}.
\label{eq:bl}
\end{equation}
La identidad se aplica al \emph{precio ajustado}, nunca a diferencias de cotizaciones crudas: una segunda
derivada amplifica el ruido de redondeo y de spread hasta producir densidades negativas
(figuras~\ref{fig:densidad_cruda} y~\ref{fig:densidad_ajustada}).

\grafica[0.62]{s2_densidad_cruda}{Segundas diferencias de los mids crudos (mundo sintético con strikes cada
punto). La curva discontinua es la densidad verdadera, desconocida en la práctica; los puntos rojos son
valores negativos, imposibles para una densidad.}{fig:densidad_cruda}

\grafica[0.62]{s2_densidad_ajustada}{La misma identidad sobre precios ajustados. La densidad SSVI integra
0.9998 en la malla y su media es $1.0001F$; el ajuste monótono--convexo, independiente, coincide en el
cuerpo.}{fig:densidad_ajustada}

Los controles mínimos de cada densidad son tres, y se registran con su valor numérico:

\begin{enumerate}
\item \textbf{No negatividad} en toda la malla evaluada.
\item \textbf{Masa próxima a uno}. La masa que falta está en colas sin strikes y se reporta; no se
redistribuye en silencio.
\item \textbf{Primer momento compatible con el forward}: $\E^{\Q}[S_T]=F$.
\end{enumerate}

También se documentan el error numérico (paso de la malla, diferencias finitas) y la sensibilidad a la
extrapolación de strikes extremos.

\paragraph{Cuantiles sin masa inventada} Los cuantiles se obtienen invirtiendo una CDF \emph{anclada}: la
masa a la izquierda de la malla sale de la pendiente del precio (segunda identidad de~\eqref{eq:bl}), no
de integrar la densidad desde un extremo desconocido. No se normaliza la masa ni se recortan valores
negativos. Cada nivel $u$ se devuelve con su estado: «identificada» si cae entre $\Q(X\le x_{\min})$ y
$\Q(X\le x_{\max})$, con $[x_{\min},x_{\max}]$ el tramo respaldado por strikes fiables, y «no
identificada», sin valor, en otro caso. Una densidad negativa invalida todos los niveles. Normalizar
habría sido engañoso: a una normal a la que le falta el 2.3\% de la masa izquierda, la normalización le
desplaza la mediana 0.029 desviaciones típicas y le atribuye una cola que nadie observó.

\begin{noconcluir}
La media de $\Q$ está fijada por el forward. No es una expectativa de rentabilidad real y no contiene, por
construcción, ninguna opinión del mercado sobre el signo del rendimiento esperado.
\end{noconcluir}

\subsection{Conjunto de superficies plausibles y colas no identificadas}\label{sec:bandas}

Un solo ajuste oculta cuánto podría variar la distribución sin dejar de ser compatible con las
cotizaciones. Por eso se mantiene un \emph{conjunto} de superficies plausibles, coherentes con los spreads
y las restricciones: se sortean precios objetivo dentro de cada $[b_i,a_i]$, se reajustan ambos métodos y
se repite con supuestos alternativos de dividendos y desamericanización. De ese conjunto salen bandas para
cuantiles y probabilidades de cola (figuras~\ref{fig:bandas_densidad} y~\ref{fig:bandas_cuantiles}).

\grafica[0.62]{s2_bandas_densidad}{Bandas de densidades compatibles con los spreads, en escala logarítmica.
La banda SSVI es estrecha incluso donde no hay strikes fiables (más allá de los cuadrados): esa precisión
la pone la forma paramétrica, no los datos. El ajuste monótono--convexo muestra cuánto se ignora
realmente en las colas.}{fig:bandas_densidad}

\grafica[0.62]{s2_bandas_cuantiles}{Ancho de la banda de sensibilidad de cada cuantil. Los cuadrados marcan
el rango identificado $[u_{\min},u_{\max}]$, definido por los strikes fiables más extremos; fuera de él
los cuantiles se reportan como no identificados.}{fig:bandas_cuantiles}

\begin{noconcluir}
Estas bandas son de \emph{sensibilidad}: describen qué distribuciones son compatibles con los datos y los
supuestos. No son intervalos de confianza estadísticos, porque no hay un modelo de observación que lo
justifique. Si falta soporte en una cola, se marca como no identificada; no se rellena con precisión
artificial.
\end{noconcluir}

\begin{criterio}
Coherencia matemática (controles de la \secref{sec:bl} superados en cada sesión) y estabilidad de
cuantiles y colas frente a perturbaciones plausibles de spreads y supuestos.
\end{criterio}

% ============================================================================
\clearpage
\section{Transporte, dinámica y extracción de señal}\label{sec:transporte}
% ============================================================================

\ficha{Distribuciones $\Q_t$ a plazo constante con bandas y soporte identificado; información base de
precio, volatilidad y mercado.}
{Representar por cuantiles, medir distancias y cambios firmados, reducir dimensión, filtrar causalmente
y calcular la innovación.}
{Series de distancias, factores filtrados con incertidumbre, variables de cola e innovaciones.}
{Todo ajuste usa solo el pasado; reconstrucciones ordenadas; medidas recortadas declaradas como tales.}

\subsection{Representación por cuantiles relativos al forward}\label{sec:cuantiles}

Cada distribución se representa por su función cuantil de rendimientos logarítmicos relativos al
forward, a plazo constante $\tau$:
\begin{equation}
x_t(u,\tau)=F^{-1}_{\Q_t}(u),\qquad \text{con } x=\ln\!\bigl(S_{t+\tau}/F_{t,\tau}\bigr),\quad u\in(0,1).
\label{eq:cuantil}
\end{equation}
Medir respecto del forward retira el desplazamiento puramente mecánico del nivel de precios: si el
subyacente sube 1\% y nada más cambia, $x_t(u)$ no se mueve. La vista en precios absolutos se conserva
por separado para interpretar la posición. El plazo constante se obtiene interpolando $\theta_T$ entre
vencimientos listados; como $\theta$ interpolado linealmente sigue siendo no decreciente, se preserva la
ausencia de arbitraje de calendario. Por construcción $\E^{\Q}\bigl[e^{x}\bigr]=1$.

\subsection{Distancias entre distribuciones y cambios firmados}\label{sec:distancias}

En una dimensión, el transporte óptimo entre dos distribuciones es el \emph{reordenamiento monótono}: cada
cuantil se desplaza a su homólogo, $T=F_t^{-1}\circ F_{t-1}$ \citep{peyre2019}. La distancia de
Wasserstein de orden 2 se calcula entonces directamente sobre las funciones cuantil:
\begin{equation}
W_2^2(\Q_t,\Q_{t-1})=\int_0^1\bigl[x_t(u)-x_{t-1}(u)\bigr]^2\,du,\qquad
W_1(\Q_t,\Q_{t-1})=\int_0^1\bigl|x_t(u)-x_{t-1}(u)\bigr|\,du .
\label{eq:w2}
\end{equation}
$W_1$ es menos sensible que $W_2$ a desplazamientos extremos y equivale al área entre las dos funciones
de distribución. Para familias de posición y escala, $W_2^2=(\Delta\mu)^2+(\Delta\sigma)^2$, lo que ayuda
a leer su magnitud. Si solo existe cobertura fiable del 5\% al 95\%, se reporta la versión recortada
$W_{2,[a,b]}^2=\int_a^b[\Delta x(u)]^2du$ con su rango, nunca como distancia de toda la distribución.

\grafica[0.6]{s3_densidades_choque}{Distribución implícita a 30 días antes y después de un choque común en
el mundo sintético: más dispersión y una cola izquierda más larga.}{fig:densidades_choque}

\grafica[0.6]{s3_mapa_transporte}{Mapa de transporte óptimo entre las dos sesiones de la
\figref{fig:densidades_choque}: cada cuantil anterior (eje horizontal) se envía al mismo cuantil posterior
(eje vertical). La pendiente mayor que uno indica más dispersión; la diagonal sería ``sin cambio''.}{fig:mapa_transporte}

\grafica[0.6]{s3_funciones_cuantil}{Funciones cuantil antes y después del choque. $W_2$ es la distancia
$L^2$ entre ellas; la leyenda compara la versión completa, la recortada a $[0.05,0.95]$ y $W_1$.}{fig:funciones_cuantil}

La magnitud sola no dice qué cambió. Por eso se registran también los \emph{cambios firmados}
$\Delta x_t(u)$ en percentiles clave, que distinguen deformaciones con distancias parecidas
(figuras~\ref{fig:cambios_firmados} y~\ref{fig:distancias_arquetipos}).

\grafica[0.6]{s3_cambios_firmados}{Perfil firmado de tres deformaciones arquetípicas de una misma
distribución: más dispersión, cola izquierda más pesada y cambio solo en las alas. Un $\Delta x<0$ con
$u<1/2$ significa que la cola izquierda se alarga.}{fig:cambios_firmados}

\grafica[0.6]{s3_distancias_arquetipos}{Distancias de los tres arquetipos. $W_2$ pesa más el cambio en las
alas que $W_1$; la versión recortada casi no lo ve. Ninguna cifra aislada identifica el tipo de cambio.}{fig:distancias_arquetipos}

\grafica[0.64]{s3_mapa_cuantiles}{Mapa temporal de cuantiles: desviación de cada $x_t(u)$ respecto de su
media de entrenamiento, por sesión. El choque común de la sesión 170 alarga ambas colas; las columnas en
blanco son snapshots faltantes.}{fig:mapa_cuantiles}

\grafica[0.64]{s3_distancias_diarias}{Distancias entre sesiones consecutivas. Los puntos rojos son sesiones
con spreads amplios, donde parte de la distancia es ruido de ajuste y no cambio de distribución: una
razón para filtrar antes de interpretar.}{fig:distancias_diarias}

\subsection{Cadencia, velocidad y aceleración}\label{sec:cadencia}

Las velocidades se definen sobre una cadencia fija y con diferencias hacia atrás,
$v_t(u)=x_t(u)-x_{t-1}(u)$, en canales separados: \emph{cierre--apertura} y \emph{apertura--apertura}.
Fines de semana, feriados y ventanas intradía son intervalos de tipo distinto y se etiquetan como tales;
no se mezclan como si fueran idénticos ni se escalan por tiempo calendario sin un modelo. La aceleración
$a_t(u)=v_t(u)-v_{t-1}(u)$ queda como candidata secundaria por una razón aritmética: si cada $x_t$ tiene
ruido de medición independiente de varianza $\sigma^2$, la velocidad hereda $2\sigma^2$ y la aceleración
$6\sigma^2$ (figuras~\ref{fig:senal_ruido} y~\ref{fig:aceleracion}).

\grafica[0.6]{s3_senal_ruido}{Cociente entre la varianza de la señal y la del ruido de medición para el
nivel, la velocidad y la aceleración de un factor sintético. Diferenciar multiplica la desviación estándar
del ruido por $\sqrt2$ y $\sqrt6$ mientras la señal se vuelve más pequeña.}{fig:senal_ruido}

\grafica[0.6]{s3_aceleracion}{Aceleración latente frente al ruido de medición que contiene la aceleración
observada: fuera del choque son de magnitud comparable.}{fig:aceleracion}

\subsection{Reducción de dimensión: componentes funcionales y variables de cola}\label{sec:fpca}

Para resumir la forma se aplica un análisis de componentes principales funcional a los cuantiles
centrados, con pesos de cuadratura en la malla de $u$ y ajustado \emph{únicamente con datos de
entrenamiento}:
\begin{equation}
x_t(u)\approx\bar x(u)+\sum_{j=1}^{J}s_{tj}\,\phi_j(u),\qquad
s_{tj}=\int_0^1\bigl[x_t(u)-\bar x(u)\bigr]\phi_j(u)\,du .
\label{eq:fpca}
\end{equation}
Se conservan factores suficientes para explicar dispersión y asimetría. Pero una PCA global puede ignorar
las colas precisamente porque explican poca varianza total (figuras~\ref{fig:fpca_autofunciones}
y~\ref{fig:fpca_varianza}); por eso se añaden variables explícitas de cola, por ejemplo el ancho de la cola
izquierda $x_t(0.10)-x_t(0.01)$ dentro del rango identificado, o la asimetría por cuantiles
$[x_t(0.9)+x_t(0.1)-2x_t(0.5)]/[x_t(0.9)-x_t(0.1)]$.

\grafica[0.6]{s3_fpca_autofunciones}{Primeras autofunciones de la PCA funcional, estimadas solo con las
sesiones de entrenamiento del mundo sintético. La primera captura dispersión; la segunda, asimetría; la
tercera, las alas.}{fig:fpca_autofunciones}

\grafica[0.6]{s3_fpca_varianza}{Varianza explicada por componente, en escala logarítmica. La componente
de colas explica una fracción ínfima de la varianza total aunque pueda ser la más relevante para el
riesgo.}{fig:fpca_varianza}

Toda reconstrucción debe seguir siendo una función cuantil válida, es decir, creciente en $u$. Con pocos
factores eso no está garantizado; se corrige con el reordenamiento creciente \citep{chernozhukov2010}, que
no aumenta la distancia $L^p$ a la función verdadera, o se trabaja en una transformación que garantice
monotonía, como el logaritmo de la densidad cuantílica \citep{petersen2016}.

\subsection{Filtrado causal con incertidumbre de observación}\label{sec:kalman}

Los factores observados cada sesión mezclan cambio real y error de ajuste. Se filtran con un modelo de
estado lineal y el filtro de Kalman \citep{kalman1960}:
\begin{equation}
f_t=A f_{t-1}+\eta_t,\quad \eta_t\sim\mathcal N(0,Q);\qquad
y_t=f_t+\varepsilon_t,\quad \varepsilon_t\sim\mathcal N(0,R_t),
\label{eq:kalman}
\end{equation}
donde la varianza de observación $R_t$ \emph{depende de la calidad} de la sesión: ancho de spreads, número
de strikes válidos y error de ajuste. Con calidad baja, el filtro confía menos en la observación y más en
la dinámica. Si falta el snapshot, solo se ejecuta la predicción. Los parámetros $(A,Q,R_0)$ se estiman
por máxima verosimilitud con datos de entrenamiento.

Solo se usa el \emph{filtro}, que emplea información hasta $t$. Un suavizador que incorpore observaciones
futuras produce series más limpias y backtests inválidos. La comparación obligada es con una media
exponencial (EWMA) y con la persistencia: un filtro más complejo puede borrar señales rápidas
(figuras~\ref{fig:kalman} y~\ref{fig:kalman_error}).

\grafica[0.64]{s3_kalman}{Filtrado del primer factor en el mundo sintético. El filtro de Kalman, con $R_t$
mayor en las sesiones de spreads amplios, sigue al factor latente; la EWMA suaviza el ruido pero llega
tarde al choque de la sesión 170.}{fig:kalman}

\grafica[0.6]{s3_kalman_error}{Error frente al factor latente por ventana de evaluación. Ningún método
domina: la persistencia es casi tan buena con calidad normal, el filtro ayuda cuando los spreads son
amplios y la EWMA falla tras el choque.}{fig:kalman_error}

\subsection{La propuesta diferenciadora: la innovación}\label{sec:innovacion}

Buena parte del cambio diario de la distribución implícita es predecible con información que ya está en
el precio: si el mercado cae, la cola izquierda suele alargarse. Lo potencialmente informativo es lo que
\emph{no} se explica así. Se define la innovación
\begin{equation}
\epsilon_t(u)=\Delta x_t(u)-\widehat{\E}\bigl[\Delta x_t(u)\mid\mathcal{I}_t^{base}\bigr],
\label{eq:innovacion}
\end{equation}
donde $\mathcal I_t^{base}$ incluye el retorno ya observado del subyacente, su volatilidad, el movimiento
de mercado y sector, el gap de apertura y el calendario conocido de resultados y dividendos
(\figref{fig:innovacion_diagrama}). El modelo de expectativa se entrena \emph{solo con fechas anteriores};
para una acción, además, se calcula la innovación relativa al mercado,
$\epsilon^A_t-\beta\,\epsilon^M_t$, para separar señal idiosincrática de choque común.

\begin{figure}[H]
\centering
\resizebox{\linewidth}{!}{% Construcción de la innovación de la distribución implícita.
\begin{tikzpicture}[
  n/.style={caja=#1,text width=3.1cm,minimum height=13mm},
  node distance=6mm and 7mm]
\node[n=qazul] (obs) at (0,0) {\textbf{Cambio observado}\\ $\Delta x_t(u)=x_t(u)-x_{t-1}(u)$};
\node[n=qnaranja] (base) at (0,-2.3) {\textbf{Información base} $\mathcal{I}^{base}_t$\\ retorno ya observado,
  volatilidad, mercado y sector, gap, calendario};
\node[n=tenue] (mod) at (4.1,-2.3) {\textbf{Modelo de expectativa}\\ ajustado solo con fechas\\ anteriores a $t$};
\node[n=tenue] (esp) at (8.2,-2.3) {\textbf{Parte esperada}\\ $\widehat{\E}\bigl[\Delta x_t(u)\mid\mathcal{I}^{base}_t\bigr]$};
\node[circle,draw=tinta2,line width=0.7pt,minimum size=6mm,font=\sffamily\bfseries] (resta) at (8.2,0) {$-$};
\node[n=qazul] (inn) at (12.2,0) {\textbf{Innovación}\\ $\epsilon_t(u)$};
\node[n=qnaranja] (uso) at (12.2,-2.3) {\textbf{Variable candidata}\\ para pronosticar $\P$\\ después de $t$};
\draw[flecha] (obs) -- (resta);
\draw[flecha] (base) -- (mod);
\draw[flecha] (mod) -- (esp);
\draw[flecha] (esp) -- (resta);
\draw[flecha] (resta) -- (inn);
\draw[flecha] (inn) -- (uso);
\node[nota,anchor=north] at (6.1,-3.25) {La innovación no prueba información privada ni causalidad: es un candidato cuyo valor
  predictivo se mide fuera de muestra.};
\end{tikzpicture}
}
\caption{Construcción de la innovación. La resta se hace con una expectativa estimada exclusivamente con
el pasado.}
\label{fig:innovacion_diagrama}
\end{figure}

\grafica[0.6]{s3_innovacion_dispersion}{Cambio observado del cuantil 5\% de la acción A frente al esperado
dado la información base. La distancia vertical a la diagonal es la innovación. El evento idiosincrático
(círculo) alarga la cola sin movimiento de precio; el choque común (cuadrado) se explica solo en
parte.}{fig:innovacion_dispersion}

\grafica[0.64]{s3_innovacion_serie}{Innovación del mercado e innovación de la acción A relativa al mercado.
En la versión relativa el choque común casi se cancela y el evento propio de A sobresale.}{fig:innovacion_serie}

\begin{noconcluir}
La innovación no prueba información privada ni causalidad. Es una variable candidata para pronosticar
rendimientos \emph{posteriores} al instante de decisión, y su utilidad se mide fuera de muestra con el
protocolo de la \secref{sec:evidencia}. En el mundo sintético de estas figuras, por diseño, la innovación
no predice nada: solo se muestra cómo se calcula.
\end{noconcluir}

\subsection{Experimentos concretos}\label{sec:experimentos}

La \tabref{tab:experimentos} lista los primeros experimentos. En ninguno se asume una dirección rentable
antes de medirla.

\begin{table}[H]
\centering\small
\caption{Primeros experimentos sobre la señal.}
\label{tab:experimentos}
\begin{tabularx}{\linewidth}{@{}>{\RaggedRight}p{3.4cm} Y Y@{}}
\toprule
\encab{Experimento} & \encab{Variable} & \encab{Pregunta}\\
\midrule
Cola izquierda frente al mercado & $\epsilon^A_t(u)-\beta\,\epsilon^M_t(u)$ para $u\le 0.10$ & ¿Aporta la
deformación idiosincrática de la cola a los pronósticos de riesgo o rendimiento?\\
Deformación entre plazos & Diferencia de $\Delta x_t(u)/\sqrt{\tau}$ entre 7, 30 y 60 días & ¿Distingue
eventos de corto plazo de choques de nivel? Requiere cotizaciones suficientes en cada plazo.\\
Persistencia & Autocorrelación de $\epsilon_t$ y su vida media & ¿La innovación se revierte, persiste o se
incorpora al precio?\\
Idiosincrático frente a común & Descomposición de $\epsilon^A_t$ en componente de mercado y residual & ¿Qué
parte, si alguna, contiene información que no está en el mercado?\\
\bottomrule
\end{tabularx}
\end{table}

\grafica[0.6]{s3_plazos_comun}{Cambio de cuantiles por plazo, normalizado por $\sqrt{\tau}$, ante un
choque común de volatilidad: el perfil es parecido en los tres plazos, algo mayor en el corto.}{fig:plazos_comun}

\grafica[0.6]{s3_plazos_evento}{Ante un evento dentro de los próximos 7 días, la distribución a 7 días se
deforma mucho más que las de 30 y 60: la diferencia entre plazos es la huella del
evento.}{fig:plazos_evento}

\subsection{Dinámica de distribuciones como objetos}\label{sec:war}

Como modelo avanzado se compara la autorregresión de factores con una autorregresión en espacio de
Wasserstein \citep{zhang2022}. La idea: el espacio de distribuciones no es lineal, pero alrededor del
baricentro de Wasserstein $\bar Q$ tiene un espacio tangente que sí lo es. En una dimensión el baricentro
es el promedio de funciones cuantil, $\bar x(u)=\E[x_t(u)]$, y el ``logaritmo'' de $Q_t$ es simplemente
$v_t(u)=x_t(u)-\bar x(u)$. Se ajusta un autorregresivo sobre $v_t$ y se regresa a funciones cuantil,
proyectando cuando el pronóstico deja de ser creciente (\figref{fig:tangente}).

\begin{figure}[H]
\centering
\resizebox{0.86\linewidth}{!}{% Autorregresión en el espacio tangente del baricentro de Wasserstein.
\begin{tikzpicture}[x=1cm,y=1cm]
% Superficie curva: espacio de distribuciones con la métrica W2
\fill[qazul!6] (0,0) .. controls (2.5,1.6) and (6.5,1.6) .. (9,0.2)
  .. controls (9.6,-1.2) and (9.2,-2.6) .. (8.4,-3.4)
  .. controls (6,-2.1) and (2.8,-2.2) .. (0.4,-3.2)
  .. controls (-0.3,-2.2) and (-0.4,-1) .. (0,0);
\draw[qazul!60,line width=0.6pt] (0,0) .. controls (2.5,1.6) and (6.5,1.6) .. (9,0.2)
  .. controls (9.6,-1.2) and (9.2,-2.6) .. (8.4,-3.4)
  .. controls (6,-2.1) and (2.8,-2.2) .. (0.4,-3.2)
  .. controls (-0.3,-2.2) and (-0.4,-1) .. (0,0);
\node[nota,anchor=west] at (0.3,-2.75) {distribuciones con la métrica $W_2$ (no lineal)};
% Plano tangente en el baricentro
\coordinate (B) at (4.6,-0.9);
\fill[qnaranja!10,draw=qnaranja!70,line width=0.6pt] (1.9,0.35) -- (7.9,0.35) -- (7.2,-2.05) -- (1.2,-2.05) -- cycle;
\node[nota,anchor=east,text=tinta] at (7.05,-1.8) {espacio tangente $T_{\bar{Q}}$ (lineal)};
\fill[tinta] (B) circle (2.2pt);
\node[font=\sffamily\scriptsize,anchor=north] at (B) {$\bar{Q}$};
% Log de dos sesiones y predicción
\draw[flecha,draw=qazuloscuro] (B) -- (2.6,-0.3) node[anchor=east,font=\sffamily\scriptsize]{$v_{t-1}$};
\draw[flecha,draw=qazuloscuro] (B) -- (3.4,0.05) node[anchor=south,font=\sffamily\scriptsize]{$v_t$};
\draw[flecha,draw=qnaranja,line width=0.9pt] (B) -- (5.4,0.1) node[anchor=south west,font=\sffamily\scriptsize]{$\hat v_{t+1}=A\,v_t$};
% Distribuciones en la superficie
\fill[qazul] (1.7,1.05) circle (2pt) node[anchor=south,font=\sffamily\scriptsize]{$Q_{t-1}$};
\fill[qazul] (3.2,1.25) circle (2pt) node[anchor=south,font=\sffamily\scriptsize]{$Q_t$};
\fill[qnaranja] (6.3,1.2) circle (2pt) node[anchor=south,font=\sffamily\scriptsize]{$\widehat{Q}_{t+1}$};
\draw[flecha,draw=tenue,densely dotted] (1.7,1.0) -- (2.55,-0.2);
\draw[flecha,draw=tenue,densely dotted] (3.2,1.2) -- (3.4,0.15);
\draw[flecha,draw=qnaranja,densely dotted] (5.45,0.2) -- (6.25,1.1);
\node[nota,anchor=west,align=left] at (9.4,1.0) {$\exp$ y proyección a\\ funciones cuantil crecientes};
\node[nota,anchor=east,align=right] at (-0.45,0.6) {$\log$ en una dimensión:\\ $v_t(u)=x_t(u)-\bar{x}(u)$};
\end{tikzpicture}
}
\caption{Autorregresión en el espacio tangente del baricentro. En una dimensión el espacio tangente
coincide con el de cuantiles centrados, así que la PCA funcional de la \secref{sec:fpca} y este modelo son
parientes cercanos: difieren en cómo parametrizan la dinámica.}
\label{fig:tangente}
\end{figure}

Esta familia permite pronosticar distribuciones como objetos. Su uso aquí es una adaptación de
investigación: no demuestra capacidad para anticipar el precio de las acciones.

\subsection{Límites del transporte}\label{sec:limites}

El transporte óptimo entre marginales \emph{no} identifica un proceso temporal del precio ni probabilidades
de tocar un stop. Las figuras~\ref{fig:trayectorias_difusion} y~\ref{fig:trayectorias_transporte} lo
muestran con 20\,000 trayectorias: ambas tienen exactamente las mismas distribuciones a 30 y 60 días, pero
en una se reordenan como si cada cuantil siguiera su propio camino.

\grafica[0.6]{s3_trayectorias_difusion}{Trayectorias de una difusión (muestra de 70). La probabilidad de
tocar un stop en 90 antes de 60 días es la que indica el título.}{fig:trayectorias_difusion}

\grafica[0.6]{s3_trayectorias_transporte}{Las mismas marginales leídas como trayectorias del mapa monótono:
la probabilidad de tocar el stop cae casi a la mitad. Las marginales no deciden cuál de las dos es la
correcta.}{fig:trayectorias_transporte}

\begin{ideaclave}
El campo de transporte es una descripción geométrica de cómo cambió la distribución, no una ley física
que deba imponerse al mercado. Las probabilidades de trayectoria (stops, barreras, máximos) requieren un
modelo dinámico explícito y su propia validación.
\end{ideaclave}

\begin{criterio}
Separar ruido de cambios persistentes con filtros causales; todas las series de esta etapa calculables sin
información futura y reproducibles con el snapshot de cada sesión.
\end{criterio}

% ============================================================================
\clearpage
\section{De la información implícita Q a pronósticos reales P}\label{sec:pronostico}
% ============================================================================

\ficha{Variables de $\Q$ (factores, colas, innovaciones), información base y rendimientos realizados
después de cada snapshot.}
{Entrenar modelos supervisados por objetivo y horizonte, corregir cruces de cuantiles y combinar.}
{Distribuciones físicas $\P$ por horizonte: cuantiles, probabilidades de evento y volatilidad
esperada.}
{Etiquetas alineadas con el instante de decisión; mejora incremental medida dentro de una misma
familia de modelos.}

\subsection{Por qué Q no es P}\label{sec:qp}

$\Q$ y $\P$ son dos distribuciones distintas del mismo rendimiento. $\Q$ es la que hace coherentes los
precios de las opciones; $\P$ describe con qué frecuencia ocurren las cosas. Las une un núcleo de
valoración (factor estocástico de descuento proyectado) $m$:
\begin{equation}
q(x)=\frac{m(x)\,p(x)}{\E^{\P}[m]}.
\label{eq:nucleo}
\end{equation}
Sin supuestos adicionales $m$ no está identificado: la misma $q$ es compatible con muchas parejas $(m,p)$.
Los intentos de ``recuperar'' $\P$ a partir de $\Q$ requieren supuestos fuertes \citep{ross2015} cuya
validez está en discusión \citep{borovicka2016}. Para índices bursátiles se documenta, en promedio, una
prima de varianza: la dispersión implícita tiende a superar a la realizada \citep{carrwu2009}. La
\figref{fig:q_vs_p} ilustra un supuesto de ese tipo y la \figref{fig:nucleo} el núcleo que implicaría.

\grafica[0.62]{s4_q_vs_p}{Una distribución $\Q$ y una $\P$ supuesta para el mismo activo. La media de $\Q$
está fijada por el forward; la de $\P$ incorpora una prima. La cola izquierda de $\Q$ es más pesada.
Ilustrativo: la relación real entre ambas es lo que se estima, no lo que se supone.}{fig:q_vs_p}

\grafica[0.62]{s4_nucleo}{Cociente $q/p$ implícito en el supuesto de la \figref{fig:q_vs_p}. Otra elección
de $p$ daría otro cociente con la misma $q$: por eso $\P$ se estima con rendimientos realizados y no se
deduce de $\Q$.}{fig:nucleo}

\subsection{Diseño supervisado y alineación temporal}\label{sec:alineacion}

Los modelos se entrenan con rendimientos efectivamente realizados \emph{después} de cada snapshot. Las
variables de $\Q$ entran como regresores; no se etiquetan como probabilidades físicas por decreto. La
etiqueta de horizonte $h$ solo existe en $t+h$, y hasta entonces no puede usarse para entrenar, calibrar ni
actualizar intervalos (\figref{fig:alineacion}).

\begin{figure}[H]
\centering
\resizebox{\linewidth}{!}{% Alineación temporal del aprendizaje supervisado: variables, decisión y etiqueta.
\begin{tikzpicture}[x=0.62cm,y=1cm]
% Ventana de variables (pasado)
\fill[qazul!14] (-4,0.3) rectangle (0,1.9);
\node[nota,text=tinta,align=center] at (-2,1.1) {\textbf{variables}\\ información hasta\\ el corte};
% Instante de decisión
\draw[tinta,line width=0.9pt] (0,2.35) -- (0,-0.1);
\node[nota,anchor=south,text=tinta] at (0,2.35) {\textbf{decisión en $t$}};
% Etiquetas por horizonte
\foreach \y/\h in {1.55/1,1.0/5,0.45/20}{
  \fill[qnaranja!25] (0,\y-0.15) rectangle (\h,\y+0.15);
  \draw[qnaranja,line width=0.7pt] (\h,\y-0.15) -- (\h,\y+0.15);
  \node[nota,anchor=west,text=tinta] at (\h+0.2,\y) {$h=\h$};}
% Eje de sesiones
\draw[guia,line width=0.7pt] (-4,0) -- (22,0);
\foreach \s in {0,1,...,21} \draw[guia] (\s,0) -- (\s,-0.08);
\foreach \s/\e in {0/t,1/t+1,5/t+5,20/t+20} \node[nota,anchor=north] at (\s,-0.1) {$\e$};
\node[nota,anchor=north,text width=12cm,align=center] at (9,-0.5) {La etiqueta de horizonte $h$ es el rendimiento
  en $(t,\,t+h]$. Solo después de madurar, en $t+h$, puede entrar al entrenamiento, a la calibración o a la
  actualización conformal.};
\end{tikzpicture}
}
\caption{Alineación temporal. Las variables se calculan con información hasta el corte; la etiqueta es el
rendimiento posterior al instante de decisión y madura al final de su horizonte.}
\label{fig:alineacion}
\end{figure}

\subsection{Escalera de modelos}\label{sec:escalera}

El orden propuesto va de lo simple a lo flexible, y cada peldaño debe ganarle al anterior fuera de muestra
para justificar su complejidad (\figref{fig:escalera}):

\begin{enumerate}
\item \textbf{Referencia histórica con volatilidad dinámica}: cuantiles de residuos estandarizados
históricos escalados por una volatilidad EWMA o GARCH.
\item \textbf{Regresión cuantílica regularizada} \citep{koenker1978}: para cada nivel $\tau$,
\begin{equation}
\hat\beta_\tau=\arg\min_\beta\sum_t\rho_\tau\bigl(y_t-z_t^\top\beta\bigr)+\lambda\lVert\beta_{-0}\rVert_1,
\qquad \rho_\tau(e)=e\,\bigl(\tau-\mathbb 1\{e<0\}\bigr),
\label{eq:qr}
\end{equation}
un programa lineal (\figref{fig:pinball}).
\item \textbf{Boosting cuantílico} con la misma pérdida y corrección de cruces.
\item \textbf{NGBoost} \citep{duan2020}, que estima con gradiente natural los parámetros de una familia de
distribuciones condicional. La familia y sus colas se eligen por validación, sin imponer normalidad por
comodidad.
\end{enumerate}

\begin{figure}[H]
\centering
\resizebox{\linewidth}{!}{% Escalera de modelos P y combinación de pronósticos.
\begin{tikzpicture}[
  m/.style={caja=qnaranja,text width=3.05cm,minimum height=16mm},
  node distance=5mm and 5mm]
\node[m] (m1) {\textbf{Referencia histórica}\\ volatilidad dinámica y residuos estandarizados};
\node[m,right=of m1] (m2) {\textbf{Regresión cuantílica}\\ regularizada ($L_1$), un modelo por $\tau$};
\node[m,right=of m2] (m3) {\textbf{Boosting cuantílico}\\ con corrección de cruces};
\node[m,right=of m3] (m4) {\textbf{NGBoost}\\ distribución condicional; familia elegida en validación};
\draw[flecha] (m1) -- (m2); \draw[flecha] (m2) -- (m3); \draw[flecha] (m3) -- (m4);
\node[nota,above=1mm of m2.north east,anchor=south] {cada peldaño debe ganarle al anterior fuera de muestra};
\coordinate (medio) at ($(m2.south)!0.5!(m3.south)$);
\node[caja=qnaranja,text width=7.2cm,minimum height=12mm,below=11mm of medio] (comb)
  {\textbf{Combinación}: pesos $w_k\ge 0$, $\sum_k w_k=1$, aprendidos en validación temporal
  y comparados con pesos iguales};
\coordinate (bus) at ($(comb.north)+(0,0.45)$);
\foreach \m in {m1,m2,m3,m4} \draw[line width=0.6pt,draw=tinta2] (\m.south) |- (bus);
\draw[flecha] (bus) -- (comb.north);
\node[caja=tenue,text width=13.9cm,minimum height=9mm,below=6mm of comb] (obj)
  {\textbf{Un objetivo explícito por horizonte} (1, 5 y 20 sesiones): rendimiento $\cdot$ probabilidad de superar
  costos $\cdot$ cuantiles de pérdida $\cdot$ volatilidad futura};
\draw[flecha] (comb) -- (obj);
\end{tikzpicture}
}
\caption{Escalera de modelos $\P$, combinación y objetivos por horizonte.}
\label{fig:escalera}
\end{figure}

\grafica[0.56]{s4_pinball}{Pérdida cuantílica para $\tau=0.05$, $0.50$ y $0.95$. Su asimetría es la que
hace que el minimizador sea el cuantil $\tau$.}{fig:pinball}

Estimar cada cuantil por separado puede producir \emph{cruces}: en una zona de extrapolación el cuantil
del 25\% queda por debajo del 5\% (\figref{fig:cruces}). El reordenamiento creciente de los cuantiles
estimados en cada punto corrige el problema sin empeorar la estimación
\citep{chernozhukov2010} (\figref{fig:reordenado}).

\grafica[0.6]{s4_cruces}{Regresiones cuantílicas estimadas por separado con 70 observaciones sintéticas. Fuera del rango observado
de la variable, dos cuantiles se cruzan (punto rojo).}{fig:cruces}

\grafica[0.6]{s4_reordenado}{Los mismos cuantiles después de reordenarlos en cada valor de la variable: siempre crecientes en
$\tau$.}{fig:reordenado}

\subsection{Objetivos por horizonte}\label{sec:objetivos}

Cada objetivo tiene su pérdida y se entrena por separado. La información de opciones podría mejorar el
riesgo sin mejorar el signo del rendimiento; ambas conclusiones son útiles y se muestran por separado.

\begin{table}[H]
\centering\small
\caption{Objetivos, pronósticos y pérdidas. Horizontes candidatos: 1, 5 y 20 sesiones.}
\label{tab:objetivos}
\begin{tabularx}{\linewidth}{@{}l Y Y@{}}
\toprule
\encab{Objetivo} & \encab{Pronóstico} & \encab{Pérdida y diagnóstico}\\
\midrule
Rendimiento & Distribución completa; mediana & CRPS; pinball en $\tau=0.5$; calibración de cuantiles\\
Superar costos & Probabilidad de que el rendimiento neto sea positivo & Brier y log score; diagrama de fiabilidad\\
Cuantiles de pérdida & $q_{0.05}$ y $q_{0.10}$ & Pinball; cobertura y ancho de intervalos\\
Volatilidad futura & Varianza realizada esperada & QLIKE, robusta a proxies ruidosos \citep{patton2011}\\
\bottomrule
\end{tabularx}
\end{table}

\subsection{Combinación de pronósticos}\label{sec:combinacion}

Los pronósticos se combinan con pesos no negativos que suman uno, aprendidos en validación temporal
---por ejemplo, minimizando el CRPS de la mezcla o la pérdida cuantílica del promedio de cuantiles--- y se
comparan siempre con pesos iguales, que en la práctica son difíciles de superar \citep{timmermann2006}.
El acuerdo entre modelos parecidos no es evidencia independiente: si comparten variables y ventana de
entrenamiento, sus errores están correlacionados.

\subsection{Una referencia estructural: Martin y Wagner}\label{sec:martin}

\citet{martin2019} proponen una expresión del rendimiento esperado de una acción a partir de índices de
varianza implícita, derivada de \citet{martin2017}:
\begin{equation}
\frac{\E_t[R_{i,t+1}]-R_{f,t+1}}{R_{f,t+1}}=\mathrm{SVIX}_t^2+\tfrac12\Bigl(\mathrm{SVIX}_{i,t}^2-\overline{\mathrm{SVIX}}_t^2\Bigr),
\label{eq:martin}
\end{equation}
donde $\mathrm{SVIX}_t$ es el del mercado, $\mathrm{SVIX}_{i,t}$ el de la acción y
$\overline{\mathrm{SVIX}}_t^2$ el promedio ponderado por valor de la sección transversal. Queda como
referencia estructural si se dispone de esos insumos de mercado y de sección transversal; no sustituye la
calibración diaria descrita arriba.

\begin{noconcluir}
Ni la \eqref{eq:martin} ni ningún modelo de esta sección convierte a $\Q$ en $\P$. Todo pronóstico $\P$ se
valida con rendimientos realizados y con las reglas de puntuación de la \secref{sec:calibracion}.
\end{noconcluir}

\begin{criterio}
Mejora incremental fuera de muestra respecto de la referencia, dentro de la misma familia de modelos, o
una delimitación honesta de dónde no hay utilidad.
\end{criterio}

% ============================================================================
\clearpage
\section{Regímenes y confianza medible}\label{sec:confianza}
% ============================================================================

\ficha{Errores de pronóstico ya madurados, factores filtrados, indicadores de calidad de datos y de
ajuste.}
{Detectar cambios de comportamiento, medir calibración, construir intervalos adaptativos y decidir si
abstenerse.}
{Probabilidad de cambio reciente, métricas de calibración, intervalos con nivel efectivo y estado de
confianza por dimensión.}
{Calibración en un tramo temporal separado; actualización solo con resultados madurados; abstención con
motivo registrado.}

\subsection{Detección de cambios de régimen}\label{sec:bocpd}

La detección bayesiana en línea de puntos de cambio \citep{adams2007} mantiene la distribución posterior
de la \emph{longitud de corrida} $r_t$, el número de sesiones desde el último cambio. Con una probabilidad
de riesgo $H$ por sesión y la predictiva $\pi_t^{(r)}=p(x_t\mid r_{t-1}=r,\ x^{(r)})$ calculada con los
datos de la corrida,
\begin{align}
P(r_t=r+1,\,x_{1:t}) &= P(r_{t-1}=r,\,x_{1:t-1})\;\pi_t^{(r)}\;(1-H),\label{eq:bocpd1}\\
P(r_t=0,\,x_{1:t}) &= \textstyle\sum_{r}P(r_{t-1}=r,\,x_{1:t-1})\;\pi_t^{(r)}\;H .\label{eq:bocpd2}
\end{align}
Con una previa normal--gamma inversa, la predictiva es una $t$ de Student y la recursión es exacta. Se
aplica sobre errores de pronóstico o factores, y se compara primero con un detector simple por umbral
(figuras~\ref{fig:umbral} a~\ref{fig:cambio_reciente}).

\grafica[0.6]{s5_umbral}{Errores estandarizados de pronóstico en un mundo sintético con un cambio de
varianza en la sesión 160, y alarmas de un detector por umbral (dos de las últimas cinco sesiones con
$|z|>2.5$).}{fig:umbral}

\grafica[0.6]{s5_corrida}{Posterior de la longitud de corrida de BOCPD en escala logarítmica. Tras el
cambio, la masa colapsa a corridas cortas y vuelve a crecer desde cero.}{fig:corrida}

\grafica[0.6]{s5_cambio_reciente}{Probabilidad posterior de que haya ocurrido un cambio en las últimas diez
sesiones. Es alta justo después del cambio verdadero y baja el resto del tiempo.}{fig:cambio_reciente}

\begin{noconcluir}
Una probabilidad alta de cambio de régimen dice que el comportamiento reciente difiere del aprendido; no
es una probabilidad de caída del precio. Su uso natural es degradar la confianza y ampliar la
incertidumbre, no generar señales direccionales.
\end{noconcluir}

\subsection{Cuatro dimensiones de confianza}\label{sec:dimensiones}

La salida del sistema distingue cuatro fuentes de duda que no deben mezclarse en un solo número
(\figref{fig:confianza}). Un alto número de contratos observados no equivale a muchas sesiones
independientes: cien strikes de un mismo día son una sola observación del estado del mercado
(\secref{sec:particiones}).

\begin{figure}[H]
\centering
\resizebox{\linewidth}{!}{% Cuatro dimensiones de confianza, cada una con indicadores y estados.
\begin{tikzpicture}[
  tarjeta/.style={draw=#1,fill=#1!5,rounded corners=3pt,line width=0.6pt,text width=6.9cm,
    minimum height=2.55cm,align=left,inner sep=6pt,font=\sffamily\scriptsize,text=tinta,anchor=north west},
  estado/.style={circle,fill=#1,minimum size=2.2mm,inner sep=0pt}]
\node[tarjeta=qazul] (c1) at (0,0) {\textbf{\small 1. Calidad de datos}\\[2pt]
  ¿Se puede reconstruir $\Q$ en el corte?\\[2pt]
  \textcolor{tinta2}{Spread relativo, strikes válidos por vencimiento, retraso de la fuente, snapshots
  faltantes, sincronía con el subyacente.}};
\node[tarjeta=qazul] (c2) at (7.45,0) {\textbf{\small 2. Sensibilidad del ajuste}\\[2pt]
  ¿Cuánto cambia $\Q$ dentro de lo compatible con los datos?\\[2pt]
  \textcolor{tinta2}{Ancho de bandas de cuantiles, discrepancia SSVI frente a ajuste convexo, restricciones
  activas, error de desamericanización.}};
\node[tarjeta=qnaranja] (c3) at (0,-2.85) {\textbf{\small 3. Incertidumbre predictiva}\\[2pt]
  ¿Qué tan dispersa y qué tan calibrada es la predicción $\P$?\\[2pt]
  \textcolor{tinta2}{Ancho de intervalos, cobertura y PIT recientes, probabilidad de cambio de régimen,
  desempeño del último tramo maduro.}};
\node[tarjeta=qaqua] (c4) at (7.45,-2.85) {\textbf{\small 4. Estabilidad de la decisión}\\[2pt]
  ¿Cambiaría la recomendación ante perturbaciones plausibles?\\[2pt]
  \textcolor{tinta2}{Dispersión de pesos óptimos remuestreados, distancia a la región casi óptima,
  sensibilidad a costos y a escenarios de estrés.}};
% Leyenda de estados (icono + etiqueta; el color nunca va solo)
\node[estado=bueno] (e1) at (0.25,-5.95) {};
\node[nota,anchor=west,text=tinta] at (e1.east) {\ \textbf{apto}: se usa sin restricciones};
\node[estado=alerta] (e2) at (5.0,-5.95) {};
\node[nota,anchor=west,text=tinta] at (e2.east) {\ \textbf{degradado}: solo reducciones o mantener};
\node[estado=critico] (e3) at (10.15,-5.95) {};
\node[nota,anchor=west,text=tinta] at (e3.east) {\ \textbf{bloqueado}: sin decisiones basadas en señal};
\end{tikzpicture}
}
\caption{Dimensiones de confianza y sus indicadores. Cada dimensión se reporta con un estado explícito y
su etiqueta; el color nunca comunica el estado por sí solo.}
\label{fig:confianza}
\end{figure}

\subsection{Calibración y reglas de puntuación}\label{sec:calibracion}

La calibración de probabilidades se ajusta en un tramo temporal separado del de entrenamiento y del de
prueba. Las reglas de puntuación son propias \citep{gneiting2007}:
\begin{align}
\text{Brier} &= \tfrac1N\textstyle\sum_i (p_i-o_i)^2 \quad\text{para eventos binarios,}\label{eq:brier}\\
\mathrm{CRPS}(F,y) &= \int_{-\infty}^{\infty}\bigl(F(z)-\mathbb 1\{y\le z\}\bigr)^2dz
  = 2\int_0^1\rho_\tau\bigl(y-F^{-1}(\tau)\bigr)\,d\tau ,\label{eq:crps}
\end{align}
y la pérdida cuantílica $\rho_\tau$ para percentiles individuales. La segunda igualdad de \eqref{eq:crps}
conecta directamente los pronósticos cuantílicos con el CRPS \citep{gneiting2011}. Las puntuaciones se
acompañan de curvas de calibración, cobertura y ancho de intervalos (figuras~\ref{fig:fiabilidad}
y~\ref{fig:pit}).

\grafica[0.58]{s5_fiabilidad}{Diagrama de fiabilidad de un pronóstico de evento en el mundo sintético.
El modelo sobreconfiado empuja las probabilidades hacia los extremos: su curva es más plana que la
diagonal y su Brier, peor.}{fig:fiabilidad}

\grafica[0.58]{s5_pit}{Histogramas PIT de pronósticos de distribución. Un pronóstico calibrado da valores
$F(y)$ uniformes; intervalos demasiado estrechos producen forma de U y demasiado anchos, una
joroba.}{fig:pit}

\subsection{Conformal adaptativo}\label{sec:conformal}

Como envoltura de intervalos se prueba la inferencia conformal adaptativa \citep{gibbs2021}. Si
$\mathrm{err}_t=\mathbb 1\{y_t\notin\widehat C_t(\alpha_t)\}$, el nivel efectivo se actualiza como
\begin{equation}
\alpha_{t+1}=\alpha_t+\gamma\,(\alpha-\mathrm{err}_t),
\qquad
\Bigl|\tfrac1N{\textstyle\sum_{t}}\mathrm{err}_t-\alpha\Bigr|\le\frac{\max\{\alpha_1,1-\alpha_1\}+h\,\gamma}{\gamma\,N},
\label{eq:aci}
\end{equation}
donde $N$ es el número de intervalos evaluados y $h$ el horizonte. Con horizonte $h$, el error de la
etiqueta que vence en $t$ solo se conoce en $t$: la actualización usa exclusivamente errores madurados
(\figref{fig:alineacion}), y los $h-1$ intervalos ya emitidos con un nivel anterior explican el término
$h\gamma$; con $h=1$ es la cota de \citet{gibbs2021}. El umbral de cada intervalo es el
$\lceil(n+1)(1-\alpha_t)\rceil$-ésimo puntaje de calibración.

La cota vale para \emph{cualquier} secuencia, sin supuestos sobre la distribución, pero exige dos
convenciones: con $\alpha_t\ge1$ el conjunto es vacío (error seguro) y con $\alpha_t\le0$ es la recta
completa. Así $\alpha_t$ queda en $[-h\gamma(1-\alpha),\,1+h\gamma\alpha]$. Sin la primera convención la
cota falla: si los puntajes empatan en cero, por ejemplo con retornos nulos, un intervalo de ancho cero
siempre cubre, $\alpha_t$ crece sin límite y la tasa de error se queda en cero. La primera versión del
núcleo tenía ese defecto. Ahora las pruebas comprueban la cota en cada prefijo con secuencias límite
(empates, puntajes siempre crecientes, alternancias, un adversario que ve el intervalo y huecos en los
datos), además de datos aleatorios, que por sí solos no lo detectaban. Parte de la cobertura puede venir
de intervalos vacíos o infinitos, que no informan: la tasa agregada no basta para juzgar los intervalos.

\grafica[0.62]{s5_conformal_intervalos}{Intervalos del 90\% ante un cambio de volatilidad en la sesión 500
y otro en la 750 (datos sintéticos). El método de nivel fijo tarda en ajustar su ventana; el adaptativo
reacciona antes.}{fig:conformal_intervalos}

\grafica[0.62]{s5_conformal_cobertura}{Cobertura móvil de 50 sesiones. Ambos métodos terminan cerca del
90\% en el promedio de largo plazo, pero el de nivel fijo pasa muchas sesiones por debajo del 75\%: la
garantía agregada oculta fallas locales.}{fig:conformal_cobertura}

\grafica[0.58]{s5_conformal_alfa}{Trayectoria del nivel efectivo $\alpha_t$: baja cuando los errores se
acumulan y sube cuando sobra cobertura.}{fig:conformal_alfa}

\begin{noconcluir}
El control de cobertura a largo plazo no implica que mañana, o en cada régimen, exista exactamente la
cobertura nominal; tampoco convierte el intervalo en una probabilidad de rentabilidad.
\end{noconcluir}

\subsection{Abstención}\label{sec:abstencion}

Una entrada nueva requiere tres cosas a la vez: datos aptos, estabilidad suficiente del ajuste y evidencia
predictiva relevante (\figref{fig:abstencion}). Si falta cualquiera, se bloquean los incrementos basados
en señal y se comunica el problema. Las reducciones exigidas por límites de riesgo definidos de antemano se
evalúan aparte, con datos adecuados.

\begin{figure}[H]
\centering
\resizebox{0.92\linewidth}{!}{% Regla de abstención: los incrementos exigen tres condiciones; las reducciones por límites van aparte.
\begin{tikzpicture}[
  pregunta/.style={diamond,aspect=2.6,draw=qazul,fill=qazul!7,line width=0.6pt,align=center,
    font=\sffamily\scriptsize,text=tinta,inner sep=1.2pt,text width=2.7cm},
  bloque/.style={caja=critico,text width=3.6cm},
  permite/.style={caja=bueno,text width=3.6cm},
  node distance=6mm and 9mm]
\node[caja=tenue,text width=3.4cm] (ini) {Nueva sesión: snapshot, ajuste y pronósticos};
\node[pregunta,below=of ini] (d1) {¿Datos aptos?\\ cobertura, spreads, retraso};
\node[pregunta,below=of d1] (d2) {¿Ajuste estable?\\ bandas, contraste de métodos};
\node[pregunta,below=of d2] (d3) {¿Evidencia predictiva\\ relevante y calibrada?};
\node[permite,below=of d3] (ok) {Se permiten incrementos basados en señal, dentro de límites y fuera de la zona de no operación};
\node[bloque,right=of d1] (b1) {Bloquear incrementos basados en señal y comunicar el problema};
\node[caja=alerta,text width=3.6cm,right=of d3] (b3) {Sin nuevas entradas: mantener o reducir; registrar el motivo};
\draw[flecha] (ini) -- (d1);
\draw[flecha] (d1) -- node[nota,right]{sí} (d2);
\draw[flecha] (d2) -- node[nota,right]{sí} (d3);
\draw[flecha] (d3) -- node[nota,right]{sí} (ok);
\draw[flecha] (d1) -- node[nota,above]{no} (b1);
\draw[flecha] (d2.east) -| node[nota,pos=0.25,above]{no} (b1.south);
\draw[flecha] (d3) -- node[nota,above]{no} (b3);
% Rama independiente de límites de riesgo
\node[caja=qaqua,text width=3.5cm,left=14mm of d2] (lim) {Límites de riesgo definidos de antemano (concentración, CVaR, pérdidas)};
\node[caja=qaqua,text width=3.5cm,below=of lim] (red) {Evaluar reducciones con datos adecuados, aunque el modelo se abstenga};
\draw[flecha] (lim) -- (red);
\node[nota,text width=3.6cm,below=1mm of red] {abstenerse del modelo no significa ignorar una exposición existente};
\end{tikzpicture}
}
\caption{Regla de abstención. La rama de la izquierda no depende del modelo: abstenerse de la señal no
significa ignorar una exposición existente.}
\label{fig:abstencion}
\end{figure}

\begin{criterio}
Calibración aceptable en el tramo reservado, cobertura estable por subperiodos y regla de abstención
ejercitada en sesiones con datos degradados sin producir decisiones basadas en señal.
\end{criterio}

% ============================================================================
\clearpage
\section{Decisión de posición y riesgo conjunto}\label{sec:decision}
% ============================================================================

\ficha{Pronósticos $\P$ calibrados por activo y horizonte, posición actual, costos, límites y estado de
confianza.}
{Construir escenarios conjuntos, optimizar rendimiento, CVaR y costos, y aplicar una zona de no
operación.}
{Región de pesos estables o acción discreta (mantener, aumentar, reducir), con su razón y su riesgo
conjunto.}
{Unidades y horizontes compatibles; ninguna recomendación sin pasar por límites y abstención.}

\subsection{Escenarios conjuntos físicos}\label{sec:escenarios}

El riesgo de una cartera depende de cómo se mueven juntos sus activos; la suma de riesgos individuales no
lo reproduce. Los escenarios conjuntos se construyen con un modelo de factores de mercado y sector,
$R_i=\alpha_i+\beta_i^\top f+\varepsilon_i$, y una covarianza regularizada por contracción hacia una matriz
estructurada \citep{ledoit2004}. Como competidor se usa dependencia $t$ \citep{demarta2005}, que produce
más movimientos extremos conjuntos con la misma correlación lineal
(figuras~\ref{fig:copula_gauss} y~\ref{fig:copula_t}), y escenarios de estrés en los que la dependencia es
distinta de la de calma (\figref{fig:calma_crisis}).

\grafica[0.54]{s6_copula_gauss}{Cópula gaussiana con correlación 0.5 (datos sintéticos). En rojo, los
escenarios en que ambos activos quedan por debajo de su cuantil 2.5\%.}{fig:copula_gauss}

\grafica[0.54]{s6_copula_t}{Cópula $t$ con tres grados de libertad y la misma correlación: la frecuencia de
caídas conjuntas extremas casi se duplica.}{fig:copula_t}

\grafica[0.54]{s6_calma_crisis}{Dos regímenes con dependencia distinta. Un modelo que impone la misma
dependencia en calma y en crisis subestima el riesgo conjunto justo cuando importa.}{fig:calma_crisis}

\subsection{El problema costo--riesgo--rendimiento}\label{sec:problema}

Para acciones y ETF, con pesos $w$, retorno aleatorio conjunto $R$ al horizonte de decisión y posición
actual $w_{\text{actual}}$, el problema inicial es
\begin{equation}
\min_{w\in\mathcal C_t}\ -\E_{\widehat P_t}\bigl[w^\top R\bigr]
+\lambda\,\CVaR_{\alpha,\widehat P_t}\bigl(-w^\top R\bigr)+c\bigl(w-w_{\text{actual}}\bigr).
\label{eq:objetivo}
\end{equation}
Con $S$ escenarios $r_s$ y la representación de \citet{rockafellar2000},
$\CVaR_\alpha(L)=\min_\zeta\ \zeta+\frac{1}{(1-\alpha)S}\sum_s(L_s-\zeta)^+$, y con costos proporcionales
$c(\Delta w)=\kappa^\top|\Delta w|$, el problema es un programa lineal:
\begin{equation}
\begin{aligned}
\min_{w,\,b,\,v,\,\zeta,\,z}\quad & -\hat\mu^\top w+\lambda\Bigl(\zeta+\frac{1}{(1-\alpha)S}\sum_{s=1}^{S}z_s\Bigr)+\kappa^\top(b+v)\\
\text{s.a.}\quad & z_s\ge -r_s^\top w-\zeta,\quad z_s\ge 0,\qquad s=1,\dots,S,\\
& w=w_{\text{actual}}+b-v,\quad b,v\ge 0,\\
& \mathbf 1^\top w+\kappa^\top(b+v)\le 1,\quad 0\le w\le w_{\max},\quad w\in\mathcal C_t .
\end{aligned}
\label{eq:lp}
\end{equation}
$\CVaR_{0.95}$ es el promedio del 5\% peor de las pérdidas simuladas (\figref{fig:perdidas}). El conjunto
factible $\mathcal C_t$ recoge las restricciones de la \tabref{tab:restricciones}. Rendimiento, riesgo y
costos deben expresarse en las mismas unidades y el mismo horizonte; los costos incluyen compra y venta,
fracciones de acción, mínimos por operación y, cuando corresponda, conversión de moneda.

\begin{table}[H]
\centering\small
\caption{Restricciones del conjunto factible $\mathcal C_t$.}
\label{tab:restricciones}
\begin{tabularx}{\linewidth}{@{}l Y@{}}
\toprule
\encab{Restricción} & \encab{Forma lineal típica}\\
\midrule
Presupuesto y caja & $\mathbf 1^\top w+\text{costos}\le 1$; sin apalancamiento en la primera versión\\
Concentración & $w_i\le w_{\max}$ por activo\\
Exposición sectorial & $\sum_{i\in\text{sector}}w_i\le s_{\max}$\\
Beta de la cartera & $\beta^\top w\le\beta_{\max}$\\
Rotación & $\mathbf 1^\top(b+v)\le\text{rot}_{\max}$ por sesión\\
Liquidez & valor negociado $\le$ fracción del volumen diario promedio\\
Presupuesto de riesgo (opcional) & $\CVaR_\alpha(-w^\top R)\le b_{\text{riesgo}}$\\
\bottomrule
\end{tabularx}
\end{table}

\begin{notatecnica}[Homogeneidad del objetivo]
Sin restricciones activas, el objetivo \eqref{eq:objetivo} con costos nulos es positivamente homogéneo:
duplicar $w$ duplica el rendimiento esperado y el CVaR. La solución tiende entonces a los extremos
---toda la caja o la frontera de $\mathcal C_t$--- y es muy sensible a $\lambda$. En un experimento
sintético de la implementación de referencia (prueba \texttt{test\_objetivo\_homogeneo\_lleva\_a\_extremos}),
con escenarios a cinco sesiones, $\lambda=0.01$ lleva a los límites de concentración y $\lambda=0.02$ a
caja total. Conviene fijar la escala con un presupuesto de riesgo
explícito (última fila de la \tabref{tab:restricciones}), que además es más fácil de interpretar.
\end{notatecnica}

\subsection{CVaR: qué mide y qué no}\label{sec:cvar}

\grafica[0.6]{s6_perdidas}{Distribución simulada de pérdidas a cinco sesiones de una cartera sintética de
cuatro activos. El VaR 95\% es el umbral del 5\% peor; el CVaR 95\% es el promedio de esa cola.}{fig:perdidas}

\grafica[0.56]{s6_suma_riesgos}{Suma ponderada de los CVaR individuales frente al CVaR de la cartera. La
diferencia, que es el beneficio de diversificación, casi desaparece en el escenario de estrés.}{fig:suma_riesgos}

\begin{noconcluir}
El CVaR evalúa el promedio de pérdidas en la cola especificada \emph{bajo los escenarios supuestos}. No
es un límite garantizado de pérdida: si los escenarios subestiman la dependencia o las colas, el CVaR
subestima el riesgo.
\end{noconcluir}

\subsection{Zona de no operación y estabilidad de la decisión}\label{sec:no_operacion}

Un cambio pequeño de peso debe justificar sus costos \emph{y} superar una banda de incertidumbre. El
rendimiento esperado es la entrada más ruidosa del problema; al remuestrear escenarios y perturbar
$\hat\mu$, el óptimo recorre buena parte de la frontera del presupuesto de riesgo
(\figref{fig:region_estable}). Por eso la salida no es un punto sino una \emph{región casi óptima}
\begin{equation}
\mathcal R_\delta=\bigl\{w\in\mathcal C_t:\ \hat\mu^\top w\ \ge\ \hat\mu^\top w^\ast-\delta\bigr\},
\label{eq:region}
\end{equation}
con $\delta$ del orden de la dispersión bootstrap del valor del óptimo. La regla es simple: si
$w_{\text{actual}}\in\mathcal R_\delta$, mantener; si no, moverse solo hasta la región, con el traslado de
menor costo $\arg\min_{w\in\mathcal R_\delta}\kappa^\top|w-w_{\text{actual}}|$, no hasta el punto
$w^\ast$. La recomendación se expresa como región o como acción discreta (mantener, aumentar, reducir,
cerrar), evitando falsa precisión como ``7.431\%'' cuando el rango razonable es amplio
(\figref{fig:histograma_pesos}).

\grafica[0.62]{s6_region_estable}{Región casi óptima en el plano de pesos de dos activos (supuesto: el
rendimiento esperado tiene una desviación estándar del 30\% de su valor). Los puntos son óptimos
remuestreados; la línea discontinua, el borde interior de la región. La cartera actual dentro de la región
se mantiene; la que excede el presupuesto se reduce solo hasta la región.}{fig:region_estable}

\grafica[0.58]{s6_histograma_pesos}{Distribución del peso óptimo en A entre remuestreos. La cifra puntual
aparenta una precisión que el rango del 10\% al 90\% desmiente.}{fig:histograma_pesos}

\subsection{Fase posterior: optimización robusta}\label{sec:dro}

En una fase posterior la distribución puntual $\widehat P_t$ se reemplaza por un conjunto de
distribuciones plausibles y se optimiza el peor caso dentro de él:
\begin{equation}
\min_{w\in\mathcal C_t}\ \sup_{Q\in\mathcal B_\varepsilon(\widehat P_t)}\ \E_Q\bigl[\ell(w,R)\bigr],
\qquad \mathcal B_\varepsilon(\widehat P_t)=\{Q:\ W(Q,\widehat P_t)\le\varepsilon\}.
\label{eq:dro}
\end{equation}
Una bola de Wasserstein es una opción natural (\figref{fig:ambiguedad}); para pérdidas lineales por
tramos, como la del CVaR, el problema admite reformulaciones tratables \citep{esfahani2018}. El radio se
selecciona con validación temporal y pruebas de sensibilidad, no a partir de un porcentaje de confianza
inventado.

\begin{figure}[H]
\centering
\resizebox{0.86\linewidth}{!}{% Conjunto de ambigüedad de Wasserstein alrededor de la distribución estimada.
\begin{tikzpicture}[x=1cm,y=1cm]
\fill[qaqua!8] (0,0) circle (2.1);
\draw[qaqua,line width=0.8pt] (0,0) circle (2.1);
\fill[tinta] (0,0) circle (2.2pt);
\node[font=\sffamily\scriptsize,anchor=north] at (0,-0.08) {$\widehat{P}_t$};
\foreach \a/\r in {20/1.2,75/0.8,130/1.5,200/1.1,250/1.7,310/0.9,345/1.6}
  \fill[tenue] (\a:\r) circle (1.6pt);
\fill[qnaranja] (35:2.1) circle (2.6pt);
\node[font=\sffamily\scriptsize,anchor=south west] at (35:2.15) {$Q^\ast(w)$: peor caso para la cartera $w$};
\draw[flecha,draw=tinta2] (0,0) -- (-145:2.1);
\node[nota,anchor=north east] at (-150:1.2) {radio $\varepsilon$};
\node[font=\sffamily\scriptsize,star,star points=5,fill=qazul,inner sep=1.4pt] at (-40:1.35) {};
\node[nota,anchor=west] at (-40:1.5) {$P$ verdadera (desconocida)};
\node[nota,anchor=west,text width=5.2cm,align=left] at (3.2,0.4) {\textbf{Conjunto de ambigüedad}\\[1pt]
  $\mathcal{B}_\varepsilon(\widehat P_t)=\{Q: W(Q,\widehat P_t)\le\varepsilon\}$.\\[3pt]
  $\varepsilon$ se elige por validación temporal y pruebas de sensibilidad: pequeño, optimiza para la
  estimación; grande, se vuelve excesivamente conservador.};
\end{tikzpicture}
}
\caption{Conjunto de ambigüedad de Wasserstein. La decisión se protege contra el peor caso dentro de la
bola, que puede no contener a la distribución verdadera.}
\label{fig:ambiguedad}
\end{figure}

\begin{noconcluir}
Las garantías teóricas de esta familia dependen de supuestos (por ejemplo, muestras independientes y
colas controladas) que no se trasladan automáticamente a mercados dependientes y cambiantes.
\end{noconcluir}

\subsection{Control predictivo con horizonte móvil}\label{sec:mpc}

Después se prueba el control predictivo: planear varias decisiones, ejecutar solo la primera y recalcular
con datos nuevos \citep{boyd2017} (\figref{fig:mpc}):
\begin{equation}
\max_{w_t,\dots,w_{t+H-1}}\ \sum_{k=0}^{H-1}\Bigl(\hat\mu_{t+k}^\top w_{t+k}-\lambda\,\mathcal{R}_{t+k}(w_{t+k})
-c\bigl(w_{t+k}-w_{t+k-1}\bigr)\Bigr).
\label{eq:mpc}
\end{equation}
Requiere escenarios temporales coherentes entre horizontes, no simplemente marginales de 1, 5 y 20 días
ensambladas.

\begin{figure}[H]
\centering
\resizebox{0.8\linewidth}{!}{% Control predictivo con horizonte móvil: planear varias decisiones, ejecutar solo la primera.
\begin{tikzpicture}[x=1.45cm,y=1cm,
  plan/.style={circle,draw=qaqua,fill=white,line width=0.8pt,minimum size=2.6mm,inner sep=0pt},
  ejec/.style={circle,draw=qaqua,fill=qaqua,line width=0.8pt,minimum size=3.2mm,inner sep=0pt}]
\node[titulo,anchor=east] at (-0.35,1.5) {Plan en $t$};
\node[titulo,anchor=east] at (-0.35,0) {Plan en $t+1$};
\foreach \k/\y in {0/1.75,1/1.45,2/1.3,3/1.25,4/1.2} \node[plan] (a\k) at (\k,\y) {};
\foreach \k in {0,...,3}{\pgfmathtruncatemacro{\n}{\k+1}\draw[qaqua!60,line width=0.6pt] (a\k) -- (a\n);}
\node[ejec] at (0,1.75) {};
\foreach \k/\y in {1/0.2,2/-0.05,3/-0.15,4/-0.2,5/-0.22} \node[plan] (b\k) at (\k,\y) {};
\foreach \k in {1,...,4}{\pgfmathtruncatemacro{\n}{\k+1}\draw[qaqua!60,line width=0.6pt] (b\k) -- (b\n);}
\node[ejec] at (1,0.2) {};
\draw[guia,line width=0.7pt] (-0.2,-0.75) -- (5.4,-0.75);
\foreach \k/\e in {0/t,1/t+1,2/t+2,3/t+3,4/t+4,5/t+5}{\draw[guia] (\k,-0.75) -- (\k,-0.83);
  \node[nota,anchor=north] at (\k,-0.85) {$\e$};}
\draw[flecha,draw=qnaranja] (0.35,1.05) .. controls (0.6,0.8) .. (0.85,0.5);
\node[nota,anchor=west,text=tinta] at (0.72,0.95) {llegan datos nuevos: se recalcula todo el plan};
\node[nota,anchor=west] at (5.55,1.2) {solo se ejecuta\\ el primer paso\\ (círculo lleno)};
\end{tikzpicture}
}
\caption{Horizonte móvil. En cada sesión se optimiza una trayectoria de pesos y solo se ejecuta su primer
paso.}
\label{fig:mpc}
\end{figure}

\subsection{Si se gestionan opciones}\label{sec:opciones}

Si la posición incluye opciones, el P\&L lineal se sustituye por la \emph{revaluación completa} bajo
escenarios conjuntos de subyacente, superficie, paso del tiempo y ejercicio o asignación. La aproximación
delta--gamma--vega sirve como diagnóstico, no como valoración fiable de toda opción en movimientos grandes
(figuras~\ref{fig:opcion_spot} y~\ref{fig:opcion_vol}).

\grafica[0.56]{s6_opcion_spot}{P\&L a un día de un call a 30 días en el dinero: revaluación completa frente
a aproximaciones por griegas. Para movimientos grandes, delta--gamma sobrestima la ganancia y, a la
baja, inventa una pérdida que el call no puede tener.}{fig:opcion_spot}

\grafica[0.56]{s6_opcion_vol}{El mismo call cuando la volatilidad sube al caer el precio. Añadir vega no
arregla la aproximación en los extremos.}{fig:opcion_vol}

\begin{pendiente}
Confirmar si la posición del usuario incluye o incluirá opciones antes de construir esta rama.
\end{pendiente}

\begin{criterio}
Riesgo factible bajo todas las restricciones y resultados estables en pruebas de estrés: la recomendación
no cambia de signo ante perturbaciones plausibles de escenarios, costos y $\hat\mu$.
\end{criterio}

% ============================================================================
\clearpage
\section{Protocolo de evidencia}\label{sec:evidencia}
% ============================================================================

\ficha{Registro previo de objetivos, transformaciones y variantes; series de pronósticos y decisiones
fuera de muestra.}
{Particiones temporales con purga, comparación incremental, pruebas de pronóstico y de política,
corrección por selección y estrés.}
{Estimaciones de mejora con intervalos por bloques, veredicto por objetivo y horizonte, informe de
estrés.}
{Umbrales fijados antes del holdout; ``evidencia insuficiente'' cuando el intervalo no se separa de cero.}

\subsection{Registro previo}\label{sec:preregistro}

Antes de probar nada se registran los objetivos, las transformaciones de variables y el número de variantes
que se evaluarán. Cada experimento posterior queda en el registro, haya funcionado o no: el número total
de pruebas es un insumo de la corrección por selección (\secref{sec:seleccion}).

\subsection{Particiones temporales, purga y tamaño efectivo}\label{sec:particiones}

Se separan cuatro tramos con funciones distintas: entrenamiento, ajuste de hiperparámetros, calibración y
prueba final. La validación es \emph{walk-forward}, sin particiones aleatorias, y con purga de las
observaciones cuya etiqueta cruza una frontera entre tramos \citep{lopezdeprado2018}. Cuando el esquema
pueda reintroducir observaciones correlacionadas se añade un embargo (\figref{fig:walkforward}).

\begin{figure}[H]
\centering
\resizebox{\linewidth}{!}{% Walk-forward con purga y embargo; tramos de entrenamiento, ajuste, calibración y prueba.
\begin{tikzpicture}[x=1cm,y=1cm,
  tramo/.style={minimum height=5.2mm,inner sep=0pt,font=\sffamily\tiny,text=tinta,anchor=west},
  purga/.style={minimum height=5.2mm,inner sep=0pt,anchor=west,pattern=north east lines,pattern color=tenue,
    draw=tenue,line width=0.3pt}]
\def\bloque#1#2{% #1 = fila (y), #2 = longitud del entrenamiento
  \node[tramo,fill=qazul!40,minimum width=#2 cm] (e) at (0,#1) {entrenamiento};
  \node[purga,minimum width=3.5mm] (g1) at (e.east) {};
  \node[tramo,fill=qnaranja!40,minimum width=1.35cm] (h) at (g1.east) {\ ajuste hiperp.};
  \node[purga,minimum width=3.5mm] (g2) at (h.east) {};
  \node[tramo,fill=qaqua!40,minimum width=1.35cm] (c) at (g2.east) {\ calibración};
  \node[purga,minimum width=3.5mm] (g3) at (c.east) {};
  \node[tramo,fill=tinta2!45,minimum width=1.2cm,text=white] (p) at (g3.east) {\ prueba};}
\node[titulo,anchor=east] at (-0.15,0) {Ventana 1};
\bloque{0}{3.2}
\node[titulo,anchor=east] at (-0.15,-0.75) {Ventana 2};
\bloque{-0.75}{4.4}
\node[titulo,anchor=east] at (-0.15,-1.5) {Ventana 3};
\bloque{-1.5}{5.6}
% Prueba final reservada
\fill[pattern=crosshatch,pattern color=eje] (11.35,0.35) rectangle (14.3,-1.85);
\draw[tinta2,line width=0.6pt] (11.35,0.35) rectangle (14.3,-1.85);
\node[nota,text=tinta,align=center,fill=white,inner sep=2pt] at (12.82,-0.75) {\textbf{prueba final}\\ se abre una\\ sola vez};
% Eje y leyenda
\draw[flecha,draw=tinta2] (0,-2.2) -- (14.5,-2.2);
\node[nota,anchor=north east] at (14.5,-2.25) {tiempo};
\node[purga,minimum width=3.5mm] at (0,-2.75) {};
\node[nota,anchor=west] at (0.45,-2.75) {purga y embargo: se eliminan observaciones cuya etiqueta cruza la frontera, más un margen};
\end{tikzpicture}
}
\caption{Validación walk-forward. Cada ventana entrena con el pasado, ajusta, calibra y prueba en tramos
sucesivos separados por purga y embargo; la prueba final queda reservada hasta el final.}
\label{fig:walkforward}
\end{figure}

La dependencia reduce la información disponible más de lo que sugiere el número de filas. Con etiquetas
solapadas de $h$ sesiones, la autocorrelación inducida es $1-k/h$ y el tamaño efectivo es
aproximadamente $n/h$ (\figref{fig:obs_efectivas}); con $N$ activos de correlación media $\bar\rho$, el
número de activos independientes equivalentes es $N/[1+(N-1)\bar\rho]$ (\figref{fig:activos_efectivos}).
Por eso los intervalos se calculan con bloques temporales y teniendo en cuenta la dependencia entre activos.

\grafica[0.54]{s7_obs_efectivas}{Observaciones efectivas en 500 sesiones según el horizonte de la
etiqueta: un pronóstico a 20 sesiones tiene del orden de 25 observaciones independientes.}{fig:obs_efectivas}

\grafica[0.54]{s7_activos_efectivos}{Activos independientes equivalentes según la correlación media. Cinco
acciones con correlación 0.6 equivalen a menos de dos.}{fig:activos_efectivos}

\subsection{Comparación incremental}\label{sec:incremental}

La pregunta de investigación se responde con una escalera de conjuntos de información
(\figref{fig:incremental}): A, precio, volumen y volatilidad; B, A más la distribución $\Q$ actual; C, B
más cambios simples; D, C más transporte e innovaciones; E, D más aceleración y regímenes. Se compara
dentro de una misma familia de predicción para atribuir el cambio a la información adicional y no al
aumento simultáneo de capacidad del modelo.

\begin{figure}[H]
\centering
\resizebox{\linewidth}{!}{% Comparación incremental A-E: cada paso añade información, no capacidad del modelo.
\begin{tikzpicture}[x=1cm,y=1cm,
  paso/.style={draw=qazul,fill=qazul!#1,rounded corners=2pt,line width=0.6pt,anchor=south west,
    text width=2.55cm,align=center,font=\sffamily\scriptsize,text=tinta,inner sep=3pt}]
\node[paso=6,minimum height=1.0cm] (A) at (0,0) {\textbf{A}\\ precio, volumen\\ y volatilidad};
\node[paso=12,minimum height=1.55cm] (B) at (2.95,0) {\textbf{B = A +}\\ distribución $\Q$ actual};
\node[paso=18,minimum height=2.1cm] (C) at (5.9,0) {\textbf{C = B +}\\ cambios simples};
\node[paso=24,minimum height=2.65cm] (D) at (8.85,0) {\textbf{D = C +}\\ transporte e\\ innovaciones};
\node[paso=30,minimum height=3.2cm] (E) at (11.8,0) {\textbf{E = D +}\\ aceleración y\\ regímenes};
\foreach \a/\b in {A/B,B/C,C/D,D/E} \draw[flecha] (\a.east) ++(0,0.3) -- ++(0.36,0);
\draw[guia,line width=0.7pt] (-0.1,-0.05) -- (14.5,-0.05);
\node[nota,anchor=north] at (7.2,-0.15) {misma familia de predicción e hiperparámetros equivalentes en cada comparación:
  la mejora se atribuye a la información añadida, no a más capacidad};
\end{tikzpicture}
}
\caption{Comparación incremental A--E. La hipótesis principal corresponde al paso de C a D.}
\label{fig:incremental}
\end{figure}

\subsection{Dos pruebas independientes: pronóstico y política}\label{sec:dos_pruebas}

Una señal puede predecir mejor y no pagar sus costos. Por eso hay dos pruebas separadas:

\begin{itemize}
\item \textbf{Mejora de pronóstico}: diferencial de reglas de puntuación \eqref{eq:diferencial} por
objetivo y horizonte.
\item \textbf{Mejora de política de posición}: rendimiento neto, drawdown, pérdidas de cola, rotación y
exceso frente a un benchmark de riesgo comparable, contra las referencias de la \tabref{tab:referencias}.
\end{itemize}

\begin{table}[H]
\centering\small
\caption{Referencias para la prueba de política. Se alinean moneda, aportaciones, horarios y exposición.}
\label{tab:referencias}
\begin{tabularx}{\linewidth}{@{}l Y@{}}
\toprule
\encab{Referencia} & \encab{Qué controla}\\
\midrule
Mantener las posiciones & El valor de no hacer nada, incluidos costos evitados.\\
SPY con dividendos & La exposición pasiva al mercado.\\
Control de volatilidad sencillo & Si la mejora es solo gestión de riesgo genérica.\\
Pesos iguales & Para selección de acciones: si el modelo supera una asignación ingenua.\\
\bottomrule
\end{tabularx}
\end{table}

\subsection{Sesgo de selección}\label{sec:seleccion}

Probar muchas variantes y quedarse con la mejor infla el resultado aunque ninguna tenga ventaja real. Bajo
la hipótesis nula, el máximo esperado de $N$ ratios de Sharpe independientes con varianza $V$ es
aproximadamente \citep{bailey2014}
\begin{equation}
\widehat{SR}_0=\sqrt{V}\Bigl[(1-\gamma)\,\Phi^{-1}\!\Bigl(1-\tfrac1N\Bigr)+\gamma\,\Phi^{-1}\!\Bigl(1-\tfrac{1}{Ne}\Bigr)\Bigr],
\label{eq:sr0}
\end{equation}
con $\gamma$ la constante de Euler--Mascheroni. El Sharpe deflactado es la probabilidad de que el Sharpe
verdadero supere ese umbral, corrigiendo por asimetría $\hat\gamma_3$ y curtosis $\hat\gamma_4$ de $T$
rendimientos:
\begin{equation}
\mathrm{DSR}=\Phi\!\left(\frac{(\widehat{SR}-\widehat{SR}_0)\sqrt{T-1}}{\sqrt{1-\hat\gamma_3\,\widehat{SR}+\frac{\hat\gamma_4-1}{4}\,\widehat{SR}^2}}\right).
\label{eq:dsr}
\end{equation}
Se usa como diagnóstico adicional, junto con la evaluación fuera de muestra y el bootstrap por bloques; no
como certificado definitivo (figuras~\ref{fig:maximo_sharpe} y~\ref{fig:sharpe_deflactado}).

\grafica[0.56]{s7_maximo_sharpe}{Sharpe anual máximo esperado entre $N$ estrategias sin ninguna ventaja
real. Con dos años de datos diarios y cien variantes, el mejor resultado ronda 1.8 por puro azar; los
puntos confirman la aproximación por simulación.}{fig:maximo_sharpe}

\grafica[0.56]{s7_sharpe_deflactado}{Sharpe deflactado de un mismo resultado observado según el número
de variantes probadas (supuestos en el eje).}{fig:sharpe_deflactado}

\subsection{Intervalos por bloques y evidencia insuficiente}\label{sec:bootstrap}

La mejora media del diferencial de pérdida se estima con bootstrap de bloques móviles \citep{kunsch1989},
que conserva la dependencia dentro de cada bloque. Si el intervalo no permite distinguir la mejora de cero,
el veredicto es \emph{evidencia insuficiente}, no ``casi significativo''
(figuras~\ref{fig:bootstrap_distinguible} y~\ref{fig:bootstrap_insuficiente}).

\grafica[0.56]{s7_bootstrap_distinguible}{Distribución bootstrap de la mejora media del modelo D sobre el C
(serie sintética con autocorrelación). El intervalo del 90\% excluye el cero.}{fig:bootstrap_distinguible}

\grafica[0.56]{s7_bootstrap_insuficiente}{Una mejora media menor con la misma variabilidad: el intervalo
contiene el cero y el veredicto es evidencia insuficiente.}{fig:bootstrap_insuficiente}

\subsection{Pruebas de estrés}\label{sec:estres}

\begin{table}[H]
\centering\small
\caption{Pruebas de estrés. Ningún horario ni umbral se selecciona con el periodo final.}
\label{tab:estres}
\begin{tabularx}{\linewidth}{@{}l Y@{}}
\toprule
\encab{Perturbación} & \encab{Qué se observa}\\
\midrule
Choques de mercado y de correlación & CVaR, límites y cambio de recomendación.\\
Spreads más amplios & Ancho de bandas, estados de confianza y abstención.\\
Retrasos y cotizaciones incompletas & Que la decisión use solo datos disponibles en el corte.\\
Interrupciones y degradación del proveedor & Que un fallo no se convierta en orden.\\
Resultados trimestrales y cambios de dividendo & Conversión americana, forward y deformación por plazo.\\
Revisiones de datos & Reproducibilidad con los datos conocidos en cada momento.\\
Órdenes parcialmente ejecutadas & Reconciliación de posiciones y caja.\\
\bottomrule
\end{tabularx}
\end{table}

\subsection{Regla de avance}\label{sec:avance}

\begin{criterio}
Un modelo avanza si mejora la métrica definida de antemano, mantiene una calibración aceptable y produce
decisiones estables frente a perturbaciones plausibles. Los umbrales numéricos se fijan antes del holdout
con el presupuesto de riesgo acordado.
\end{criterio}

\begin{noconcluir}
Llegar a diciembre con el sistema operativo no equivale a haber demostrado ventaja estadística. Un sistema
que funciona sin fallas y una hipótesis confirmada son dos entregables distintos.
\end{noconcluir}

% ============================================================================
\clearpage
\section{Arquitectura y operación diaria}\label{sec:arquitectura}
% ============================================================================

\ficha{Snapshots, parámetros versionados, posiciones y caja, calendario de mercado.}
{Ejecutar la cadena de módulos en horarios fijos, registrar cada resultado y exponerlo por API e
interfaz.}
{Informe diario, recomendación o abstención, órdenes simuladas idempotentes y reconciliadas.}
{Corrida reproducible con el mismo snapshot; ningún fallo aguas arriba se convierte en orden.}

\subsection{Componentes}\label{sec:componentes}

La implementación propuesta separa un motor cuantitativo en Python, almacenamiento de snapshots en
Parquet consultable mediante DuckDB, metadatos de corridas y posiciones en SQLite, una API independiente y
una interfaz web para distribuciones y decisiones (\figref{fig:arquitectura}). Es una propuesta de
arquitectura, sin selección final de versiones ni benchmark de rendimiento.

\begin{figure}[H]
\centering
\resizebox{\linewidth}{!}{% Arquitectura: módulos separados, almacenamiento, registro, API e interfaz.
\begin{tikzpicture}[
  mod/.style={caja=#1,text width=2.05cm,minimum height=10mm},
  bd/.style={cylinder,shape border rotate=90,aspect=0.22,draw=tinta2,fill=grisclaro,line width=0.6pt,
    font=\sffamily\scriptsize,text=tinta,align=center,text width=2.3cm,minimum height=15mm,inner sep=2pt},
  x=1cm,y=1cm]
% Módulos del motor (serpiente)
\node[mod=qazul] (m1) at (3.9,0.9) {Ingestión};
\node[mod=qazul] (m2) at (6.45,0.9) {Control de\\ calidad};
\node[mod=qazul] (m3) at (9.0,0.9) {Superficies};
\node[mod=qazul] (m4) at (11.55,0.9) {Distribuciones};
\node[mod=qazul] (m5) at (14.1,0.9) {Transporte\\ y factores};
\node[mod=qnaranja] (m6) at (14.1,-0.85) {Predicción $\P$};
\node[mod=qnaranja] (m7) at (11.55,-0.85) {Calibración};
\node[mod=qaqua] (m8) at (9.0,-0.85) {Escenarios};
\node[mod=qaqua] (m9) at (6.45,-0.85) {Decisión};
\node[mod=qaqua] (m10) at (3.9,-0.85) {Simulación};
\foreach \a/\b in {m1/m2,m2/m3,m3/m4,m4/m5,m6/m7,m7/m8,m8/m9,m9/m10} \draw[flecha] (\a) -- (\b);
\draw[flecha] (m5) -- (m6);
\node[nota,anchor=south] at (9.0,1.5) {motor cuantitativo en Python: cada módulo con entradas y salidas versionadas};
% Almacenamiento
\node[bd] (s1) at (0.6,0.9) {Snapshots\\ inmutables\\ \tiny Parquet + DuckDB};
\node[bd] (s2) at (0.6,-1.6) {Corridas y\\ posiciones\\ \tiny SQLite};
\draw[flecha] (s1) -- (m1);
% Registro transversal
\node[caja=tenue,text width=12.1cm,minimum height=8mm] (reg) at (9.0,-2.3) {\textbf{Registro}: hashes de entrada,
  versión de modelo, parámetros, hora de corte y motivo de abstención de cada resultado};
\draw[flecha] (reg.west) -- (s2.east |- reg.west);
% Servicios
\node[caja=tenue,text width=2.6cm] (api) at (5.3,-3.75) {API independiente};
\node[caja=tenue,text width=2.6cm] (ui) at (9.0,-3.75) {Interfaz web};
\node[caja=qaqua,text width=2.6cm] (paper) at (12.7,-3.75) {Conector de paper};
\draw[flecha] (reg.south -| api.north) -- (api.north);
\draw[flecha] (api) -- (ui);
\draw[flecha] (reg.south -| paper.north) -- (paper.north);
\end{tikzpicture}
}
\caption{Módulos del motor y servicios. Los colores siguen las tres capas: medición $\Q$ (azul),
predicción $\P$ (naranja) y acción (verde).}
\label{fig:arquitectura}
\end{figure}

\begin{table}[H]
\centering\small
\caption{Responsabilidad de cada módulo. El paquete de referencia \texttt{quantileflow/} del repositorio
contiene versiones mínimas de los bloques matemáticos, usadas para las figuras.}
\label{tab:modulos}
\begin{tabularx}{\linewidth}{@{}l Y l@{}}
\toprule
\encab{Módulo} & \encab{Responsabilidad} & \encab{Referencia}\\
\midrule
Ingestión & Capturar snapshots inmutables con ambos sellos de tiempo & ---\\
Control de calidad & Filtros, estados de confianza y motivos de exclusión & ---\\
Superficies & SSVI restringido, ajuste convexo, desamericanización & \texttt{superficies}, \texttt{opciones}\\
Distribuciones & Breeden--Litzenberger, controles, cuantiles y soporte & \texttt{distribuciones}\\
Transporte y factores & $W_1$, $W_2$, cambios firmados, PCA funcional, filtros & \texttt{transporte}, \texttt{filtrado}\\
Predicción $\P$ & Modelos por objetivo y horizonte & \texttt{pronostico}\\
Calibración & Reglas de puntuación, PIT, conformal adaptativo & \texttt{calibracion}\\
Escenarios y decisión & Escenarios conjuntos, CVaR, costos, región casi óptima & \texttt{decision}\\
Regímenes y evidencia & BOCPD, Sharpe deflactado, bootstrap, particiones & \texttt{regimenes}, \texttt{evidencia}\\
Simulación y registro & Paper, reconciliación y trazabilidad & ---\\
\bottomrule
\end{tabularx}
\end{table}

\subsection{Trazabilidad y reproducibilidad}\label{sec:trazabilidad}

Cada resultado conserva hashes de entrada, versión del modelo, parámetros, hora de corte y motivo de
abstención. Con el mismo snapshot, la corrida se reproduce exactamente. El siguiente registro es un ejemplo
ilustrativo de su forma, no una salida real:

\begin{lstlisting}[style=registro]
{
  "corrida": "2026-09-25T09:45:00-04:00/SPY",
  "hora_corte": "2026-09-25T09:45:00-04:00",
  "snapshot": {"sha256": "3f9a...c21e", "fuente": "proveedor-a", "retraso_s": 0},
  "codigo": {"version": "0.1.0", "commit": "a1b2c3d"},
  "modelos": {"superficie": "ssvi-potencia", "p_h5": "qr-l1", "intervalos": "aci-gamma-0.01"},
  "parametros_sha256": "9d04...7b10",
  "confianza": {"datos": "apto", "ajuste": "degradado", "prediccion": "apto", "decision": "apto"},
  "salida": {"accion": "mantener", "region_peso": [0.30, 0.46], "cvar95_5s": 0.031},
  "abstencion": {"incrementos": "bloqueados", "motivo": "ajuste degradado: bandas de cola anchas"}
}
\end{lstlisting}

\subsection{Calendario diario}\label{sec:calendario}

La captura empieza en la apertura regular; la primera evaluación se hace alrededor de las 09:45 ET si la
calidad alcanza el umbral, con una comprobación posterior a las 10:00 y un registro próximo al cierre
(\figref{fig:calendario}). Las medias sesiones y los cambios de horario ajustan la agenda. Los canales
cierre--apertura y apertura--apertura se mantienen separados.

\begin{figure}[H]
\centering
\resizebox{\linewidth}{!}{% Calendario diario (hora del Este) y canales de cambio.
\begin{tikzpicture}[x=1cm,y=1cm,
  hito/.style={circle,fill=#1,minimum size=2.6mm,inner sep=0pt},
  nota/.style={font=\sffamily\scriptsize,text=tinta2,align=center}]
% Agenda vertical
\draw[guia,line width=0.8pt] (0,0.3) -- (0,-4.6);
\node[hito=qazul] at (0,0) {};
\node[titulo,anchor=east] at (-0.3,0) {09:30};
\node[nota,anchor=west,text=tinta,align=left] at (0.35,0) {\textbf{Apertura regular}: empieza la captura de snapshots};
\node[hito=qazul] at (0,-1.1) {};
\node[titulo,anchor=east] at (-0.3,-1.1) {09:45};
\node[nota,anchor=west,text=tinta,align=left] at (0.35,-1.1) {\textbf{Primera evaluación}, solo si la calidad\\ alcanza el umbral};
\node[hito=qazul] at (0,-2.2) {};
\node[titulo,anchor=east] at (-0.3,-2.2) {10:00};
\node[nota,anchor=west,text=tinta,align=left] at (0.35,-2.2) {\textbf{Comprobación posterior}: ¿la recomendación\\ es estable respecto de la de 09:45?};
\node[hito=qaqua] at (0,-3.3) {};
\node[titulo,anchor=east] at (-0.3,-3.3) {$\approx$15:50};
\node[nota,anchor=west,text=tinta,align=left] at (0.35,-3.3) {\textbf{Registro} próximo al cierre};
\node[hito=tenue] at (0,-4.35) {};
\node[titulo,anchor=east] at (-0.3,-4.35) {16:00};
\node[nota,anchor=west,text=tinta,align=left] at (0.35,-4.35) {Cierre; medias sesiones y cambios de horario\\ ajustan esta agenda};
% Canales de cambio
\begin{scope}[xshift=8.3cm,yshift=-0.4cm]
\node[titulo,anchor=west] at (0,1.0) {Dos canales que no se mezclan};
\fill[grisclaro] (0,-0.9) rectangle (2.2,-0.5);
\fill[grisclaro] (3.2,-0.9) rectangle (5.4,-0.5);
\node[nota,anchor=north] at (1.1,-0.95) {sesión $t-1$};
\node[nota,anchor=north] at (4.3,-0.95) {sesión $t$};
\draw[flecha,draw=qnaranja,line width=0.8pt] (2.2,-0.45) .. controls (2.5,0.15) and (2.9,0.15) .. (3.2,-0.45);
\node[nota,anchor=south] at (2.7,0.2) {cierre $\to$ apertura};
\draw[flecha,draw=qazul,line width=0.8pt] (0.25,-1.35) .. controls (1.4,-2.0) and (2.3,-2.0) .. (3.45,-1.35);
\node[nota,anchor=north] at (1.85,-1.95) {apertura $\to$ apertura};
\node[nota,anchor=north west,text width=6.2cm,align=left] at (0,-2.45) {Con una fuente retrasada 15 minutos, el
  instante \emph{ejecutable} de decisión se mueve en la misma cantidad: la evaluación de 09:45 solo puede
  usar datos hasta 09:30.};
\end{scope}
\end{tikzpicture}
}
\caption{Agenda de una sesión y canales de cambio. Con una fuente retrasada, el instante ejecutable de
decisión se mueve.}
\label{fig:calendario}
\end{figure}

\subsection{Pantalla}\label{sec:pantalla}

La pantalla muestra distribuciones con bandas, el mapa temporal de cuantiles, los cambios de cola y sus
innovaciones, los pronósticos $\P$ diferenciados de $\Q$, la posición actual frente a la propuesta con costos
y riesgo conjunto, la razón de la decisión y el historial de calibración (\figref{fig:pantalla}).

\begin{figure}[H]
\centering
\resizebox{0.95\linewidth}{!}{% Maqueta de la pantalla principal (solo estructura; sin datos).
\begin{tikzpicture}[x=1cm,y=1cm,
  panel/.style={draw=eje,fill=white,rounded corners=2pt,line width=0.6pt,anchor=north west,
    font=\sffamily\scriptsize,text=tinta,align=left,inner sep=5pt},
  etiq/.style={font=\sffamily\scriptsize\bfseries,text=tinta,anchor=north west}]
\draw[eje,fill=grisfondo,rounded corners=3pt,line width=0.7pt] (0,0) rectangle (15.4,-9.2);
% Barra superior
\fill[qazul!10] (0.05,-0.05) rectangle (15.35,-0.75);
\node[etiq] at (0.25,-0.18) {SPY \textperiodcentered{} 25 sep 2026 \textperiodcentered{} corte 09:45 ET};
\node[font=\sffamily\scriptsize,text=tinta,anchor=north east] at (15.2,-0.18)
  {\tikz[baseline=-0.6ex]\node[circle,fill=bueno,minimum size=2mm,inner sep=0pt]{}; datos: apto
   \quad \tikz[baseline=-0.6ex]\node[circle,fill=alerta,minimum size=2mm,inner sep=0pt]{}; ajuste: degradado};
% Fila 1
\node[panel,minimum width=7.45cm,minimum height=3.05cm,text width=7.1cm] at (0.25,-0.95)
  {\textbf{Distribución $\Q$ con bandas}\\ densidad y cuantiles a 30 días; colas no identificadas marcadas};
\node[panel,minimum width=7.45cm,minimum height=3.05cm,text width=7.1cm] at (7.9,-0.95)
  {\textbf{Mapa temporal de cuantiles}\\ últimas 60 sesiones; choques y datos faltantes señalados};
% Fila 2
\node[panel,minimum width=7.45cm,minimum height=2.2cm,text width=7.1cm] at (0.25,-4.2)
  {\textbf{Cambios de cola e innovaciones}\\ $\Delta x_t(u)$ firmado; innovación propia y relativa al mercado};
\node[panel,minimum width=7.45cm,minimum height=2.2cm,text width=7.1cm] at (7.9,-4.2)
  {\textbf{Pronósticos $\P$ (distintos de $\Q$)}\\ cuantiles a 1, 5 y 20 sesiones; calibración reciente};
% Fila 3
\node[panel,minimum width=4.85cm,minimum height=2.45cm,text width=4.5cm] at (0.25,-6.6)
  {\textbf{Posición actual frente a propuesta}\\ región de pesos estables; costos; riesgo conjunto};
\node[panel,minimum width=4.85cm,minimum height=2.45cm,text width=4.5cm] at (5.3,-6.6)
  {\textbf{Razón de la decisión}\\ texto generado a partir de cifras estructuradas; motivo de abstención};
\node[panel,minimum width=4.85cm,minimum height=2.45cm,text width=4.5cm] at (10.35,-6.6)
  {\textbf{Historial de calibración}\\ cobertura, PIT y aciertos por horizonte ya maduro};
\end{tikzpicture}
}
\caption{Maqueta de la pantalla principal (estructura, sin datos). Los estados de confianza aparecen en la
barra superior con icono y etiqueta.}
\label{fig:pantalla}
\end{figure}

\begin{noconcluir}
Un componente de lenguaje puede explicar resultados estructurados, pero no inventar cifras ni decidir
tamaños por texto libre. Todo número que aparezca en la explicación debe provenir del registro de la
corrida.
\end{noconcluir}

\subsection{Conector de paper y órdenes}\label{sec:ordenes}

El conector de paper es un componente separado: órdenes idempotentes, reconciliación de posiciones y caja,
cancelación y estado de ejecuciones (\figref{fig:ordenes}). Inicialmente el producto solo recomienda y
registra; la simulación de ejecución se incorpora después de comprobar el motor.

\begin{figure}[H]
\centering
\resizebox{\linewidth}{!}{% Ciclo de una orden en el conector de paper: validación, idempotencia y reconciliación.
\begin{tikzpicture}[
  est/.style={caja=#1,text width=2.35cm,minimum height=10mm},
  node distance=7mm and 8mm]
\node[est=qaqua] (rec) {Recomendación};
\node[diamond,aspect=2.4,draw=qazul,fill=qazul!7,line width=0.6pt,font=\sffamily\scriptsize,align=center,
  text=tinta,inner sep=1pt,text width=2.4cm,right=of rec] (puerta) {Puerta de\\ validación};
\node[est=tenue,right=of puerta] (prop) {Orden propuesta\\ \tiny clave idempotente};
\node[est=tenue,right=of prop] (env) {Enviada};
\node[est=bueno,below=of env] (ejec) {Ejecutada};
\node[est=tenue,left=of ejec] (parc) {Parcialmente\\ ejecutada};
\node[est=alerta,right=of ejec] (canc) {Cancelada o\\ rechazada};
\node[est=qaqua,below=of parc] (recon) {Reconciliación\\ posiciones y caja};
\node[est=tenue,left=of recon] (regi) {Registro};
\node[est=critico,below=of puerta] (sin) {Sin orden:\\ abstención registrada};
\draw[flecha] (rec) -- (puerta);
\draw[flecha] (puerta) -- node[nota,above]{pasa} (prop);
\draw[flecha] (puerta) -- node[nota,right,align=left]{datos o superficie\\ inválidos, límites} (sin);
\draw[flecha] (prop) -- (env);
\draw[flecha] (env) -- (parc);
\draw[flecha] (env) -- (ejec);
\draw[flecha] (env) -- (canc);
\draw[flecha] (parc) -- (ejec);
\draw[flecha] (parc) -- (recon);
\draw[flecha] (ejec) |- (recon);
\draw[flecha] (canc) |- (recon);
\draw[flecha] (recon) -- (regi);
\draw[flecha] (sin) |- (regi);
\draw[flecha,draw=tenue,densely dashed] (env.north) .. controls +(0,0.9) and +(0,0.9) .. node[nota,above]{reintento con la misma clave: no duplica} (prop.north);
\end{tikzpicture}
}
\caption{Ciclo de una orden simulada. La puerta de validación impide que un fallo de superficie o de datos
se convierta silenciosamente en una orden; la clave idempotente evita duplicados en los reintentos.}
\label{fig:ordenes}
\end{figure}

\begin{criterio}
Operación reproducible en paper durante un periodo acordado, con fallos inyectados (datos faltantes,
proveedor caído, ejecuciones parciales) manejados sin órdenes duplicadas ni posiciones descuadradas.
\end{criterio}

% ============================================================================
\clearpage
\section{Secuencia de construcción}\label{sec:secuencia}
% ============================================================================

La construcción avanza por etapas con un criterio verificable de salida (\tabref{tab:secuencia}). Ninguna
etapa se considera terminada porque su código exista: se considera terminada cuando supera su criterio.

\begin{table}[H]
\centering\small
\caption{Etapas, entregables y criterios para avanzar.}
\label{tab:secuencia}
\begin{tabularx}{\linewidth}{@{}>{\RaggedRight}p{2.6cm} Y Y@{}}
\toprule
\encab{Etapa} & \encab{Entregable} & \encab{Criterio para avanzar}\\
\midrule
Datos y especificación & Muestra auditada, eventos, horarios y contrato de datos & Cadenas reconstruibles en
el instante de decisión; cobertura y costo conocidos\\
Distribuciones & Superficies restringidas, cuantiles, bandas y \emph{replay} & Coherencia matemática y
estabilidad frente a spreads y colas\\
Dinámica & Transporte, factores filtrados e innovaciones & Separar ruido de cambios persistentes; cálculo
sin información futura\\
Pronóstico & Referencias y modelos $\P$ con calibración temporal & Mejora incremental fuera de muestra o
delimitación honesta de su utilidad\\
Decisión & Escenarios, CVaR, costos y zona de no operación & Riesgo factible y resultados estables en
pruebas de estrés\\
Paper & Informe diario y reconciliación & Operación reproducible y fallos manejados sin órdenes duplicadas\\
\bottomrule
\end{tabularx}
\end{table}

La estimación inicial de ingeniería es de 6 a 10 semanas para un prototipo integrado con un universo
pequeño, si los datos adecuados están disponibles (\figref{fig:calendario_obra}). No incluye esperar
suficientes observaciones prospectivas para evaluar habilidad predictiva, y se revisa después de auditar la
muestra de datos.

\grafica[0.62]{s9_calendario}{Secuencia tentativa. El reparto por etapa es ilustrativo y se superpone
parcialmente; lo que la propuesta fija es el total de 6 a 10 semanas de ingeniería, que se revisará tras la
auditoría de datos.}{fig:calendario_obra}

\begin{noconcluir}
La investigación puede concluir que el transporte ayuda al riesgo y no al rendimiento, o que no añade valor
sobre las referencias. Ambos son resultados válidos del programa, no fracasos del calendario.
\end{noconcluir}

% ============================================================================
\clearpage
\section{Qué priorizar y qué posponer}\label{sec:prioridades}
% ============================================================================

La prioridad inmediata es todo lo que hace creíble una conclusión, sea positiva o negativa: datos
correctos, SSVI restringido con una comparación alternativa, cuantiles con bandas, $W_1$ y $W_2$ con la
dirección de las colas, filtrado causal, innovación condicionada, pronósticos físicos sencillos,
calibración, CVaR con costos, pruebas incrementales y trazabilidad (\figref{fig:prioridades}).

En una segunda etapa entran la detección de régimen, la autorregresión en espacio de Wasserstein, el
conjunto de modelos y la optimización robusta. Como investigación avanzada quedan los modelos neuronales de
dinámica de opciones sin arbitraje, si el volumen y la diversidad de datos lo permiten; existen trabajos
específicos de neural SDE para superficies y riesgo \citep{cohen2021}, pero no constituyen evidencia de
rentabilidad de esta estrategia.

Se posponen el reinforcement learning, los transformers extensos, la generación de trayectorias por
difusión y el criterio de Kelly agresivo. Su complejidad y su sensibilidad a simuladores y pronósticos no se
justifican antes de demostrar la utilidad de las señales y la calidad de los datos. Ningún modelo se
incorpora solo por ser más sofisticado.

\begin{figure}[H]
\centering
\resizebox{0.94\linewidth}{!}{% Hoja de ruta: qué priorizar, qué dejar para después y qué posponer.
\begin{tikzpicture}[
  col/.style={cajallena=#1,text width=3.25cm,minimum height=9mm,font=\sffamily\scriptsize\bfseries},
  item/.style={caja=#1,text width=3.25cm,minimum height=7mm,font=\sffamily\tiny},
  node distance=1.6mm and 3mm]
\node[col=qazul] (c1) {INMEDIATO};
\node[col=qnaranja,right=of c1] (c2) {SEGUNDA ETAPA};
\node[col=qvioleta,right=of c2] (c3) {INVESTIGACIÓN AVANZADA};
\node[col=tenue,right=of c3] (c4) {POSPUESTO};
% Inmediato
\node[item=qazul,below=of c1] (i1) {Datos correctos y contrato de datos};
\node[item=qazul,below=of i1] (i2) {SSVI restringido y comparación alternativa};
\node[item=qazul,below=of i2] (i3) {Cuantiles y bandas de sensibilidad};
\node[item=qazul,below=of i3] (i4) {$W_1$ y $W_2$ con dirección de colas};
\node[item=qazul,below=of i4] (i5) {Filtrado causal};
\node[item=qazul,below=of i5] (i6) {Innovación condicionada};
\node[item=qazul,below=of i6] (i7) {Pronósticos físicos sencillos y calibración};
\node[item=qazul,below=of i7] (i8) {CVaR con costos};
\node[item=qazul,below=of i8] (i9) {Pruebas incrementales y trazabilidad};
% Segunda etapa
\node[item=qnaranja,below=of c2] (s1) {Detección de régimen};
\node[item=qnaranja,below=of s1] (s2) {Autorregresión Wasserstein};
\node[item=qnaranja,below=of s2] (s3) {Conjunto de modelos};
\node[item=qnaranja,below=of s3] (s4) {Optimización robusta};
% Investigación avanzada
\node[item=qvioleta,below=of c3,text width=3.25cm] (a1) {Modelos neuronales de dinámica de opciones sin
  arbitraje, si el volumen y la diversidad de datos lo permiten};
% Pospuesto
\node[item=tenue,below=of c4] (p1) {Reinforcement learning};
\node[item=tenue,below=of p1] (p2) {Transformers extensos};
\node[item=tenue,below=of p2] (p3) {Trayectorias por difusión};
\node[item=tenue,below=of p3] (p4) {Kelly agresivo};
\node[nota,text width=3.25cm,below=2mm of p4] {su complejidad y su sensibilidad a simuladores y pronósticos no
  se justifican antes de demostrar la utilidad de las señales};
\node[nota,text width=3.25cm,below=2mm of a1] {existen trabajos de neural SDE para superficies y riesgo, sin
  evidencia de rentabilidad para esta estrategia};
\end{tikzpicture}
}
\caption{Hoja de ruta. Cada columna exige que la anterior haya superado sus criterios.}
\label{fig:prioridades}
\end{figure}

\begin{ideaclavefinal}[En resumen]
La decisión arquitectónica central es conservar separadas la medición $\Q$, la predicción $\P$ y la acción
sobre la posición. La hipótesis más valiosa para investigar primero es si la innovación de las colas,
medida con incertidumbre y relativa al mercado, mejora la gestión de una posición existente. Todo el resto
del sistema existe para poder responder esa pregunta con honestidad, incluida la posibilidad de que la
respuesta sea no.
\end{ideaclavefinal}

% ============================================================================
\clearpage
\section{Estado del proyecto al 28 de septiembre de 2026}\label{sec:estado}
% ============================================================================

Las secciones anteriores describen el sistema tal como se propuso el 26 de septiembre. Esta sección cuenta
qué se construyó desde entonces, qué funciona ya, qué cambió respecto de la propuesta y qué falta. Refleja
el estado de la mañana del lunes 28 de septiembre de 2026, antes de la primera sesión capturada de forma
programada. El detalle de cada paso está en los documentos del repositorio: el resumen de la sesión con
Alpaca, la descripción de la fuente y las respuestas a cada revisión.

\begin{ideaclave}[En pocas palabras]
\begin{itemize}
\item \textbf{La idea.} Cada mañana, a las 09:45 de Nueva York, se mide cuánto más cuesta protegerse de una
caída del S\&P~500 que apostar a una subida (el RR25 de las opciones SPXW). Después se comprueba si su
cambio de un día a otro anticipa la dirección o el riesgo de la bolsa, medida con SPY.
\item \textbf{Lo logrado.} Hay fuente de datos, un programa que mide y calcula lo que la bolsa hizo después,
y un proceso automático que guarda los precios cada día hábil. Está instalado, probado y confirmado por una
verificación independiente.
\item \textbf{Lo que falta.} Vigilar las primeras 3--5 sesiones, juntar 20--40 para el piloto y, con bastante
más historia, hacer la prueba predictiva con reglas fijadas de antemano.
\item \textbf{Lo que todavía no se sabe.} Si la señal predice algo. No hay ningún resultado sobre eso.
\end{itemize}
\end{ideaclave}

\subsection{Qué se construyó}\label{sec:estado_construido}

La \tabref{tab:construido} resume las piezas. Las comprueban 203 pruebas automáticas que no usan la red. La
batería también pasa en Windows y con Python obligado a declarar la codificación de cada texto.

\begin{table}[H]
\centering\small
\caption{Piezas construidas y su estado.}
\label{tab:construido}
\begin{tabularx}{\linewidth}{@{}>{\RaggedRight}p{3.3cm} Y >{\RaggedRight\arraybackslash}p{3.0cm}@{}}
\toprule
\encab{Pieza} & \encab{Qué hace} & \encab{Estado}\\
\midrule
Núcleo de referencia & Los métodos de este documento, con filtros causales, cuantiles con mesetas y
volatilidad implícita robusta. & Corregido\\
Pipeline del piloto & Ingesta, controles, medidas (RR25, asimetría y paridad), etiquetas, informe y
manifiesto reproducible. & Listo, con plantilla sintética\\
Adaptador de Alpaca & Reintentos, crudo inmutable con sus hashes, normalización al contrato y nivel
implícito de SPX. & Verificado con datos reales\\
Captura diaria & Ráfaga 5~s antes de las 09:45 y de las 10:00, página a página, con recuperación y un
estado por hora. & Instalada el 27 de septiembre\\
Histórico del objetivo & SIP de SPY en cada corte, 15 minutos después, y dividendos de SPY. & Instalado el
27 de septiembre\\
Etiquetas y dividendos & Rendimiento total y de precio de SPY, con dividendos versionados y sus
discrepancias. & Listo\\
Revisión de la operación & Estado, puntualidad, respaldo, SIP y dividendos de cada corte. & En uso desde el
28 de septiembre\\
\bottomrule
\end{tabularx}
\end{table}

\subsection{Datos: Alpaca y la captura hacia adelante}\label{sec:estado_datos}

Alpaca se verificó el 26 de septiembre con su API y su documentación (\tabref{tab:alpaca}). El hallazgo que
más condiciona el plan es que \textbf{no guarda cotizaciones pasadas de opciones}: el bid/ask de las 09:45
de sesiones anteriores no se puede reconstruir. Por eso la muestra se captura hacia adelante, desde el 28 de
septiembre.

\begin{table}[H]
\centering\small
\caption{Qué ofrece Alpaca para este proyecto (verificado el 26 de septiembre de 2026).}
\label{tab:alpaca}
\begin{tabularx}{\linewidth}{@{}>{\RaggedRight}p{3.6cm} Y@{}}
\toprule
\encab{Tema} & \encab{Resultado}\\
\midrule
SPXW & Disponible: opciones europeas, liquidación PM y vencimientos todos los días hábiles.\\
Historia de cotizaciones & \textbf{No hay.} Solo barras y operaciones de opciones desde febrero de 2024.\\
Nivel de SPX & \textbf{No hay.} Se estima por paridad con el vencimiento SPXW más cercano.\\
Feed gratuito (\texttt{indicative}) & Según Alpaca, sus cotizaciones están \emph{modificadas}: no son el NBBO
de OPRA.\\
Feed OPRA & Requiere el plan Algo Trader Plus (99~USD al mes) y firmar el acuerdo de OPRA.\\
SPY & Cotización IEX en tiempo real; NBBO consolidado (SIP) consultable 15 minutos después; dividendos en los
eventos corporativos, sin garantía de cuándo se publican.\\
\bottomrule
\end{tabularx}
\end{table}

La primera verificación con datos reales usó las últimas cotizaciones del viernes 25 al cierre. No es una
sesión de apertura, pero mostró cómo es el feed gratuito:
\begin{itemize}
\item llegaron cotizaciones para los 3\,250 contratos pedidos;
\item solo el 9--15\,\% de los precios de SPXW cae en la grilla de ticks de la bolsa, cuando en un NBBO real
caerían todos;
\item entre el 20 y el 28\,\% de los pares de SPXW viola la banda de la paridad;
\item en el modo descriptivo del cierre, que no es elegible para el piloto, el RR25 a 30 días fue de
$-3.08$ puntos en SPXW y de $-3.20$ en SPY, con una sonrisa coherente.
\end{itemize}
Con este feed, los residuos de paridad no sirven como señal ni como control fino de calidad, y sin OPRA no se
puede medir cuánto sesga el RR25. El piloto mide el feed \texttt{indicative}; la decisión sobre OPRA se toma
con las primeras sesiones reales.

\subsection{La captura diaria}\label{sec:estado_captura}

La captura corre en GitHub Actions, dentro de un repositorio \textbf{privado}: los datos de Alpaca y de OPRA no
se pueden redistribuir, y el repositorio del código es público. Los workflows descargan el código del
repositorio público en un commit fijo, hoy \texttt{4ffde8a}, y cada registro guarda el commit usado. Cambiar
ese commit es cambiar la versión de medición.

\begin{figure}[H]
\centering
\resizebox{\linewidth}{!}{% Jornada de captura en el repositorio de datos (hora de Nueva York, EDT). t = minutos desde las 09:00.
\begin{tikzpicture}[x=1mm,y=1cm,
  disparo/.style={circle,fill=#1,minimum size=2.5mm,inner sep=0pt},
  respaldo/.style={circle,draw=#1,fill=white,line width=0.8pt,minimum size=2.5mm,inner sep=0pt},
  ventana/.style={fill=#1!6,draw=#1!30,line width=0.4pt},
  espera/.style={fill=#1!28,draw=none},
  carril/.style={titulo,anchor=east,align=right}]
% Ventanas en que la compuerta deja arrancar
\filldraw[ventana=qazul] (0,2.42) rectangle (43,2.78);
\filldraw[ventana=qazuloscuro] (0,1.52) rectangle (58,1.88);
\filldraw[ventana=qnaranja] (76,0.62) rectangle (135,0.98);
% Esperas hasta el corte
\fill[espera=qazul] (11,2.52) rectangle (45,2.68);
\fill[espera=qazuloscuro] (21,1.62) rectangle (60,1.78);
% Apertura y cortes
\draw[tenue,dashed,line width=0.6pt] (30,0) -- (30,3.1);
\node[nota,anchor=south] at (30,3.1) {apertura};
\draw[qazul,line width=0.9pt] (45,0) -- (45,3.1);
\node[nota,anchor=south,text=qazul] at (45,3.1) {\textbf{corte}};
\draw[qazuloscuro,line width=0.9pt] (60,0) -- (60,3.1);
\node[nota,anchor=south,text=qazuloscuro] at (60,3.1) {\textbf{corte}};
% Disparos y respaldos
\node[disparo=qazul] at (11,2.6) {};
\node[respaldo=qazul] at (26,2.6) {};
\node[disparo=qazuloscuro] at (21,1.7) {};
\node[respaldo=qazuloscuro] at (36,1.7) {};
\node[disparo=qnaranja] at (81,0.8) {};
% Carriles
\node[carril] at (-2,2.6) {captura-0945};
\node[carril] at (-2,1.7) {captura-1000};
\node[carril] at (-2,0.8) {historico-alpaca};
% Notas de cada carril
\node[nota,anchor=west,align=left] at (64,2.6) {ráfaga en paralelo 5 s antes del corte de las 09:45;\\cada página se guarda al llegar};
\node[nota,anchor=west,align=left] at (64,1.7) {la misma lógica para el corte de las 10:00,\\en otra máquina y otro proceso};
\node[nota,anchor=north west,align=left] at (76,0.56) {SIP de SPY en cada corte (consultable 15 min después) y dividendos};
% Eje
\draw[guia,line width=0.7pt] (0,0) -- (150,0);
\foreach \t/\h in {0/09:00,15/09:15,30/09:30,45/09:45,60/10:00,75/10:15,90/10:30,105/10:45,120/11:00,135/11:15,150/11:30}{
  \draw[guia] (\t,0) -- (\t,-0.1);
  \node[nota,anchor=north] at (\t,-0.12) {\h};}
% Leyenda
\begin{scope}[yshift=-0.85cm]
\node[disparo=tinta2] at (2,0) {};
\node[nota,anchor=west] at (4,0) {disparo programado};
\node[respaldo=tinta2] at (32,0) {};
\node[nota,anchor=west] at (34,0) {respaldo, 15 min después};
\fill[espera=tinta2] (68,-0.08) rectangle (76,0.08);
\node[nota,anchor=west] at (77,0) {espera hasta el corte};
\filldraw[ventana=tinta2] (104,-0.18) rectangle (112,0.18);
\node[nota,anchor=west] at (113,0) {la compuerta deja arrancar};
\end{scope}
\end{tikzpicture}
}
\caption{Jornada de captura en hora de Nueva York (EDT). Cada hora de corte tiene su disparo y su respaldo,
en máquinas separadas, y una compuerta deja arrancar solo a los que llegan a tiempo. Los horarios son los
programados: GitHub puede retrasarlos, y los registros miden el arranque real.}
\label{fig:jornada}
\end{figure}

Cada corte termina en un estado: \texttt{completa}, \texttt{parcial}, \texttt{fallida} o, si no hubo captura,
\texttt{perdida}. Cada ejecución deja un registro, también cuando falla al iniciar o después de capturar.
Tres mecanismos protegen la captura:
\begin{itemize}
\item un diario por página permite recuperar una captura interrumpida;
\item un plazo de preparación deja llegar a tiempo al respaldo;
\item un plazo absoluto impide que un proceso colgado siga corriendo.
\end{itemize}
Solo se guarda el crudo: las tablas se reconstruyen a partir de él, y reprocesarlo da los mismos bytes. Las
credenciales nunca se escriben en un archivo.

\subsection{Objetivo, etiquetas y dividendos}\label{sec:estado_objetivo}

El piloto separa tres papeles que no se mezclan:
\begin{itemize}
\item \textbf{señal}: las opciones SPXW;
\item \textbf{referencia de las opciones}: el nivel implícito de SPX, que usan los controles y las medidas de
la cadena;
\item \textbf{objetivo}: SPY observado en el SIP, con el rendimiento total (con dividendos) y el de precio por
separado.
\end{itemize}

Las etiquetas miden 1 y 5 sesiones, con la misma hora inicial y final. Alpaca no garantiza cuándo publica un
dividendo, así que cada dividendo se guarda en versiones, con la fecha desde la que se conoce cada una. Si
una consulta ya no trae un dividendo, o lo trae sin fecha ex o sin monto, se abre una discrepancia. Las
etiquetas que podría afectar quedan \emph{pendientes} hasta que la resuelva evidencia fechada. Una etiqueta
se acepta cuando una consulta hecha al menos 60 días después de su fin no deja discrepancias abiertas: es
una regla, no una garantía. El entrenamiento solo usa las etiquetas que ya habrían madurado en la fecha
simulada.

\subsection{Revisiones externas}\label{sec:estado_revisiones}

Cada entrega pasó por una revisión externa. Cada hallazgo se reprodujo, se convirtió en una prueba de
regresión y se corrigió (\tabref{tab:revisiones}). Entre tanto, la batería pasó de 142 a 203 pruebas.

\begin{table}[H]
\centering\small
\caption{Rondas de revisión y su resultado.}
\label{tab:revisiones}
\begin{tabularx}{\linewidth}{@{}>{\RaggedRight}p{3.1cm} Y >{\RaggedRight\arraybackslash}p{2.5cm}@{}}
\toprule
\encab{Revisión} & \encab{Qué encontró o pidió} & \encab{Corrección}\\
\midrule
\texttt{e2b92f0} & Ocho hallazgos, entre ellos la elegibilidad estricta, los estados de captura y la
separación de fuentes. & \texttt{79c1554}\\
\texttt{d815bdd} & Separar señal, referencia y objetivo; objetivo SPY observado; dividendos; recuperación de
capturas. & \texttt{03f2230} y siguientes\\
\texttt{8c97b2b} & Cobertura de dividendos aparte; versiones de eventos y etiquetas; plazo para el respaldo. &
\texttt{6f0e748}, \texttt{f2a5663}\\
\texttt{a5e2748} & Una ausencia no retira un dividendo, sino que abre una discrepancia; estado puntual de las
etiquetas. & \texttt{c5aaded}, \texttt{8316540}\\
\texttt{5c15028} & Un dividendo sin fecha ex o sin monto deja pendientes las etiquetas afectadas. &
\texttt{7fb0585}\\
Cierre de \texttt{b240dc2} & Sin defectos nuevos. Después: registro de cada ejecución a prueba de errores y
revisión de la operación. & \texttt{7a59d04}, \texttt{bbbdc6e}\\
Cambios operativos & Registro también al fallar al iniciar; ruta explícita de resoluciones; pruebas en
Windows. & \texttt{4ffde8a}\\
Verificación de \texttt{4ffde8a} & Sin defectos; recomienda fijarlo en los workflows. & \texttt{e1a7832}\\
Instalación & Verificación independiente del repositorio de datos y de la prueba manual. & Confirmada\\
\bottomrule
\end{tabularx}
\end{table}

\subsection{Instalación y aceptación}\label{sec:estado_instalacion}

El repositorio de datos se instaló la noche del 27 de septiembre, con las plantillas revisadas, idénticas
byte a byte a las del commit \texttt{e1a7832}. La prueba manual salió bien:
\begin{itemize}
\item la captura inmediata quedó \texttt{completa}, con 3\,250 cotizaciones, y su registro guarda el commit
\texttt{4ffde8a} sin cambios locales;
\item el histórico descargó completo el SIP de SPY en los 10 cortes del 21 al 25 de septiembre, además de los
dividendos;
\item no apareció ninguna credencial en los datos guardados.
\end{itemize}
Una verificación independiente descargó el repositorio y reprodujo los mismos resultados.

\begin{criterio}[Criterio para aceptar la operación]
Se revisan 3--5 sesiones reales, desde el 28 de septiembre, y en los dos cortes:
\begin{itemize}
\item el estado, el retraso del arranque, la preparación, el margen de la ráfaga y las respuestas tardías;
\item el comportamiento del respaldo cuando actúe;
\item el SIP y los dividendos, no solo el código de salida;
\item que las capturas del horario regular sean elegibles para el piloto y que cada registro conserve el
commit fijado.
\end{itemize}
\end{criterio}

\begin{notatecnica}[Mantenimiento antes del 19 de octubre]
GitHub pasará la etiqueta \texttt{ubuntu-latest} de Ubuntu 24.04 a 26.04 de forma gradual entre el 19 de
octubre y el 19 de noviembre \citep{github2026}. Eso cae en medio del piloto, y unas corridas usarían una
imagen y otras, la otra. Antes de esa fecha, las plantillas se fijarán en \texttt{ubuntu-24.04}, la imagen que
ya se usa, junto con versiones de las acciones hechas para Node.js 24. El cambio se instalará con una corrida
manual de prueba.
\end{notatecnica}

\subsection{Qué falta}\label{sec:estado_falta}

La \figref{fig:ruta} ordena lo que queda. Cada paso exige que el anterior haya superado su criterio.

\begin{figure}[H]
\centering
\resizebox{\linewidth}{!}{% Estado y ruta al 28 de septiembre de 2026: qué está hecho, qué sigue y qué queda para después.
\begin{tikzpicture}[
  col/.style={cajallena=#1,text width=2.75cm,minimum height=9mm,font=\sffamily\scriptsize\bfseries},
  item/.style={caja=#1,text width=2.75cm,minimum height=7mm,font=\sffamily\tiny},
  node distance=1.6mm and 2.4mm]
\node[col=bueno] (c1) {HECHO};
\node[col=qazul,right=of c1] (c2) {AHORA};
\node[col=qvioleta,right=of c2] (c3) {PRÓXIMO};
\node[col=qnaranja,right=of c3] (c4) {DESPUÉS};
\node[col=tenue,right=of c4] (c5) {MÁS ADELANTE};
% Hecho
\node[item=bueno,below=of c1] (h1) {Propuesta y núcleo de referencia};
\node[item=bueno,below=of h1] (h2) {Pipeline del piloto con datos sintéticos};
\node[item=bueno,below=of h2] (h3) {Fuente de datos verificada: Alpaca};
\node[item=bueno,below=of h3] (h4) {Objetivo SPY, etiquetas y dividendos};
\node[item=bueno,below=of h4] (h5) {Captura diaria instalada y verificada};
% Ahora
\node[item=qazul,below=of c2] (a1) {Aceptación: 3--5 sesiones, del 28 de septiembre al 2 de octubre};
\node[item=qazul,below=of a1] (a2) {Mantenimiento de las plantillas antes del 19 de octubre};
\node[item=qazul,below=of a2] (a3) {Decisión sobre OPRA};
% Próximo
\node[item=qvioleta,below=of c3] (p1) {Piloto de 20--40 sesiones (hasta el 23 de octubre o el 20 de noviembre)};
\node[item=qvioleta,below=of p1] (p2) {Informe y dictamen de los datos};
\node[item=qvioleta,below=of p2] (p3) {Umbrales congelados en una versión nueva};
% Después
\node[item=qnaranja,below=of c4] (d1) {Protocolo predictivo fijado antes de mirar};
\node[item=qnaranja,below=of d1] (d2) {Historia ampliada: más de un año o comprada};
\node[item=qnaranja,below=of d2] (d3) {Prueba predictiva con periodo reservado};
% Más adelante
\node[item=tenue,below=of c5] (m1) {Cuantiles y transporte (variante F)};
\node[item=tenue,below=of m1] (m2) {Confianza, conformal y régimen};
\node[item=tenue,below=of m2] (m3) {Decisión sobre la posición};
\node[nota,text width=2.75cm,below=2mm of m3] {cada bloque entra solo si el anterior superó su criterio};
\end{tikzpicture}
}
\caption{Estado y ruta al 28 de septiembre de 2026.}
\label{fig:ruta}
\end{figure}

\begin{enumerate}
\item \textbf{Aceptación de la captura}, del 28 de septiembre al 2 de octubre: revisar 3--5 sesiones con el
criterio anterior.
\item \textbf{Mantenimiento de las plantillas}, antes del 19 de octubre.
\item \textbf{Piloto de medición}, de 20--40 sesiones: la sesión 20 es el 23 de octubre y la 40, el 20 de
noviembre. Su salida es un informe reproducible con un dictamen de los datos: «apto para ampliar historia»,
«apto con limitaciones» o «insuficiente». Cualquier relación con los rendimientos que se vea ahí es
exploratoria.
\item \textbf{Protocolo predictivo}, escrito antes de mirar el periodo reservado
(\secref{sec:estado_cambios}).
\item \textbf{Historia ampliada.} Entre 20 y 40 sesiones no bastan para contrastar la hipótesis. Si el efecto
es pequeño, hará falta más de un año de datos, salvo que se compre historia (por ejemplo, a Cboe DataShop).
\item \textbf{Después}, y solo con evidencia: cuantiles y transporte, confianza y la decisión sobre la
posición.
\end{enumerate}

\begin{pendiente}
\begin{itemize}
\item Contratar OPRA (99~USD al mes) o seguir con el feed gratuito, según las primeras sesiones.
\item Alargar el piloto de 20 a 40 sesiones.
\item Comprar historia para la prueba predictiva o acumularla hacia adelante.
\end{itemize}
\end{pendiente}

\subsection{Qué cambió respecto de esta propuesta}\label{sec:estado_cambios}

\begin{itemize}
\item \textbf{Fuente de datos.} La \tabref{tab:fuentes} dejaba a Alpaca para explorar; hoy es la fuente del
piloto, con captura hacia adelante porque no guarda cotizaciones pasadas. Cboe DataShop sigue como
candidata para la historia ampliada.
\item \textbf{Objetivo.} Alpaca no publica el nivel de SPX, así que el objetivo es SPY observado en el SIP, con
rendimiento total. El SPX implícito queda solo como referencia de las opciones.
\item \textbf{Primera pregunta.} La hipótesis general de la \secref{sec:hipotesis} se concretó para la primera
prueba. En la ecuación~\eqref{eq:diferencial}, la señal es el cambio del RR25 respecto de la sesión anterior,
el objetivo es el signo del rendimiento de SPY a una sesión y la puntuación es la de Brier. El horizonte de
cinco sesiones y el riesgo de cola quedan como análisis secundarios.
\item \textbf{Comparación por pasos.} Con la misma familia de modelos, se agrega una fuente de información
cada vez: primero la base (rendimientos pasados, salto de apertura, volatilidad reciente y calendario);
luego la volatilidad implícita, el nivel y el cambio de la asimetría y la paridad; y, al final, los
cuantiles y el transporte.
\item \textbf{Paridad.} Con el feed gratuito, los residuos de paridad son un diagnóstico de calidad, no una
señal.
\end{itemize}

\begin{noconcluirfinal}
\begin{itemize}
\item La instalación prueba que la captura funciona, no que la señal sirva.
\item La primera verificación con datos reales es del cierre de un viernes, no de una apertura.
\item El piloto evaluará la calidad de los datos, no la capacidad predictiva. Todavía no hay base para
recomendar posiciones.
\end{itemize}
\end{noconcluirfinal}


\appendix
% ============================================================================
\clearpage
\section{Notación}\label{ap:notacion}
% ============================================================================

\begin{table}[H]
\centering\small
\caption{Símbolos principales. Algunos se reutilizan con significado local ($\alpha$, $\varepsilon$, $r$);
el contexto de cada sección los desambigua.}
\label{tab:notacion}
\begin{tabularx}{\linewidth}{@{}l Y@{}}
\toprule
\encab{Símbolo} & \encab{Significado}\\
\midrule
$S_t$, $F_{t,\tau}$, $D(t,T)$ & Precio del subyacente, forward al plazo $\tau$ y factor de descuento\\
$K$, $T$, $\tau$ & Strike, vencimiento y plazo constante de la distribución implícita (30 días)\\
$k=\ln(K/F)$, $w=\sigma^2T$ & Log-moneyness y varianza implícita total\\
$\theta_T$, $\rho$, $\varphi$, $\eta$, $\gamma$ & Parámetros SSVI: varianza ATM, asimetría, curvatura y función de potencia\\
$b_i$, $a_i$, $m_i$, $s_i$ & Bid, ask, mid y medio spread de la cotización $i$\\
$q^{\Q}(K)$, $g(k)$ & Densidad implícita del precio y función de Durrleman\\
$x_t(u)$, $u$ & Cuantil de $\ln(S_{t+\tau}/F_{t,\tau})$ bajo $\Q_t$ y nivel de probabilidad\\
$W_1$, $W_2$, $\Delta x_t(u)$ & Distancias de Wasserstein y cambio firmado de cuantiles\\
$\phi_j(u)$, $s_{tj}$ & Autofunciones y puntajes de la PCA funcional\\
$f_t$, $A$, $Q$, $R_t$ & Estado, transición, ruido de estado y ruido de observación del filtro\\
$\mathcal I^{base}_t$, $\epsilon_t(u)$ & Información base e innovación de la distribución implícita\\
$Y_{t,h}$, $d_{t,h}$ & Objetivo a horizonte $h$ y diferencial de pérdida entre modelos\\
$\rho_\tau$, CRPS, Brier & Pérdida cuantílica y reglas de puntuación\\
$r_t$, $H$ & Longitud de corrida y probabilidad de riesgo de BOCPD\\
$\alpha_t$, $\gamma$ & Nivel efectivo y paso del conformal adaptativo\\
$w$, $\kappa$, $\lambda$, $\mathcal C_t$ & Pesos, costos proporcionales, aversión al riesgo y conjunto factible\\
$\CVaR_\alpha$, $\mathcal R_\delta$ & Valor en riesgo condicional al nivel $\alpha$ y región casi óptima\\
$\mathcal B_\varepsilon(\widehat P_t)$ & Bola de Wasserstein de radio $\varepsilon$ alrededor de la distribución estimada\\
\bottomrule
\end{tabularx}
\end{table}

% ============================================================================
\section{Glosario}\label{ap:glosario}
% ============================================================================

\begin{description}[leftmargin=0pt,itemsep=3pt,font=\sffamily\bfseries\small]
\item[Arbitraje estático.] Combinación de opciones de un mismo instante que asegura ganancia sin riesgo. Sus
dos formas son de \emph{calendario} (entre vencimientos) y de \emph{mariposa} (entre strikes).
\item[Breeden--Litzenberger.] Identidad que obtiene la densidad implícita como segunda derivada del precio
de las calls respecto del strike.
\item[BOCPD.] Detección bayesiana en línea de puntos de cambio mediante la posterior de la longitud de corrida.
\item[Conformal adaptativo.] Método de intervalos que ajusta su nivel efectivo según los errores recientes y
garantiza la tasa de error promedio a largo plazo.
\item[CRPS.] Regla de puntuación de pronósticos de distribución; integra la pérdida cuantílica en todos los niveles.
\item[CVaR.] Promedio de las pérdidas en la cola especificada (por ejemplo, el 5\% peor) bajo los escenarios supuestos.
\item[Desamericanización.] Conversión de precios de opciones americanas en equivalentes europeos mediante un
modelo con ejercicio anticipado y dividendos.
\item[Embargo y purga.] Eliminación de observaciones cuya etiqueta cruza una frontera temporal entre tramos, más
un margen de seguridad.
\item[Idempotencia.] Propiedad de una orden cuyo reintento con la misma clave no crea duplicados.
\item[Innovación.] Parte del cambio de la distribución implícita no explicada por la información base.
\item[PIT.] Valor $F(y)$ del pronóstico en el resultado observado; uniforme si el pronóstico está calibrado.
\item[$\Q$ y $\P$.] Distribución implícita en los precios de las opciones y distribución física de los
rendimientos.
\item[Región casi óptima.] Conjunto de carteras factibles cuyo valor estimado está dentro de la incertidumbre
del óptimo; base de la zona de no operación.
\item[Sharpe deflactado.] Probabilidad de que el Sharpe verdadero supere el máximo esperado por azar dado el
número de variantes probadas.
\item[SSVI.] Parametrización de la superficie de varianza implícita con condiciones suficientes sencillas de
ausencia de arbitraje estático.
\item[Transporte óptimo.] Forma de mover la masa de una distribución a otra con costo mínimo; en una dimensión
es el reordenamiento monótono de cuantiles.
\end{description}

% ============================================================================
\section{Decisiones pendientes}\label{ap:pendientes}
% ============================================================================

\begin{table}[H]
\centering\small
\caption{Decisiones que parametrizan el sistema y dónde se usan.}
\label{tab:pendientes}
\begin{tabularx}{\linewidth}{@{}c l Y@{}}
\toprule
& \encab{Decisión} & \encab{Dónde se usa}\\
\midrule
\iconopendiente & Universo definitivo & Datos, escenarios conjuntos, límites de concentración\\
\iconopendiente & Moneda de los 50,000 & Costos, conversión, benchmark comparable\\
\iconopendiente & Posiciones actuales & Comparación mantener frente a modificar; zona de no operación\\
\iconopendiente & Horizonte de decisión principal & Objetivos, pérdidas, etiquetas y purga\\
\iconopendiente & Presupuesto y proveedor de datos & Historia disponible, latencia y horario de decisión\\
\iconopendiente & Modelo de costos & Programa lineal y zona de no operación\\
\iconopendiente & Presupuesto de riesgo y límites & Conjunto factible $\mathcal C_t$ y reducciones obligatorias\\
\iconopendiente & ¿Se gestionarán opciones? & Rama de revaluación completa\\
\iconopendiente & Umbrales de la regla de avance & Deben fijarse antes de abrir la prueba final\\
\bottomrule
\end{tabularx}
\end{table}

% ============================================================================
\section{Reproducibilidad de las figuras}\label{ap:reproducibilidad}
% ============================================================================

Todas las gráficas se generan con datos sintéticos, semillas fijas y el paquete de referencia
\texttt{quantileflow/} del repositorio. Se dibujan con matplotlib en su estilo por defecto, una figura por
gráfica y con textos en inglés; los pies, en español, explican cada una. Desde la raíz del repositorio:

\begin{lstlisting}[style=registro]
pip install -r requirements.txt
python -m pytest                              # pruebas de propiedades del núcleo
python docs/propuesta/figuras_src/generar.py  # regenera las 62 figuras PNG
make pdf                                      # compila este documento
\end{lstlisting}

\begin{table}[H]
\centering\small
\caption{Origen de las figuras.}
\label{tab:origen}
\begin{tabularx}{\linewidth}{@{}l l Y@{}}
\toprule
\encab{Figuras} & \encab{Script} & \encab{Mundo sintético}\\
\midrule
\texttt{s1\_*}, \texttt{s2\_*} & \texttt{fig\_s1\_s2.py} & Calendario de 30 días; cadena SSVI con spreads y bid nulo\\
\texttt{s3\_*} & \texttt{fig\_s3.py} & Panel de 260 sesiones: mercado y acción A, choque común y evento idiosincrático\\
\texttt{s4\_*}, \texttt{s5\_*} & \texttt{fig\_s4\_s5.py} & Supuesto $\Q$/$\P$; regresión cuantílica; cambio de varianza; desplazamiento de volatilidad\\
\texttt{s6\_*}, \texttt{s7\_*}, \texttt{s9\_*} & \texttt{fig\_s6\_s7.py} & Cópulas; cartera de cuatro activos; remuestreo de óptimos; nulas de Sharpe\\
\bottomrule
\end{tabularx}
\end{table}

Las pruebas verifican, entre otras propiedades, que la superficie de referencia cumple las condiciones
suficientes de SSVI, que el contraejemplo SVI tiene $g<0$, que la densidad por Breeden--Litzenberger
coincide con la analítica e integra a uno con media igual al forward, que $W_2$ de una traslación es la
traslación, que el CVaR del programa lineal coincide con el empírico, que costos altos producen no
operación, que el conformal adaptativo cumple su cota de largo plazo y que el máximo Sharpe esperado
coincide con la simulación.


\clearpage
% ============================================================================
% Referencias (autor-año). Las fuentes de la propuesta original conservan su enlace.
% ============================================================================
\renewcommand{\refname}{Referencias}
\begin{thebibliography}{99}
\small
\setlength{\itemsep}{2pt}

\bibitem[Adams y MacKay(2007)]{adams2007}
Adams, R.~P. y MacKay, D.~J.~C. (2007). \emph{Bayesian online changepoint detection}. arXiv:0710.3742.
\url{https://arxiv.org/abs/0710.3742}

\bibitem[A{\"i}t-Sahalia y Lo(1998)]{aitsahalia1998}
A{\"i}t-Sahalia, Y. y Lo, A.~W. (1998). Nonparametric estimation of state-price densities implicit in
financial asset prices. \emph{Journal of Finance}, 53(2), 499--547.

\bibitem[Alpaca(s.f.)]{alpaca}
Alpaca. \emph{Historical option data} (documentación). \url{https://docs.alpaca.markets/us/docs/historical-option-data}

\bibitem[Bailey y L{\'o}pez~de~Prado(2014)]{bailey2014}
Bailey, D.~H. y López de Prado, M. (2014). The deflated Sharpe ratio: correcting for selection bias,
backtest overfitting and non-normality. \emph{Journal of Portfolio Management}, 40(5), 94--107.
\url{https://www.davidhbailey.com/dhbpapers/deflated-sharpe.pdf}

\bibitem[Borovi{\v{c}}ka et~al.(2016)]{borovicka2016}
Borovička, J., Hansen, L.~P. y Scheinkman, J.~A. (2016). Misspecified recovery. \emph{Journal of
Finance}, 71(6), 2493--2544.

\bibitem[Boyd et~al.(2017)]{boyd2017}
Boyd, S., Busseti, E., Diamond, S., Kahn, R.~N., Koh, K., Nystrup, P. y Speth, J. (2017). Multi-period
trading via convex optimization. \emph{Foundations and Trends in Optimization}, 3(1), 1--76.
\url{https://stanford.edu/~boyd/papers/cvx_portfolio.html}

\bibitem[Breeden y Litzenberger(1978)]{breeden1978}
Breeden, D.~T. y Litzenberger, R.~H. (1978). Prices of state-contingent claims implicit in option prices.
\emph{Journal of Business}, 51(4), 621--651.

\bibitem[Burkovska et~al.(2018)]{burkovska2018}
Burkovska, O. et~al. (2018). Calibration to American options: numerical investigation of the
de-Americanization method. \emph{Quantitative Finance}. \url{https://pmc.ncbi.nlm.nih.gov/articles/PMC6034575/}

\bibitem[Carr y Wu(2009)]{carrwu2009}
Carr, P. y Wu, L. (2009). Variance risk premiums. \emph{Review of Financial Studies}, 22(3), 1311--1341.

\bibitem[Cboe DataShop(s.f.)]{cboe}
Cboe DataShop. \emph{Option quote intervals} (ficha de producto). \url{https://datashop.cboe.com/option-quote-intervals}

\bibitem[Chernozhukov et~al.(2010)]{chernozhukov2010}
Chernozhukov, V., Fernández-Val, I. y Galichon, A. (2010). Quantile and probability curves without
crossing. \emph{Econometrica}, 78(3), 1093--1125.

\bibitem[Cohen et~al.(2021)]{cohen2021}
Cohen, S.~N., Reisinger, C. y Wang, S. (2021). \emph{Arbitrage-free neural-SDE market models}.
arXiv:2105.11053. \url{https://arxiv.org/abs/2105.11053}

\bibitem[Cox et~al.(1979)]{cox1979}
Cox, J.~C., Ross, S.~A. y Rubinstein, M. (1979). Option pricing: a simplified approach. \emph{Journal of
Financial Economics}, 7(3), 229--263.

\bibitem[Demarta y McNeil(2005)]{demarta2005}
Demarta, S. y McNeil, A.~J. (2005). The t copula and related copulas. \emph{International Statistical
Review}, 73(1), 111--129.

\bibitem[Duan et~al.(2020)]{duan2020}
Duan, T., Avati, A., Ding, D.~Y., Thai, K.~K., Basu, S., Ng, A. y Schuler, A. (2020). NGBoost: natural
gradient boosting for probabilistic prediction. \emph{Proceedings of the 37th International Conference on
Machine Learning}. arXiv:1910.03225. \url{https://arxiv.org/abs/1910.03225}

\bibitem[Gatheral y Jacquier(2014)]{gatheral2014}
Gatheral, J. y Jacquier, A. (2014). Arbitrage-free SVI volatility surfaces. \emph{Quantitative Finance},
14(1), 59--71. arXiv:1204.0646. \url{https://arxiv.org/abs/1204.0646}

\bibitem[Gibbs y Cand{\`e}s(2021)]{gibbs2021}
Gibbs, I. y Candès, E. (2021). Adaptive conformal inference under distribution shift. \emph{Advances in
Neural Information Processing Systems}, 34. arXiv:2106.00170. \url{https://arxiv.org/abs/2106.00170}

\bibitem[GitHub(2026)]{github2026}
GitHub (2026). \emph{Ubuntu 26 generally available and latest migration}. GitHub Changelog, 17 de
septiembre de 2026.
\url{https://github.blog/changelog/2026-09-17-ubuntu-26-generally-available-and-latest-migration/}

\bibitem[Gneiting y Raftery(2007)]{gneiting2007}
Gneiting, T. y Raftery, A.~E. (2007). Strictly proper scoring rules, prediction, and estimation.
\emph{Journal of the American Statistical Association}, 102(477), 359--378.
\url{https://sites.stat.washington.edu/people/raftery/Research/PDF/Gneiting2007jasa.pdf}

\bibitem[Gneiting y Ranjan(2011)]{gneiting2011}
Gneiting, T. y Ranjan, R. (2011). Comparing density forecasts using threshold- and quantile-weighted
scoring rules. \emph{Journal of Business \& Economic Statistics}, 29(3), 411--422.

\bibitem[Kalman(1960)]{kalman1960}
Kalman, R.~E. (1960). A new approach to linear filtering and prediction problems. \emph{Journal of Basic
Engineering}, 82(1), 35--45.

\bibitem[Koenker y Bassett(1978)]{koenker1978}
Koenker, R. y Bassett, G. (1978). Regression quantiles. \emph{Econometrica}, 46(1), 33--50.

\bibitem[K{\"u}nsch(1989)]{kunsch1989}
Künsch, H.~R. (1989). The jackknife and the bootstrap for general stationary observations. \emph{Annals of
Statistics}, 17(3), 1217--1241.

\bibitem[Ledoit y Wolf(2004)]{ledoit2004}
Ledoit, O. y Wolf, M. (2004). A well-conditioned estimator for large-dimensional covariance matrices.
\emph{Journal of Multivariate Analysis}, 88(2), 365--411.

\bibitem[Lee(2004)]{lee2004}
Lee, R.~W. (2004). The moment formula for implied volatility at extreme strikes. \emph{Mathematical
Finance}, 14(3), 469--480.

\bibitem[L{\'o}pez~de~Prado(2018)]{lopezdeprado2018}
López de Prado, M. (2018). \emph{Advances in Financial Machine Learning}. Wiley.

\bibitem[Martin(2017)]{martin2017}
Martin, I. (2017). What is the expected return on the market? \emph{Quarterly Journal of Economics},
132(1), 367--433.

\bibitem[Martin y Wagner(2019)]{martin2019}
Martin, I. y Wagner, C. (2019). What is the expected return on a stock? \emph{Journal of Finance}, 74(4),
1887--1929.

\bibitem[Mohajerin~Esfahani y Kuhn(2018)]{esfahani2018}
Mohajerin Esfahani, P. y Kuhn, D. (2018). Data-driven distributionally robust optimization using the
Wasserstein metric: performance guarantees and tractable reformulations. \emph{Mathematical Programming},
171, 115--166. arXiv:1505.05116. \url{https://arxiv.org/abs/1505.05116}

\bibitem[Patton(2011)]{patton2011}
Patton, A.~J. (2011). Volatility forecast comparison using imperfect volatility proxies. \emph{Journal of
Econometrics}, 160(1), 246--256.

\bibitem[Petersen y M{\"u}ller(2016)]{petersen2016}
Petersen, A. y Müller, H.-G. (2016). Functional data analysis for density functions by transformation to a
Hilbert space. \emph{Annals of Statistics}, 44(1), 183--218.

\bibitem[Peyr{\'e} y Cuturi(2019)]{peyre2019}
Peyré, G. y Cuturi, M. (2019). Computational optimal transport. \emph{Foundations and Trends in Machine
Learning}, 11(5--6), 355--607. arXiv:1803.00567. \url{https://arxiv.org/abs/1803.00567}

\bibitem[Rockafellar y Uryasev(2000)]{rockafellar2000}
Rockafellar, R.~T. y Uryasev, S. (2000). Optimization of conditional value-at-risk. \emph{Journal of
Risk}, 2(3), 21--41.

\bibitem[Ross(2015)]{ross2015}
Ross, S. (2015). The recovery theorem. \emph{Journal of Finance}, 70(2), 615--648.

\bibitem[Timmermann(2006)]{timmermann2006}
Timmermann, A. (2006). Forecast combinations. En \emph{Handbook of Economic Forecasting}, vol.~1,
135--196. Elsevier.

\bibitem[Zhang et~al.(2022)]{zhang2022}
Zhang, C., Kokoszka, P. y Petersen, A. (2022). Wasserstein autoregressive models for density time series.
\emph{Journal of Time Series Analysis}, 43(1), 30--52. arXiv:2006.12640. \url{https://arxiv.org/abs/2006.12640}

\end{thebibliography}


\end{document}
